% !TeX spellcheck = en_GB
%\documentclass[a4paper,12pt,titlepage,draft]{book} %{report}
\documentclass[a4paper,12pt,titlepage]{book} %{report}







\usepackage[skip=4pt]{parskip}
\usepackage{amsfonts}
\usepackage{amsmath}
\usepackage{amssymb}
\usepackage{amsthm}
\usepackage{titlesec}
\usepackage{listings}
\usepackage{fancyhdr}
\usepackage[export]{adjustbox}
\usepackage{setspace}
\usepackage{hyperref}
\usepackage{imakeidx}
\usepackage[T1]{fontenc}
\usepackage[utf8]{inputenc}
\usepackage{url}
\usepackage[ngerman,english]{babel}
\usepackage{comment} 

% slashbox veraltet
%\usepackage{slashbox}
\usepackage{diagbox}
%
%Und dann
%
%\diagbox{A}{B}


\usepackage{xspace} 
\usepackage{enumitem}
\usepackage{mathrsfs}
\usepackage{subcaption}
\usepackage{color}
\usepackage{multirow}
\usepackage{colortbl}
\usepackage{booktabs}
\usepackage{citeall}
\usepackage{longtable}
\usepackage{changepage}
\usepackage{nicematrix}
\usepackage{MnSymbol}
\usepackage{makecell}
\usepackage{tcolorbox}
\usepackage{diagbox} % https://tex.stackexchange.com/questions/7262/diagonally-divided-table-cell?noredirect=1&lq=1
\usepackage{accents} 
\usepackage{dirtree}
\usepackage{tikz}
\usepackage{xparse}
\usepackage{stackengine}
\usepackage{mathtools}
\usepackage{tabularx}
\usepackage{environ}
\usepackage{cleveref}

\usetikzlibrary{shapes,shapes.geometric,arrows,arrows.meta,fit,calc,bending}
\usetikzlibrary{automata,positioning}
\usetikzlibrary{shadows.blur}
\usetikzlibrary{intersections}
\usetikzlibrary{cd}
\usetikzlibrary{matrix,patterns,shadings,backgrounds}
\usetikzlibrary{babel} % https://tex.stackexchange.com/questions/331480/missing-endcsname-in-tikzcd

\tikzstyle{transition} = [font=\small]
\tikzstyle{line} = [draw,thick,-latex]

\setlength{\extrarowheight}{3pt}



\usepackage[color=gray!27]{todonotes}
%\usepackage[disable]{todonotes}

\usepackage[url=false,sorting=anyt,backref=true]{biblatex}
\addbibresource{../tidied5.bib}


\DefineBibliographyStrings{english}{%
	backrefpage = {cited on page},% originally "cited on page"
	backrefpages = {cited on pages},% originally "cited on pages"
}


\setenumerate[1]{label={(\arabic*)}}

\newcommand{\mathaord}{} % sonst gibt es eine Fehlermeldung mit "accents"


\makeindex

%
%\usepackage{ifdraft}
%% https://tex.stackexchange.com/questions/59530/draft-option-in-tikz-like-graphicx
%\ifoptionfinal{}{%
%	\usetikzlibrary{external}\tikzexternalize
%}

\DeclareMathOperator{\lcm}{lcm}

\newcommand{\pref}[1]{\operatorname{Pref}(#1)}
\newcommand{\suff}[1]{\operatorname{Suff}\left(#1\right)}
\newcommand{\factor}[1]{\operatorname{Fact}(#1)}
\newcommand{\fact}[1]{\ensuremath{\operatorname{Fact}(#1)}}
\newcommand{\constrsync}[1]{\text{\normalfont \textsc{Sync}\textsc{[}\mbox{$#1$}\textsc{]}}\xspace}
\newcommand{\concat}[2]{#2 #1}
\newcommand{\concatcirc}[2]{\ensuremath{#2 \circ #1}}
\newcommand{\applyfunc}[2]{#2(#1)}
\newcommand{\applytrans}[3]{\ensuremath{#1(#2,#3)}}
\newcommand{\applytransexplicit}[3]{\ensuremath{#1(#2,#3)\xspace\,}}
\newcommand{\semiautomat}[4]{\mathcal #1 = (#2, #3, #4)}
\newcommand{\automat}[6]{\ensuremath{\mathcal #1 = (#2, #3, #4, #5, #6)}}  % entfernt überstrich, wenn man es so benutzt: \automat{\overline{\mathcal A}_i}{Q_i \cup \{p_i\}}{\overline{\delta_i}}{p_i}{F_i}
\newcommand{\semiautomatexpl}[4]{\ensuremath{\mathcal #1 = (#2, #3, #4)}}
\newcommand{\automatp}[6]{\ensuremath{#1 = (#2, #3, #4, #5, #6)}}
\newcommand{\perm}{\operatorname{perm}}
\newcommand{\Syn}{\operatorname{Syn}}
\newcommand{\Cerny}{\v{C}ern\'{y}\xspace}
\newcommand{\orbit}[2]{\ensuremath{#2(#1)}}
\newcommand{\Orb}{\operatorname{Orb}}
\newcommand{\rt}[1]{\operatorname{rt}{(#1)}}

\newcommand{\rank}[1]{\ensuremath{\operatorname{rank}(#1)}}
\newcommand{\defect}[1]{\ensuremath{\operatorname{def}(#1)}} 
\newcommand{\Orbit}[2]{\ensuremath{\operatorname{Orb}_{#2}(#1)}}
\newcommand{\Stab}[2]{\ensuremath{{#2}_{#1}}}
\newcommand{\lereach}{\ensuremath{\precsim_R}}
\newcommand{\eqreach}{\ensuremath{\simeq_R}}
\newcommand{\lelogm}{\ensuremath{\le_m^{\log}}}
\newcommand{\lePm}{\ensuremath{\le_m^{\PTIME}}}
\newcommand{\leACZEROm}{\ensuremath{\le_m^{\AC^0}}}
\newcommand{\leACkm}[1]{\ensuremath{\le_m^{\AC^{#1}}}}
\newcommand{\leNCkm}[1]{\ensuremath{\le_m^{\NC^{#1}}}}
\newcommand{\leNCkT}[1]{\ensuremath{\le_T^{\NC^{#1}}}}
\newcommand{\eqlogm}{\ensuremath{\equiv_m^{\log}}}
\newcommand{\eqPm}{\ensuremath{\equiv_m^{\PTIME}}}
\newcommand{\eqACZEROm}{\ensuremath{\equiv_m^{\AC^0}}}
\newcommand{\eqACkm}[1]{\ensuremath{\equiv_m^{\AC^{#1}}}}
\newcommand{\eqNCkm}[1]{\ensuremath{\equiv_m^{\NC^{#1}}}}
\newcommand{\eqNCkT}[1]{\ensuremath{\equiv_T^{\NC^{#1}}}}
\newcommand{\fword}[2]{\ensuremath{{#1}_{#2}}}
\newcommand{\SI}[2]{\ensuremath{[#1 \mapsto #2]}}
\newcommand{\move}[1]{\operatorname{move}(#1)}
\newcommand{\fix}[1]{\operatorname{fix}(#1)}
\newcommand{\dupl}[1]{\operatorname{dupl}(#1)}
\newcommand{\excl}[1]{\operatorname{excl}(#1)}
\newcommand{\TM}[1]{\mathcal T_M(#1)}
\newcommand{\TS}[1]{\mathcal T_S(#1)}
\newcommand{\Z}{\mathbb{Z}}
\newcommand{\PA}[1]{\mathcal P(#1)}

\newcommand{\rYES}{{\sc YES}\xspace}
\newcommand{\rNO}{{\sc NO}\xspace}
\newcommand{\eqLp}[1]{\mathrel{\ensuremath{\equiv_{L}^{\ge #1}}}}
\newcommand{\eqLn}{\mathrel{\eqLp{n}}}
\newcommand{\eqLpp}[2]{\mathrel{\ensuremath{\equiv_{#1}^{\ge #2}}}}
\newcommand{\eqLL}[2]{\mathrel{\ensuremath{\equiv_{#1, #2}}}}
\newcommand{\NeqLL}[2]{\mathrel{\ensuremath{\not\equiv_{#1,#2}}}}
\newcommand{\eqLsimU}[3]{\mathrel{\ensuremath{\equiv_{#1}^{#2, #3}}}}
\newcommand{\preeqLL}[2]{\mathrel{\ensuremath{\sim_{#1, #2}}}}
 \newcommand{\excess}[2]{\operatorname{exc}_{#1}(#2)}
 
%\newcommand{\shift}[1]{\mathsf{X}_{#1}} % X_{#1}
%\newcommand{\term}[1]{\textit{#1}}
%\newcommand{\leftm}{{\vartriangleright}}  % "left marker"
%\newcommand{\rightm}{{\vartriangleleft}} % "right marker"

\newcommand{\Zb}{\mathbb{Z}}
\newcommand{\minaut}[1]{\ensuremath{\operatorname{MinAut}(#1)}}
\newcommand{\restrictAutGamS}[3]{\ensuremath{#1_{|#2, #3}}}
\newcommand{\upperset}[1]{\mbox{\ensuremath{\displaystyle  \uparrow \!\!#1}}}
\newcommand{\lowerset}[1]{\ensuremath{\downarrow \!\!#1}}
\renewcommand{\vec}[1]{\ensuremath{\boldsymbol{#1}}}


\newcommand{\PTIME}{\textsf{P}\xspace}
\newcommand{\LOGSPACE}{\textsf{L}\xspace}
\newcommand{\WNL}{\textsf{WNL}\xspace}
\newcommand{\WSYNC}{\textsf{W[Sync]}\xspace}

%\renewcommand{\SAT}{\textsc{\sc SAT}\xspace}
\newcommand{\SIS}{\textsc{\sc Sync\-Into}\xspace}
\newcommand{\ST}{\textsc{\sc Set\-Trans}\xspace}
\newcommand{\DST}{\textsc{\sc Disj\-Set\-Trans}\xspace}
\newcommand{\SFS}{\text{\sc Sync\-From}\xspace}
\newcommand{\DFAS}{\text{\sc Sync}\xspace}
\newcommand{\DFASWS}{\text{\sc Sync\-Short}\xspace}
\newcommand{\DFASS}{\text{\sc Sync\-State}\xspace}
\newcommand{\DFASW}{\text{\sc Sync\-Words}\xspace}
\newcommand{\DFAINT}{\text{\sc Intersection}\xspace}
\newcommand{\DAG}{\text{\sc DAG}\xspace} 
\newcommand{\DFAST}{\textsc{\sc Sync\-Trap}\xspace}
\newcommand{\THREESAT}{\text{\sc 3SAT}\xspace}
\newcommand{\ONEDAGST}{\text{\sc $1$-DAG-ST}\xspace}
\newcommand{\ONEDAG}{\text{\sc $1$-DAG}\xspace}
\newcommand{\TWODAG}{\text{\sc $2$-DAG}\xspace}
\newcommand{\TWODAGST}{\text{\sc $2$-DAG-ST}\xspace}
\newcommand{\KDAG}{\text{\sc $k$-DAG}\xspace}
\newcommand{\SCP}{\text{\sc Single\-Cycle}\xspace}
\newcommand{\PDFAINT}{\text{\sc PDFA-Intersection}\xspace}
\newcommand{\SCC}{\text{\sc Strongly\-Conn}\xspace}
\newcommand{\kDAGST}[1]{\text{\sc $#1$-DAG-ST}\xspace}
\newcommand{\KDAGST}{\kDAGST{k}\xspace}
\newcommand{\ORD}{{\sc OrderBetweenVertices}\xspace}
\newcommand{\REACH}{\text{\sc Reach}\xspace}
\newcommand{\REACHF}{\text{\sc ReachFrom}\xspace}

\newcommand{\problemdef}[4]{
	% \problemdef{title}{non-small-caps reference/comment/... in title}{input}{problem}
	% Next line: width of the drawn box, spacing left (right with "right = ...")
	\begin{tcolorbox}[width = \textwidth,colback=white,arc=0pt,outer arc=0pt,boxrule=0.7pt,left =0.5em,right=0em]
		\textsc{#1} #2
		\\[2pt]
		% Next line: table width
		\begin{tabular}{ @{}l p{0.84\textwidth} c }
			\textsf{Input:} & #3 \\[.5pt]
			\textsf{Problem:} & #4
		\end{tabular}
		\vspace{-0.25em}
	\end{tcolorbox}
}

\crefname{lemma}{Lemma}{Lemmas}
\crefname{proposition}{Proposition}{Propositions}
\crefname{theorem}{Theorem}{Theorems}
\crefname{equation}{Equation}{Equations}
\crefname{item}{Item}{Items}
\crefname{corollary}{Corollary}{Corollaries}
\crefname{section}{Section}{Sections}
\crefname{example}{Example}{Examples}
\crefname{Fact}{Fact}{Facts}
\crefname{table}{Table}{Tables}
\crefname{remark}{Remark}{Remarks}
\crefname{part}{Part}{Parts}
\crefname{figure}{Figure}{Figures}
\crefname{chapter}{Chapter}{Chapters}
\crefname{definition}{Definition}{Definitions}
\crefname{corollary}{Corollary}{Corollaries}

% https://tex.stackexchange.com/questions/410585/cleveref-format-for-enumerate
\crefname{enumi}{Item}{Items}
\crefname{enumii}{Item}{Items}
\crefname{enumiii}{Item}{Items}

% https://latex.org/forum/viewtopic.php?t=6667
\newcommand{\crefitem}[2]{\cref{#1},~Item~\ref{#2}}
\newcommand{\Crefitem}[2]{\Cref{#1},~Item~\ref{#2}}


%%\newcommand{\THM}{Thm.}
%%\newcommand{\COR}{Cor.}
%%\newcommand{\SEC}{Sect.}
%%\newcommand{\PROP}{Prop.}
%%\newcommand{\REM}{Rem.}
%%\newcommand{\LEM}{Lem.}
%%\newcommand{\EQN}{Eqn.}
%%\newcommand{\FIG}{Fig.}
\newcommand{\THM}{Theorem}
\newcommand{\COR}{Corollary}
\newcommand{\PROP}{Proposition}
%\newcommand{\REM}{Remark}
\newcommand{\SEC}{Section}
\newcommand{\LEM}{Lemma}
%\newcommand{\EQN}{Equation}
\newcommand{\FIG}{Figure}
%\newcommand{\EX}{Example}
%\newcommand{\CHAP}{Chapter}





%% https://tex.stackexchange.com/questions/570985/how-to-use-a-fact-as-an-equation-in-latex
%\newcommand\myparbox[1]{\parbox{0.667\textwidth}{\raggedright #1}} % choose width suitably
%
%
%
%
%% \sqcup  oder \square
%% https://tex.stackexchange.com/questions/136485/empty-clause-symbol
%\newcommand{\TMblank}{%
%\text{\fboxsep=-.2pt\fbox{\rule{0pt}{1.35ex}\rule{1.15ex}{0pt}}}%
%}

\tikzstyle{every state}=[blur shadow,fill={rgb:black,0.4;white,10}]


% Nur fuer pdflatex
%\DeclareUnicodeCharacter{0301}{*************************************}
\DeclareUnicodeCharacter{0301}{\i}








\DeclareSymbolFont{rsfscript}{OMS}{rsfs}{m}{n}
\DeclareSymbolFontAlphabet{\mathrsfs}{rsfscript}

% https://tex.stackexchange.com/questions/22012/hyperref-and-imakeidx-do-not-work-together
\newcommand{\BHT}[1]{\normalsize{\hyperpage{#1}}\normalsize}
\newcommand{\IN}[1]{\index{#1|BHT}}
%\renewcommand{\index[1]{\index{#1|BH}}

% Mit \IN geht das nicht
%\IN{function composition|see{composition}}
%\index{function composition|see{composition}}






% https://tex.stackexchange.com/questions/21234/doing-something-only-when-the-draft-option-is-on
\makeatletter
\def\ifdraft{\ifdim\overfullrule>\z@
	\expandafter\@firstoftwo\else\expandafter\@secondoftwo\fi}
\makeatother

% https://tex.stackexchange.com/questions/32208/footnote-runs-onto-second-page
\interfootnotelinepenalty=10000 %% Completely prevent breaking of footnotes





% https://tex.stackexchange.com/questions/655563/what-is-the-meaning-of-no-math-alphabet-change-to-frozen-version-normal-on-inpu
\def\@font@info#1{}
\setcounter{localmathalphabets}{0}





% Andrew Stacey's code from
% https://tex.stackexchange.com/a/50054/3954
%\makeatletter
%\tikzset{%
%	remember picture with id/.style={%
%		remember picture,
%		overlay,
%		save picture id=#1,
%	},
%	save picture id/.code={%
%		\edef\pgf@temp{#1}%
%		\immediate\write\pgfutil@auxout{%
%			\noexpand\savepointas{\pgf@temp}{\pgfpictureid}}%
%	},
%	if picture id/.code args={#1#2#3}{%
%		\@ifundefined{save@pt@#1}{%
%			\pgfkeysalso{#3}%
%		}{
%			\pgfkeysalso{#2}%
%		}
%	}
%}

%\def\savepointas#1#2{%
%	\expandafter\gdef\csname save@pt@#1\endcsname{#2}%
%}
%
%\def\tmk@labeldef#1,#2\@nil{%
%	\def\tmk@label{#1}%
%	\def\tmk@def{#2}%
%}
%
%\tikzdeclarecoordinatesystem{pic}{%
%	\pgfutil@in@,{#1}%
%	\ifpgfutil@in@%
%	\tmk@labeldef#1\@nil
%	\else
%	\tmk@labeldef#1,(0pt,0pt)\@nil
%	\fi
%	\@ifundefined{save@pt@\tmk@label}{%
%		\tikz@scan@one@point\pgfutil@firstofone\tmk@def
%	}{%
%		\pgfsys@getposition{\csname save@pt@\tmk@label\endcsname}\save@orig@pic%
%		\pgfsys@getposition{\pgfpictureid}\save@this@pic%
%		\pgf@process{\pgfpointorigin\save@this@pic}%
%		\pgf@xa=\pgf@x
%		\pgf@ya=\pgf@y
%		\pgf@process{\pgfpointorigin\save@orig@pic}%
%		\advance\pgf@x by -\pgf@xa
%		\advance\pgf@y by -\pgf@ya
%	}%
%}
%\newcommand\tikzmark[2][]{%
%	\tikz[remember picture with id=#2] {#1;}}
%\makeatother
%% end of Andrew's code
%
%\newcommand\ShadeCell[4][0pt]{%
%	\begin{tikzpicture}[overlay,remember picture]%
%		\shade[#4] ( $ (pic cs:#2) + (0pt,1.9ex) $ ) rectangle ( $ (pic cs:#3) + (0pt,-#1*\baselineskip-.8ex) $ );
%	\end{tikzpicture}%
%}%
%
%\newcommand\HatchedCell[4][0pt]{%
%	\begin{tikzpicture}[overlay,remember picture]%
%		\fill[#4] ( $ (pic cs:#2) + (0,1.9ex) $ ) rectangle ( $ (pic cs:#3) + (0pt,-#1*\baselineskip-.8ex) $ );
%	\end{tikzpicture}%
%}%





% https://tex.stackexchange.com/questions/161338/can-cleveref-be-made-to-use-the-oxford-comma-for-multiple-citations
\newcommand{\creflastconjunction}{, and\nobreakspace}


%		\tikzset{every state/.style={minimum size=1pt}}


\newcommand{\tikzcell}[1]{%
	\begin{tikzpicture}[
		baseline=(current bounding box.center),
		every node/.style={align=center},
		inner sep=0pt,
		outer sep=0pt,
		overlap/.style={draw, minimum height=6mm, minimum width=1.6cm},
		normal/.style={draw, minimum height=6mm, minimum width=1.6cm}
		]
		#1
	\end{tikzpicture}%
}


% Mark thesis as confidential
% \usepackage[angle=0,fontsize=16pt,color=black,pos={20.5cm, 0.5cm} ,anchor=rt,alignment=r,text=TAJEMNICA\ PŁ\\Egz.\ nr\ X]{draftwatermark}

% Highlight thesis author in bibliobraphy
% \usepackage{xstring}
% \let\originalbibitem\bibitem
% \def\bibitem#1#2\par{%
	%   \noexpandarg
	%   \originalbibitem{#1}
	%   \StrSubstitute{#2}{I.Nazwisko}{\textbf{I.Nazwisko}}\par}





\newcommand{\cg}{\cellcolor{gray!50}}
%\newcommand{\cd}{\cellcolor{gray!20}}
\newcommand{\cd}{}


\def\ygap{2}
\def\xgap{2}
\def\arrowgap{.125}
\tikzset{
	arc center/.is choice,
	arc center/left/.code={\def\isLeft{1}\tikzset{arc to}},
	arc center/right/.code={\def\isLeft{0}\tikzset{arc to}},
	arc center offset/.code={\def\centerOffset{#1}\tikzset{arc to}},
	arc direction/.is choice,
	arc direction/cw/.code={\def\isCw{1}\tikzset{arc to}},
	arc direction/ccw/.code={\def\isCw{0}\tikzset{arc to}},
	arc to/.style={
		to path={
			[evaluate={
				coordinate \u, \v, \m, \mv, \mvOrtho, \c, \cu, \cv;
				\u = (\tikztostart);
				\v = (\tikztotarget);
				\m = (\u)!0.5!(\v);
				\mv = (\v)-(\m);
				if \isLeft then {
					\mvOrtho = (-\mvy, \mvx);
				} else {
					\mvOrtho = (\mvy, -\mvx);
				};
				\c = (\m)+(\mvOrtho);
				\c = (\m)!\centerOffset!(\c);
				\cu = (\u)-(\c);
				\cv = (\v)-(\c);
				\radius = veclen(\cu);
				\startAngle = atan2(\cuy, \cux);
				\endAngle = atan2(\cvy, \cvx);
				if \isCw then {
					if \startAngle < \endAngle then {
						\endAngle = \endAngle - 360;
					};
				} else {
					if \startAngle > \endAngle then {
						\endAngle = \endAngle + 360;
					};
				};
			}]
			arc[radius=\radius pt, start angle=\startAngle, end angle=\endAngle] \tikztonodes
		}
	},
	vertex/.style={circle, inner sep=1pt},
	left loop/.style={arc center=left, arc direction=cw, arc center offset=10pt},
	right loop/.style={arc center=right, arc direction=ccw, arc center offset=10pt},
	invisible/.style={text opacity=0},
	-{Stealth[round, flex=1]}
}








\newcommand{\decisionproblem}[3]{%
	\par\addvspace{.5\baselineskip}
	\noindent
	
	\begin{tabularx}{\textwidth}{@{\hspace{\parindent}} l X c}
		\multicolumn{2}{@{\hspace{\parindent}}l}{#1} \\
		\textsf{Input:}    & #2 \\
		\textsf{Question:} & #3
	\end{tabularx}
	
	\par\addvspace{.5\baselineskip}
}



\NewDocumentCommand{\VertexChain}{O{} m m m m m}{%
	% #1 = optionaler Startknoten
	% #2 = prefix
	% #3,#4,#5,#6 = Indizes
	
	% --- Start ---
	\IfBlankTF{#1}{%
		\node[state,label=above:$v_{#3}$] (#2-#3) {};
	}{%
		\node[state,label=above:$v_{#3}$, right=of #1, xshift=0.5em] (#2-#3) {};
	}
	
	\node[right=of #2-#3] (#2-d1) {$\cdots$};
	
	\node[state,
	label={[xshift=-7pt]above:$v_{#4}$}, 
	right=of #2-d1
	] (#2-#4) {};
	
	\node[state, label=above:$v_{{#4}+1}$, right=of #2-#4] (#2-#4n) {};	
	
	\node[right=of #2-#4n] (#2-d2) {$\cdots$};
	
	\node[state,
	label={[xshift=-7pt]above:$v_{#5}$}, 
	right=of #2-d2] (#2-#5) {};
	
	\node[state,label=above:$v_{{#5}+1}$, right=of #2-#5] (#2-#5n) {};
	
	\node[right=of #2-#5n] (#2-d3) {$\cdots$};
	
	\node[state,
	label={[xshift=-7pt]above:$v_{#6}$},
	right=of #2-d3
	] (#2-#6) {};
	
	\path[->]
	(#2-#3) edge (#2-d1)
	(#2-d1) edge (#2-#4)
	%(#2-#4) edge (#2-#4n)
	(#2-#4n) edge (#2-d2)
	(#2-d2) edge (#2-#5)
	%(#2-#5) edge (#2-#5n)
	(#2-#5n) edge (#2-d3)
	(#2-d3) edge (#2-#6);
}

\newcommand{\thesistitle}{The Set of Synchronizing Words and the Constrained Synchronization Problem}

% https://tex.stackexchange.com/questions/20576/passing-newenvironment-parameters-to-the-end-block
\NewDocumentEnvironment{proofitem}{m}
{%
	\noindent\emph{Proof of #1:}%
}
{%
	\emph{[End, Proof of #1]}%
}

% Gives naming conflict with Springer Mono Style. Thus imported here.
\usepackage{complexity}

\usepackage[
twoside=true,
left=3cm, % inner
right=2cm, % outer
top=2.5cm,
bottom=2.5cm,
]{geometry}
\linespread{1.5}



%\linespread{1.5}
\onehalfspacing

\titleformat{\chapter}{\LARGE\bf}{\thechapter}{25pt}{\LARGE\bf}
\titlespacing*{\chapter}{0pt}{-35pt}{10pt} %{-50pt}{10pt}

\titleformat*{\section}{\Large\bf}
\titleformat*{\subsection}{\large\bf}

\renewcommand*\rmdefault{ptm}
\renewcommand*\sfdefault{phv}
\renewcommand*\ttdefault{pcr}

\definecolor{mygray}{rgb}{0.97,0.97,0.97}
\lstset{basicstyle=\small\ttfamily, frame=single, backgroundcolor=\color{mygray}}
\urlstyle{same}

\widowpenalty 1800
\clubpenalty 1800

\setcounter{tocdepth}{2}
%\setcounter{secnumdepth}{2}

\graphicspath{ {} }

\setlength{\headheight}{14.49998pt}
\addtolength{\topmargin}{-2.49998pt}



% style for headers and footers
\newcommand{\setfancyhf} {
	\fancyhf{}
	%\fancyhead[r]{{\footnotesize\leftmark}}
	%\fancyhead[l]{{\footnotesize\rightmark}}
	\fancyhead[RO]{{\footnotesize\nouppercase{\leftmark}}}
	\fancyhead[LE]{{\footnotesize\nouppercase{\rightmark}}}
	\fancyfoot[RE, LO]{}
	\fancyfoot[C]{\thepage}
}

%%https://tex.stackexchange.com/questions/267362/texstudio-how-to-change-color-of-internal-pdf-viewer
%\pagecolor[HTML]{FDF6E3}
%\pagecolor[HTML]{FD0000}
%\pagecolor{black}
%\color{white}% set the default colour to white

% TeX-Probleme die auftraten:
%https://tex.stackexchange.com/questions/7953/how-to-expand-texs-main-memory-size-pgfplots-memory-overload
% >pdflatex.exe --extra-mem-bot=10000000 -synctex=1 -interaction=nonstopmode thesis-trier.tex


%\includeonly{part3/part3}



%\tikzset{>={Stealth[length=2mm,width=2mm,scale=1.2,inset=0.2mm]}}
\tikzset{>={Stealth[scale=1.3,inset=0.2mm]}}






\newcommand{\pathmain}{./..}

\theoremstyle{plain}
\newtheorem{theorem}{Theorem}[section]
\newtheorem{proposition}[theorem]{Proposition}
\newtheorem{lemma}[theorem]{Lemma}
\newtheorem{claim}[theorem]{Claim}
\newtheorem{conjecture}[theorem]{Conjecture}
\newtheorem{corollary}[theorem]{Corollary}
\newtheorem{definition}[theorem]{Definition}

\theoremstyle{remark}
\newtheorem{remark}[theorem]{Remark}
\newtheorem{example}[theorem]{Example}

\makeatletter
\newcommand{\problemtitle}[1]{\gdef\@problemtitle{#1}}% Store problem title
\newcommand{\probleminput}[1]{\gdef\@probleminput{#1}}% Store problem input
\newcommand{\problemquestion}[1]{\gdef\@problemquestion{#1}}% Store problem question
\NewEnviron{problem}{
	\problemtitle{}\probleminput{}\problemquestion{}% Default input is empty
	\BODY% Parse input
	\par\addvspace{.5\baselineskip}
	\noindent
	\begin{tabularx}{\textwidth}{@{\hspace{\parindent}} l X c}
		\multicolumn{2}{@{\hspace{\parindent}}l}{\@problemtitle} \\% Title
		\textsf{Input:} & \@probleminput \\% Input
		\textsf{Question:} & \@problemquestion% Question
	\end{tabularx}
	\par\addvspace{.5\baselineskip}
}
\makeatother

% redefine part command to allow quotations
\makeatletter
\newcommand{\quoteline}[1]{\def\@quote{#1}}
\newcommand{\quotef}[2]{
	\begin{flushright}
		\rightskip=1.8cm\textit{#1} \\
		\vspace{.2em}
		\rightskip=.8cm#2
	\end{flushright}
}
\def\@part[#1]#2{%
	\ifnum \c@secnumdepth >-2\relax
	\refstepcounter{part}%
	\addcontentsline{toc}{part}{\thepart\hspace{1em}#1}%
	\else
	\addcontentsline{toc}{part}{#1}%
	\fi
	\markboth{}{}%
	{\centering
		\interlinepenalty \@M
		\normalfont
		\ifnum \c@secnumdepth >-2\relax
		\huge\bfseries \partname\nobreakspace\thepart
		\par
		\vskip 20\p@
		\fi
		\Huge \bfseries #2\par
	}	
	\vspace{1.8cm}
	\ifdefined\@quote
	\@quote
	\def\@quote{}
	\else
	\fi   
	\@endpart}
\newcommand{\partwithquote}[3]{
	\quoteline{
		\quotef{#2}{#3}
	}
	\part{#1}
}
\makeatother




\begin{document}
	
\tableofcontents

%\clearpage
%
%- cocoon paper diskutieren
%- eine NP-hardness Bedingung mit den weakly aux state notions bzw wie die Pspace und NL-Bedingung uv^*w und u muss zurücksetzen


\subsection{Beachten und Formales}

\todo[inline]{ \footnotesize 
	
	- fehlender referenzen/fragezeichen
	
	- einige links im Index funktionieren nicht, z.B. "right coset". sieht man aber -> kein grüner Kasten. einige einträge im index keine hypoerlinks	
	
	- Überschriften in Figure-Liste "list of figures" nochmal überarbeiten. Einige sind zu lang
	
	- Uni-Affiliations von Henning und Igor adden? Oder doch nicht, siehe Vorgabe Deckblatt
	
	- Paragraphen durch Zeilenabstand oder durch Einrücken kennzeichnen?
	
	- eidesstaatliche Erklärung + wissenschaflticher Lebenslauf
	
	- dois in bib adden.
	
	- prüfen, dass alle kapitel "Discussion"-Section haben
	
	- nach comma slices gucken. 
	
	Im Notation Index:
	- Seite 7 unten in der Mitte, danach Seitennummern oben links/rechts?
	
	- oben steht list of tables auf den seiten???
	
	- Generall Seitenzahlen noch einmal prüfen. Die nummerierung der seiten haut nicht hin, am anfang 1-14 und dann nochmal von 1
	
	- Notation Index sortieren bzgl. Seitenzahlen
	
	
	- Im introductory Teil die eigenen Referenzen kürzer zitieren? Nur titel + author? irgendwie abgekürzt zitieren, ohne links und vollem konferenznamen?
	%	
	%	% https://www.google.com/search?q=cite+vs+fullcite&rlz=1C1GCEU_deDE823DE823&oq=cite+vs+fullcite&aqs=chrome..69i57j33i160l2.2398j0j7&sourceid=chrome&ie=UTF-8
	%	% https://tex.stackexchange.com/questions/159286/biblatex-how-to-fullcite-the-first-occurences-and-cite-the-others
	%	% https://github.com/jgm/pandoc/issues/2335
	%	
		
	- partielle Funktionen wirklich nicht vor Part IV benutzt (z.b. bei Definition von circuits)?
	
	- nutzt man bei cref eigentlich tilde oder nicht?
	
	- Markos wie semiautomat rausnehmen?	
	
	- warum eigentlich hier $u_i$ und in anderem theorem $v_i$ genannt? solche sachen vereinheitlichen
	
	- nachtmodus wegnehmen
	
	- drucken (vordrucken mit pdf-drucker 2-seitig)
		
	- lebenslauf einfügen	
	
	- bei der reihenfolge der konkatenation überall nochmal aufpassen. 	
	
	- spellchecker drüber laufen lassen

	- ldots zu cdots in überall in formeln!	
	
	- notation proper prefix, $pref^+(L)$ vielleicht? im auxiliary state avoiding teil brauche ich das oft.
	
	- notation max(x,y) vs $\max\{x,y\}$
	
	- notation $A^{\ge n}$ einführen?
		
	- \cref{constr_sync:thm:prefix_before}: achtung, zwei labels
	
	- \cref{constrsync:sec:aux}: in den sätzen oben mit w statt v bezeichnet
	
	- Problemname im index? oder eigener index für probleme. Irgendwo auch index aller problem? nach notation index? Abkürzungen wie PDFA auch in index. Index of open problems und conjectures.
	
			
	- zitierungen sortieren
		
	- notation alph(vi) einführen für buchstaben in wort?
	
	
	- style von paper Computational complexity of problems for deterministic presentations of sofic shifts , cai. bilder von reduktionen wie in symbolic dynamics problems paper.
	% https://arxiv.org/pdf/2112.03484
	 
	
	- bei den benutzungen von aussage nicht nur "by Cor X" zitieren, sondern kurz aussage im text nennen?
	
	
	- mal discussion and conclusion, mal nur discussion genannt, die abschlusssektion... vereinheitlichen. Oder "Related Work" nennen?
	
	
	- equalities, formulas, equations. wann welchen begriff nehmen?
}

\subsection{Konventionen}

\begin{adjustwidth}{-1.8cm}{-1.8cm}
	
	\footnotesize 
	\begin{enumerate}[itemsep=-5pt]
		
		\item Oxford-Comma	
		
		\item one-point singular idempotent  oder  idempotent one-point singular. vereinheitlichen. auch für partially ordered binary automata oder binary partially ordered. welche reihenfolge nehmen? auch: binary completely dist oder complete dist binary, einheitlich ob alpahbet zuerst oder nicht, vergleichen mit anderen stellen im dokument. english? defect one transformation vs transofrmation of defect one.
		überalle so rum commutative partially ordered nicht partially ordered commutative 
				
		\item Überall "tight bound" durch "sharp bound" ersetzen.
		
		\item statt $\eqreach$ umgekehrt, unten der geschwungene strich?
		
		\item sprechweise: "proved in Theorem/Lemma/Proposition" -> ist aber nur statement, gezeigt aber im beweis also "shown in the proof of Theorem/Lemma/Proposition" besser?
		
		\item outdegree und indegree gucken überall ohne bindestrich geschrieben.
		
		\item in Bildern von Automaten Zustände auch überall grau hinterlegt und einheitliches aussehen. Schatten in einigen bildern unter zuständen fehlen.		
		
		\item betonungszeichen, auch bei cerny. akzente bei araujo üerall besetzt?			
		
		\item $a,b,c,d$ statt $p,q,s,t$ für punkte?
		
		\item machmal s,t, manchmal p,q für zustände benutzt, vor allem in "more notions from autoamta theory"-teil
		
		\item vereinheitlichen non-trivial, non-empty, non-deterministic statt nontrivial, nonempty, nondeterministic; nondisjoint, non-emtpy durch nonempty ersetzen? labelled vs labeled?
		
		\item bei TFAE-Aussagen Punkt am Ende von Items setzen, vereinheitlichen.  Kommas am ende von formelauflistungen auch an anderen stellen weg (konventionen: keine satzsachen bei eingerückten formeln, oxford comma, cdots für operationen, ldots für auflistungen), oder frei in punctuation in eingerückten formeln. Das irgendwo sagen? (hatte mal nachhilfeschüler die verwirrt waren)
		%\todo{erwähnen keine punkte/kommata am ende von fallunterscheidungen.}
		%\todo[inline]{theorem: zentrale aussage des kapitels,
			%	proposition: eine aussage die für sich steht, aber nicht so wichtig ist,
			%	lemma: tehcnisches ergebnis, wird in beweisen von, meist theoremen, benutzt.}
		%
		% 	\todo{mal $:$ mal $\mid$ bei set builder. je nachdem ob formel $\forall : etc$ folgt oder betragsstriche }
		
		\item sink durch trap überall ersetzen	
		
		\item bei Sprechweise aufpassen "Problem in NL/NP usw", aber ein Algorithmus ist nicht in NL/NP, aber läuft auf einer nondet. TM in log-space/P-time oder so, oder "Algorithm witnesses das Problem" sagen oder so
		
		\item einheitlich von 0 oder 1 zustände benennen. notation $[n]$ bei 0 starten lassen oder nicht.
		
		\item prefix in labels in Part 3 ändern. nicht mehr $sync_max$, sondern $aut_sync_words$ oder so? Und generell
		laben-Namen vereinheitlichen
		% https://tex.stackexchange.com/questions/289038/how-do-i-show-the-labels-of-theorems-definitions-etc
		
		\item proved or proven?
		
		\item vereinheiltlichen family oder class of sets sagen 	
		
		\item Sections oder Paragraphen in Diskussion-Teilen?
		
		\item holds true kann man schon sagen? englisch?
		
		\item form, type, structure? english
		
		\item relativ unkonsequent manchmal enum umgebungen, machmal description,machmal nur fliesstext um definitionen einzuführen. vereinheitlichen?
		
		\item $[n]$ durch $X$ ersetzen und sagen X finite? gucken wo ich $[n]$ durch $X$ ersetzen kann.
		
		\item hier reset letter/word einführen? dann bräuchte man nicht mehr synchronizing 
		
		\item written oder denoted by bei Definitionen schreiben?
		
		\item Konjugation $t^g$ weiter vorne definieren? Ist ja jetzt kein genuines Permutationsgruppenkonzept.
		
		\item mit $R(A)$ statt $I(A)$ bezeichnen für Rystsov graph?	
		
		\item bei transforamtion matrix nciht adjacenzmatrix nennen, weil kein graph. Nutze ich Adjanzenmatrizen irgendwo?
		
		\item A,B für zustandsmengen, S,T für semigruppen! aber S,T auch für SCCs benutzt
		
		\item aufpassen, nicht $S$ für mengen und semigruppe benutzen! deswegen bei semigruppen immer teilmengen von domain A/B/C...
		
		\item reihenfolgen in zitierungen sortieren, also dass in cite-Befehl alles aufsteigend sortiert erscheint?
		
		\item set builder notation mal $\mid$ mal $:$, erstes bevorzugen wegen formel mit : i def, aber wenn betragsstriche dann anderes
		
		\item schreibt man "... such that" oder "... be such that"? was ist besseres englisch?
		
		\item kapitel sparse: w oder u vereinheitlichen. Sprachen in $w^*$
		
		\item statt regular rational sagen?	 
		
		\item context-free language definieren?	 
		
		\item P statt C für prefix code schreiben? \cref{constr_sync:thm:prefix_before}: P als namen eher vermeiden, da meist zustände von constraint automat mit P bezeichnet
		
		\item appendix vor literaturverzeichnis?
		
		\item überall Item vorne dranschreiben, wenn man sich auf Unterpunkte von Aussage bezieht. In rank factorization paper machen die das auch so
		
		\item minimalautomat auch $\mathcal A_{\equiv_L}$ statt $\mathcal A_L$?
		
		\item schlecht, auf spätere ergebnisse zu verweisen...
			
		\item gucken, dass überall $\Sigma$ und $Q$ in richtiger reihenfolge bei automatonspezifikation. besondern im letzten teil, bei pto als input.		
		
		\item sprechweisen We state/prove/show, konsequent unterscheiden, erstes nur aussage, letzte mit beweis?		
		
		\item submonoid of Sigma-star oder nur monoid sagen?
		
		\item SI-autoamta oder automata with simple idempotents -> vereinheitlichen (siehe auch captions, letzter teil vs vorne)
		
		\item bei mehreren Autoren "und" statt "ambersand": Also Shyr und Yu statt Shyr \& Yu. Bianca hat da auch mal eine Mail zu geschrieben. ChatGPT sagt auch eher "und".
		
		\item referenzen, namen von journalen und konferenzen alle abkürzen?
		
		\item In proof of \cref{constrinput:pto:lem:Sigma1starSigma2Sigma3star} clauses c klein schreiben? woanders groß. vereinheitlichen
		
		\item Sec. oder Sect. als Abküzrung (in zitierungen)? kommando dafür. 
		%\todo[inline]{für "Thm." usw in cite-befehlen eigens kommando, THMCITE oder so, auch in align gleichungsketten}
				
		\item makro! nconcat für n-fache konkatenation mit pünktchenschreibweise? 
		
	\end{enumerate}
\end{adjustwidth}


\subsection{Noch machen}


\begin{adjustwidth}{-1.8cm}{-1.8cm}

\footnotesize 
	\begin{enumerate}[itemsep=-5pt]
	\item Diagramm Kapitelabhängigkeiten (draw.io, chatgpt + tikz). Wie in Martin Davis Computability Lehrbuch
		
	\item spezielle notation für partielle funktion, so mit gekreuztem pfeil oder so? dann in nomenclature auch erwähnen? 
	
	\item  begriff uniformly minimal einführen?
	
	\item  pfeilspitzen in manchen bildern größer
	
	\item runtime vs running time? Und "algorithm that runs" oder "algorithm running"?
	
	\item  prüfen, dass alle decision problems mit Namen auch im Index auftauchen
	
	\item Kommando "semiautomat" durch "semiautomaton" oder "semiaut" ersetzen. Und auch "automat" durch "automaton"?
	
	\item applyfunc und concat überall konsequent nutzen
	
	\item notation für orbitalgraphen?
	
	\item kommt begriff concept irgendwo vor?
	
	\item brauche ich reset automaton irgendwo, ist verwirrend mit synchronizing automaton?	
	
	\item In Index steht cyclic automaton, SI-automaton, strongly connected automaton, usw separate Einträge. oder automaton-Eintrag und die Typen drunter?	
	
	\item In "classes of automata"-chapter section strongly connected and strongly trap-connected automata?    
		
	\item Primitive Matrizen diskutieren?
		
	\item Grundidee des P-time Algorithmus von Szykula/Ferens skizzieren
	
	\item Namen für decision problem {\sc DFA-Distingusish} oder so? 
	
	\item schönere tcolorbox statt fbox nehmen?
	
	\item die autoamtenfamilien im slowly-sync-teil mit mathscr und nicht mathcal? (dann keine überschneidung mit alternating group)	
	
	\item  besseren nconcat befehl nutzen für n-fache composition, und umbennnene in comp wegen composition statt concat; ist ja keine concatenation...
		
	\item nochmal schauen, alles in aschbacher style schreiben. oder downey/fellows
	%\todo[inline]{genauer, in section so und so; nochmal mit Aschbacher und Harary vergleichen}
	% \todo[inline]{ist das eigentlich üblich, einführung mit nicht-definierten begriffen zu schreiben und diese erst im kapitel einzuführen.
		% 	aschbacher macht es, harary? und in artikeln?
		% }
	% 
		
	\item ack,abstract,zusammenfssung auch in table of contents erwähnen, wie in Ryzhikov diss
	
	\item die Autoamten $A_t$ und $A_c$ am Anfang einführen und nicht erst in \cref{constrsync:prop:SI-binary-sync-in-L} ?
	
	\item die bezeichner a,b,c,d in der reihenfolge aussage/beweis in den bedingugnen für transitivität (Lemma~\ref{prelim:lem:2sets-transitive} und folgende in permutation groups-section) vereinheitlichen	
	
	\item circular in subclasses teil von prelim einführen, dann kann ich auch ein wenig mehr zu denen sagen (martyugin arbeiten, pin/dubuc)	
		
	\item Section "Synchronizing Autoamta and the Reachabiltiy Order", teile davon verallgemeinerung die Ergebnisse für trap-state automata davor. Vielleicht kapitel anders anordnen um dise Ergebnisse in trap-state Kapitel nutzen
		
	\item doppelpfeile an automaten für hin und zurück irgendwo erklaeren bzw einführen. kommt bei kleinen familien D,E usw oft vor (am anfang, wo state diagrams eingeführt, aber dann auch beispiel dazu? noch neben cerny anderen automaten)
	% https://tex.stackexchange.com/questions/445108/drawing-a-double-headed-latex-arrow
	
	\item wie steinberg primitive monoide definieren. primitive automaten einführen? % dann macht die definition von primitiven gruppen ohne transitiv auch mehr sinn (bzw. brauch ich di egar nicht extra definieren, aber bemerkung, dass hier ncht transitiv wäre gut)
		
	\item die Gleichungen aus Formal Languaees Teil nach Part IV verschieben?	
	
	\item Aussagen bei Einführung vom Power Automaton für beliebige Transformationen formulieren?
	
	\item für relationen auch transition matrix definieren? Wird bisher nur kurz in \cref{prelim:sec:remarks_func_comp} erwähnt.
		
	\item Sprechweise permutes SCCs, betonen wenn sets gemeint (siehe Hennings bemerkung von S.122 (92) )	
		
	\item auch nochmal später schauen, dass nicht G(A) = I(A) geschrieben (der Automatengraph hat ja noch ganz viele self-loops). Bemerkung dazu
	
	\item kapitel kommutative automaten: wie wäre es,w enn ich nithct immer fixed, sondern einfach even unary scheribe?
	
	
	\item \cref{prelim:remarks_trans_semigrp_semiautomata} umformulieren, semiautomat nicht bloß transformation semigroups mit gen set (faithful)
	%
	% -bzgl. beziehungs transformations semigroups vs automaten: vielleicht automaten immer mit start und final und semiautomaten komplett rauslassen? nur von transformatio
	%	semigroups sprechen? wobei bei automat nicht der quotient in Sigma star, also keine faithful darstellung unbedingt.
		
	% hier auch so ein graph am anfang:
	% https://wiki.epfl.ch/provenance2011/documents/foundations+of+databases-abiteboul-1995.pdf	
	
	\item in \cref{constr_sync:subsec:morphisms} ganz oft für sigma nicht applyfunc benutzt!
	
	\item kurzschreibweise, regex für mengen bzw für ausdrücke. ich mische $\{\cdot\}$ und ohne klammern notation. also regex und sprache z.B. a(a+b)* vs $a\{a,b\}^*$ einheitlich, $\{a,b\}^*$ oder $(a+b)^*$
	
	\item circular automata doch schon früher einführen weil schon in \cref{prelim:compl:sec:discussion} erwähnt.
	
	\item PDFAs einhetlich mit mathcal B? aber eingeüfhrt mit mathcal A statt mathcal B?
		
	\item Konkrete beispiele codes abschnitt: siehe journal artikel cocoon, version 30.06.2023 
		
	\item übersichtstabelle small state constraint PDFAs: bei rot striche nutzen, und bei grün punkte, wegen farbsehschwächen und s/w druck
	
	\item nutze square-and-multiply argument/algo an einigen stellen, z-.B. im Beweis von \cref{sparse:thm:bounded_in_NP}. Denn irgendwo am anfang erklären?
	
	\item auch nochmal die alten mails von nugant durchgucken
	
	\end{enumerate}
\end{adjustwidth}




%\setlength{\marginparwidth}{13cm}
\todo[inline]{ \footnotesize 
	
	- Im Abschnitt über bounded languages: oben w genannt, unten wieder $u$? vereinheitlichen?
	
	- Im Beweis von \cref{constrsync:thm:sparsemonoid-regular} (teil growth of languages): Bezout in allgeminer form, wie der shallit pighizinni satz; vorne zum zitieren. und die fussnote zu frobenius coin problem in dem beweis als eigenen text nach statement
	
	- probleme mit cref auf arbeitsrechner. erwähnen vorm einschicken
		
	- when bounded by a polynomial oder if bounded? also grammatik when vs if im englischen?
		
	- neg oder overline um negation von booleschen Variablen auszudrücken?
		
	- konstanten in circuits kann man nur mit mindestens einer eingeabe erzeugen, in section zu circuits erwähnen?
		
	- We can assume $s \ne t$ and that~$s$ has no incoming edges and~$t$ has no outgoing edges.
	vielleicht das immer sagen und auf k-DAG-ST verzichten.
	Gucken wo ich das alles brauch. falls nur 1x dann ja
	
	
	- in def solvable taucht auch notation mit anwendung auf teilmenge auf, aber das ist mehrdeutig, weil da nicht
	die erweiterung auf mengen gemeint.
	
	
	% permutes the set of vs permute each strongly connected 
	
	
	
	- vielleicht extra notation k-SFS, oder Sigma-SFS? nutze es für fixe alphabet im beweis von \cref{constr_sync:thm:introductory_examples}
	
	
	- nach fehlenden referenzen, fragezeichen im dokument suchen
	
	% kapitelabhängigkeiten
	% https://webhomes.maths.ed.ac.uk/~v1ranick/papers/cohnalg1.pdf
	
	
	
	% In Überschriften (Semi-) zu (Semi). Und dort, wo es eingeführt wird.
	%
	% https://www.reddit.com/r/asklinguistics/comments/7jgtgd/what_is_the_proper_or_closest_method_to_indicate/
	%	https://www.google.com/search?q=how+to+mark+optional+parts+of+words+in+english&rlz=1C1GCEU_deDE823DE823&oq=how+to+mark+optional+parts+of+words+in+english&gs_lcrp=EgZjaHJvbWUyBggAEEUYOTIHCAEQIRigATIHCAIQIRigAdIBCTEwMDQ3ajBqN6gCALACAA&sourceid=chrome&ie=UTF-8
	%	https://www.reddit.com/r/asklinguistics/comments/7jgtgd/what_is_the_proper_or_closest_method_to_indicate/
	%	
	%	To mark optional parts of words in English, especially prefixes or suffixes, use parentheses like (un)happy or happi(ness), or square brackets for linguistic glossing like boy[NOM.SG]. For optional words or phrases within sentences, use ellipses (...) for omitted content or braces {} (though less common) for annotations, or just omit the word if the meaning remains clear, as in informal writing. 
	%	For Prefixes & Suffixes (Morphology/Linguistics)
	%	Parentheses (): Common for showing an optional part of a word.
	%	(un)do (meaning 'undo' or 'do')
	%	anti(biotic) (meaning 'antibiotic' or 'biotic')
	%	Square Brackets [] (Linguistic Glossing): Used to show elements that aren't overtly present or are optional in a specific context.
	%	puer-Ø (meaning 'boy' with an unmarked singular noun). 
	%	For Words & Phrases (Syntax/Writing)
	%	Ellipsis ...: Indicates omitted words or phrases, common in quotes or informal speech, signaling something is left out.
	%	"Four score and seven years ago... our fathers...".
	%	Parentheses (): Can enclose optional phrases or extra, non-essential information in sentences.
	%	"The dog (a golden retriever) ran in the park.".
	%	Omission: Simply remove the word if the sentence still makes sense, common for adjuncts (optional adverbs/phrases).
	%	"John went to the store (yesterday)." -> "John went to the store.". 
	%	Other Methods
	%	Hyphens: Can sometimes link optional parts, though less standardized.
	%	Braces {}: Sometimes used in linguistic analysis to group optional elements. 
	
	%\todo{oder (semi)automaton oder (semi-)automaton?
		%https://english.stackexchange.com/questions/427661/how-to-correctly-hyphenate-a-phrase-with-clarifying-word-in-parenthesis
		%} 
	%
	
	
	%%%%%%%%%%%%%%%%%% Sonstige Bemerkungen %%%%%%%%%%%%%%%%%%%%%%%%%%%%%%
	
	% Ideas, on their own, are worthless
	% https://www.youtube.com/watch?v=5RFwX1g9bq8 
	
	% https://www.dickimaw-books.com/latex/thesis/html/biblatextroubleshooting.html
	
	% https://en.wikipedia.org/wiki/Fundamental_theorem_of_software_engineering
	
	% ://tex.stackexchange.com/questions/415395/what-is-causing-this-error-which-i-think-has-something-to-do-with-fonts
	
	
	
	
	
	
	
	%\todo{beweis über generating set von Sn ohne die relation \sim formulieren, direkt mit graph und pfaden}
	
	
	
	
	
	%	\todo{nfold product verwechslunggefhar, z.B. bei logarithmen, mit n facher mulitplikation
		%		
		%		%	https://en.wikipedia.org/wiki/Function_composition
		%	}
	%	
	
}




\todo[inline]{
	\footnotesize 
	As Ernest Hemingway famously said, “the only kind of writing is rewriting.” Few writers will produce a first draft that they're happy with. Through the process of creating a second and third draft, you will hone the language, cut out superfluous material and refine your manuscript
	
	in ack, hier ja oft neu geschreiben und paper neu geschrieben - jeder beweis does nto have much commonalities
	with everything pusblished by me so far on which
	this wordk is based. erweitert, vereinfacht, took the chance when writing this thesis to revise etc
	
	
	damit geduld der gutachter strapaziert
	
		
	prof müller-hannemann auch für verständnis danken
	
	paper ergebnisse, aber auch thesis selbst.
	i can safely say I haven written almost three different version of the present thesis until this (has-to-be) final version.
	
		
	
	- in ack, trotz revisting behaupte ich nicht, dass definite form erreicht
}

%https://tex.stackexchange.com/questions/103608/how-to-force-latex-not-to-break-the-line-after-a-hyphen
%https://tex.stackexchange.com/questions/9691/avoid-hyphenation-in-2-d



\subsection{Ideen, nicht so wichtig erstmal}

\todo[inline]{ \footnotesize

	- Noch die Subset Mapping Problems für SI-Automaten untersuchen. Für binäre sollte es machbar sein. In dem Fall erreichbare Mengen klassifizieren.

	- vielleicht noch anhang, further notions of synchronization mit NEA, partialle, parikh automata, ...
	
	- die fakten über endormorphismen in morphisme-teil mit bildern untermauern?
	
	-  Aussagen über Morphismen für gerichtete Multigraphen statt nur  Automaten formulieren, z.B. strongly connected auf strongly connected für morphismem, oder: ein Endomorphismus eines endlichen, stark zusammenhängenden gerichteten
	%Multigraphen ist ein Isomorphismus
	
		
	- reguläre ausdrücke über universelle konstruktion einführen? darum beweise mit sprachen für gleichungen
	
	- wer hat free monoid iegentlich zuerst untersucht, zeichenketten? thue? 
	quinte arbeit in lit/free-monoid ordner, tarski?
	% https://en.wikipedia.org/wiki/Concatenation_theory
	die leute, die ersetzungssysteme zuerst untersucht haben?
	

}


% beispiele solid code aus meinem FCTCS'21 paper nochmal angucken

% levis lemma bei formal langauge vielleicht erwähnen?

% noch doppelte labels bei statements suchen

% kapitel "more on prefix codes" umbenennen in factorizations and prefix codes und den PDFA-teil wieder raus nach vorne

% COCOON paper besprechen, dort die andere relation benutzt sagen für konstruktionen

% für two state strongly connected constraint PDFA  NL-c folgt einfach, wenn zustand wo eine transition undefiniert, aber an anderem definiert und geht
% rüber, dann aux state avoding mit dem am ende, also so start und final wählen, dass einer so drin ist

% - freie NOT-gates vs nur an eingängen, wenn ausgabe von teilzierkuit und negation benutzt, muss man mitunter zwei anlegen, also faktor 2

% begründen, warum feinere reduktionsbegriffe benutzt?


%\todo[inline]{integer/non-negative integer statttt number sagen?}


% gucken, dass alle probleme, die benutzt werden, auch irgendwo definiert sind




%%%%%%%%%%%%%%%%%
%% War vorher in Automata and Encodings-Teil

%	Gerne — hier ist ein **sauber formulierter, paper-tauglicher Fußnotentext** in gutem akademischem Englisch:
%
%> *One possible way to encode a list of binary strings as a single binary string is to use a simple escape-based scheme: encode each bit `0` as `00` and each bit `1` as `11`, and use `01` as a separator between successive strings. Since the separator does not occur inside any encoded string, this yields a uniquely decodable representation. Such an encoding is easy to define and sufficient for our purposes; efficiency considerations are irrelevant here.*
%
%Wenn du es **noch knapper** oder **noch formeller** brauchst (z. B. für STOC/FOCS-Stil oder ein Buch), sag kurz Bescheid.
%
%
%
%
%
%
%Für ein **Buch** würde ich eine Variante empfehlen, die
%(1) zeitlos klingt,
%(2) sehr klar ist,
%(3) unnötige Details vermeidet,
%(4) und den Eindruck vermittelt, dass dies eine **konventionelle, unkritische technische Wahl** ist.
%
%Hier ist eine **buch-taugliche, ruhige und elegante** Fußnotenfassung:
%
%> *One convenient way to encode a finite list of binary strings as a single binary string is to use a simple escape-based scheme: each bit `0` is replaced by `00`, each bit `1` by `11`, and the block `01` is used as a separator between successive strings. Since the separator cannot occur within an encoded string, the resulting representation is uniquely decodable. We fix such an encoding once and for all and abstract away from its concrete details in the sequel.*
%
%Warum diese Version gut für ein Buch ist:
%
%* **„one convenient way“** → signalisiert Austauschbarkeit
%* **„fix such an encoding once and for all“** → klassischer Buchstil
%* **„abstract away from its concrete details“** → bereitet spätere Abkürzungen vor
%* Keine Rechtfertigungen zur Effizienz → wirkt souverän
%
%Wenn du möchtest, kann ich sie auch:
%
%* stärker **formal** (Lehrbuchstil),
%* oder etwas **weicher** (einführendes Kapitel),
%* oder explizit auf **Existenzargumente** zuschneiden.



%%%%%%%%%%%%%%%%%%%%%

%\todo[inline]{beispiele und ergebnisse nennen. % https://en.wikipedia.org/wiki/Circuit_complexity#History,
%	
%	wer war eigentlich der erste, der NC mit schaltkreisen definierte?
%	
%	in welchen paper wurde zuerst der zusammenhang mit parallelcomputern formalisiert und gezeigt?
%}


%%%%%%%%%

% - vielleicht bei *-operation discuss nur sagen kleene-star meint was anderes, operation auf sprachen.
%   mit bemerkung wie in kufleitner/volker-buch
% - und für zusammenhänge das in anhang packen


%\todo{kommando automat und minaut arbeiten nicht gut zusammen} %$\automat{{\minaut{{\normalfont L}}}}{Q_L}{\Sigma}{\delta_L}{[\varepsilon]_{\equiv_L}}{F_L}$




%%%%%%%%%%%%%%

% leonardo da vinci, art is never finished, only abandond



% \todo{bemerken sum_{emptyset} = 0 und in induktionsbeweisen $x_{i-1}x_i$ für i = 1 und keine i = 0 existent}





% pfeilspitzen wie in 
% https://arxiv.org/pdf/2205.09404



%\todo{silbentrennung prüfen, ist hier vielleicht manchmal komisch
%	
%	https://www.silbentrennung24.de/wort/aktuell/
%}




\todo[inline]{bei def primitive group: vielleicht das bild aus meinem hightlights oder von grohe vortrag? }



% - halle vortrag reingucken und vergleichen charakterisierung primitivität mit uniform minimal; meine formulierung gut? 



% - bilder vereinheitlichen. auch wenn zustände kleiner dargestellt grau mit schatten



%\todo{ uniform minimal einführen?}


% kapitelabhängigkeiten:
% https://www.mimuw.edu.pl/~malcin/book/parameterized-algorithms.pdf



% preface so aufbauen mit den teilen?
% file:///C:/Users/StefanH/AppData/Local/Temp/beb5948b-cb11-43e3-9a4a-cdc1919aff92_GTM-master.zip.f92/GTM-master/GTM258.Foundation.of.Optimization.pdf



% proper suffix/prefix definiert?

% proper prefix definiert? u mit uSigma^+ \cap L != emptyset









% binary completely reachable vs completely reachable binary; überall vereinheitlichen



	%state/.style={circle,draw,minimum size=8mm}
	
	
	
	
	
% noch sagen, root state, tail state, cycle state terminology von induzierten tranformationsn





% automatenabbildungen vereinheitlichen. manchmal mit schatten, manchmal ohne in einführungsteil





% - synonym cyclic permutation für cycle rausgenommen. benutzt ich glaube ich nirgends. nochmal gucken




% \crefitem im springer-format noch klammern um letzten verweis adden
% macros bedingt, auch partwithquote anpassen



% grafiken in springer-teil zu groß. statt cm vielleicht em-Angaben nutzen?


% die paragraphen sehen komisch aus im springer-format.




% im makefile noch einbauen, dass man nur einzelne kapitel erzeugen kann


% crefitem für springer anpassen



% https://dl.acm.org/doi/10.1145/103418.103468



%\todo[inline]{stattdessen sagen, nicht transition matrices als eingabe, sondern let transition matrices und auf diese section verweisen}

 

%\begin{comment}
	% Arial Font?
	% https://www.reddit.com/r/GradSchool/comments/tu0are/arial_or_times_new_roman_for_thesis_freaking_out/
	\begin{titlepage}
		\includegraphics[height=1.65cm, right]{Logo_Universitaet.png}
		\hfill
		
		\centering
		\vspace{120pt}
		
		{
		\LARGE \bfseries \setstretch{1.2}
		\thesistitle
		\par
	    }
		
		\vspace{12pt}
		{   \setstretch{1.1}
			\large
		Vom Fachbereich IV der Universität Trier zur Verleihung des akademischen Grades Doktor der Naturwissenschaften (Dr.~rer.~nat.)
		genehmigte Dissertation 
		\par
	    }
        
        \vspace{15pt}
        {\large 
        von
        }
    
        \vspace{15pt}
        {
        	\Large
        Stefan Hoffmann
        }
        \vspace{50pt}
        
        { \large 
        Trier, 2025
        }   
			
		\vfill		
		
		% Noch affiliation in Klammern dahinter?
		\raggedright
		{   \setstretch{1.1}
			\large 
			Betreuer: Prof. Dr. Henning Fernau \\ 
			1. Berichterstatter: Prof. (Associate) Igor K.~Rystsov  \\
			Datum der Disputation: 31.01.2025
			\par
		}
		\vspace{70pt}		
	\end{titlepage}
%\end{comment}	


\pagestyle{empty}
\cleardoublepage

% set empty style (no headers and footers)
\pagestyle{empty}
\pagenumbering{roman}

\begin{center}
	\bfseries
	Abstract
\end{center}
%!TEX root = ../thesis-trier.tex



Have you ever wanted to start all over again? Has your smartphone or computer ever frozen and you had to restart it? 
If so, then you are familiar with the idea of resetting things,
the central concept of the present work.  
More specifically, it is about resetting (also called synchronizing) 
finite state transition systems, also called finite automata. 
 Finite automata process input sequences (also called words) and can be used to model a wide variety of systems, for example,
switching circuits, networks, or simple
behavior patterns in artificial intelligence. 
In all these situations, knowing whether we can reset the system is of interest. A circuit component might be out of sync with its environment, can we reset it? 
A robot might be stuck, or we are simply unsure about its current state,
so can we reset its state of operation? 
However, while this problem is tractable
for completely specified and deterministic automata,
only a small deviation from this ``optimal'' situation, i.e., having only partial information, allowing nondeterminism, or constraining the admissible input
sequences, results in intractable
algorithmic problems.

Here, we are concerned with one particular problem variant: 
constraining the admissible input
sequences (the \emph{constrained synchronization problem}).
In this setting, what computational complexity can we deduce 
from the kind of constraint put on the input sequences and/or
types of input automata?
We state various results on how to manipulate sets of admissible synchronizing input sequences
so as to retain a certain complexity of the problem.
We also identify various conditions 
that either yield tractable or intractable problems.
More specifically, the following complexity classes will arise: 
$\AC^0$, \LOGSPACE, \NL, \PTIME, \NP, \PSPACE.
Except for~$\AC^0$ and~\PTIME, we 
give sets of admissible synchronizing sequences
yielding constrained problems that are complete for these classes.
Moreover, we investigate the following classes
of admissible sets of synchronizing sequences: finite, infinite, sparse, letter-bounded, commutative, or
constraints specified by certain small-state automata; where we give
a complete complexity classifications for the latter three classes.
Lastly, we  consider the following types of input automata to the problem:
monotonic, partially ordered, commutative, solvable, SI- and $\rho$--automata.
We obtain, for example, that for commutative input automata, the problem is always in \PTIME.

Now, the set of synchronizing 
sequences, i.e., our search space, of a given input automaton
is acceptable by a finite automaton as well.
This leads naturally to the investigation of accepting automata for sets
of synchronizing sequences, specifically for the question of when they
become largest possible.
Here,  we identify and investigate two conditions
that control the size of a minimal automaton
for the set of synchronizing sequences:
reachability and distinguishability.
For both notions, we establish various structural results and connections with concepts
from the literature such as, for example, completely reachable automata.
Decision problems and algorithmic questions 
that naturally arise
from these notions are also investigated.
We state conditions yielding extremal minimal automata accepting sets of synchronizing sequences.
Moreover, we identify a connection to permutation group theory, which 
is linked to automata theory by the group of units
of transformation monoids of automata.
Specifically, via graph-theoretical constructions,
we obtain new characterizations of primitive permutation groups
based on associated synchronizing automata.
Our results link permutation groups and vertex-transitive graphs
to the theory of synchronizing automata.
Furthermore, we explore relations of the size of automata accepting sets of synchronizing
sequences to the length of shortest synchronizing sequences.
\thispagestyle{empty}
\cleardoublepage
\thispagestyle{empty}
\begin{center}
	\bfseries
	Zusammenfassung
\end{center}
%!TEX root = ../thesis-trier.tex

\selectlanguage{ngerman}



Wollten Sie schon einmal wieder von vorne anfangen? Hat sich Ihr Smartphone % mobiltelefon? Handy?
oder Computer schon einmal aufgehängt %aufgehangen  
und nur ein Neustart schuf Abhilfe? Falls ja, so wissen Sie, was es heißt etwas zurückzusetzen, ein zentrales Konzept
der vorliegenden Arbeit.
Im Speziellen handelt sie von der Zurücksetzbarkeit (oder synonym Synchronisierbarkeit) endlicher Zustandsübergangssysteme, auch endliche Automaten genannt.
Automaten interagieren über Eingabefolgen (auch Wörter genannt), welche Zustandsübergänge bewirken, mit ihrer Umwelt
und dienen der Modellierung einer Vielzahl von Systemen wie Schalt\-kreisen, 
Betriebszuständen von Maschinen oder einfachen Agenten in der künstlichen Intelligenz.
%Oft möchte man wissen, ob ein System in einen wohl\-de\-fi\-nierten Zustand versetzbar ist. 
%Beispielsweise könnte ein Schaltkreis nicht mehr synchron mit all seinen Teilen laufen; im Netzwerk könnte der aktuelle Auf\-ent\-halts\-ort unbekannt sein, man möchte aber dennoch die Möglichkeit haben stets einen Weg hinaus zu finden; ein Roboter könnte bei einer Aufgabe festhängen, oder wir kennen einfach nicht seinen aktuellen Zustand, sodass wir ihn in einen wohl\-de\-fi\-nierten Betriebszustand setzen wollen. 
%
%Oft möchte man wissen, ob ein System, beispielsweise einen Schaltkreis dessen Teile
%nicht mehr synchron laufen, in einen wohl\-de\-fi\-nierten Zustand versetzbar ist. 
%Beispielsweise könnte Teile eines Schaltkreises nicht mehr synchron laufen
%oder der aktuelle Zustand einer Maschine ist unbekannt
%und soll zurückgesetzt werden.
%Wir nennen nun einen Automaten synchronisierbar, sofern
Ein Automat heißt synchronisierbar, sofern
es eine Eingabefolge gibt, welche
ihn in einen wohldefinierten Zielzustand versetzt, unabhängig von seinem aktuellen Zustand.
Die Schwierigkeit auf Synchronisierbarkeit zu testen %, sowie die Berechnung einer (möglichst kurzen) synchronisierenden Eingabefolge,
hängt von verschiedenen Parametern ab, zum Beispiel
ob wir das System vollständig kennen, welches Verhalten beobachtbar ist,
ob Nichtdeterminismus erlaubt ist oder ob bestimmte synchronisierende Folgen ausgeschlossen sind.

In dieser Arbeit beschäftigen wir uns mit der Einschränkung 
möglicher synchronisierender Eingabefolgen, das sog.~\emph{constrained synchronization problem} (dt.~\emph{Problem der eingeschränkten Synchronisierbarkeit}).
Es stellt sich nun in natürlicher Weise die Frage, wie die Art der Einschränkung 
die algorithmische Komplexität be\-einflusst. Kann man anhand einer gegebenen Einschränkung
die resultierende Problemstellung in schwierig oder praktikabel lösbar einteilen?
Wir untersuchen komplexitätserhaltende Operationen auf Mengen von erlaubten synchronisierenden Eingabefolgen. Weiterhin
arbeiten wir verschiedene hinreichende Bedingungen  für 
praktikabel lösbare oder schwierige Probleme heraus.
Im Rahmen unserer Untersuchungen tauchen folgende Komplexitätsklassen auf: 
$\AC^0$, \LOGSPACE, \NL, \PTIME, \NP, \PSPACE,
wobei wir Bedingungen für Mitgliedschaft in einer dieser Klassen
oder Schwierigkeit bezüglich \NP\ oder \PSPACE\ 
angeben sowie, außer für~$\AC^0$ und~\PTIME, 
eine Vielzahl von 
vollständigen Problemen für die jeweiligen Komplexitätsklassen.
Die folgenden Arten von Einschränkungen für Mengen synchronisierender Eingabefolgen untersuchen wir: 
endlich, unendlich, spärlich, buchstaben-beschränkt, kommutativ oder durch kleine Automaten beschrieben.
Für kommutative, buchstaben-beschränkte und durch bestimmte kleine Automaton
beschriebene Einschränkungen  
geben wir eine komplette Klassifikation der algorithmischen Komplexität des zugehörigen eingeschränkten Synchronisierbarkeitsproblems. 
Darüber hinaus untersuchen wir die folgenden Klassen von Eingabeautomaten: monotone, partiell geordnete,
kommutative, auflösbare SI- und $\rho$-Automaten.
Unsere Ergebnisse zeigen, dass die resultierenden algorithmischen Komplexitäten
stark vom betrachteten Eingabetyp abhängen. Bei\-spiels\-wei\-se ist das Problem
für kommutative Eingabeautomaten stets in Polynomialzeit  lösbar.

Die Menge an Eingabefolgen, die einen Automaten synchronisieren, das heißt unser Suchraum, 
ist durch einen weiteren Automaton akzeptierbar, was die Untersuchung von akzeptierenden Automaton für jene Mengen nahelegt.
Von besonderem Interesse in Zusammenhang mit unseren Fragestellungen
ist der Fall, wenn ein akzeptierender Automat zur Beschreibung synchronisierender Eingabefolgen
maximal groß ist.
Wie wir hier herausarbeiten, wird die Größe solcher Automaten
durch zwei Faktoren beeinflusst: Erreichbarkeit und Unterscheidbarkeit.
Wir untersuchen beide Begriffe, stellen Zusammenhänge her
und setzen diese in Beziehung zu etablierten Kon\-zepten aus der Literatur, 
wie zum Beispiel zu
vollständig erreichbare Automaten (engl.~completely reachable automata).
Verschiedene Bedingungen für die Maximalität von Automaten für synchronisierende Eingabefolden
werden herausgearbeitet.
Weiterhin formulieren und untersuchen wir Entscheidungsprobleme und algorithmische Fragestellungen, 
welche sich aus diesen Begriffen ergeben.
Über ihren Transformationsmonoid haben Automaten einen natürlichen Bezug
zur Algebra.
Auf dieser Grundlage werden wir Verbindungen zur Theorie der Permutationsgruppen
herstellen. Im Speziellen
arbeiten wir, unter Zuhilfenahme graphentheoretischer Konstruktionen,
 verschiedene neuartige Charakterisierungen für primitive Permutationsgruppen heraus.
Unsere Ergebnisse verbinden Permutationsgruppen und ecken-transitive Graphen (engl.~vertex-transitive graphs) 
mit der Theorie synchronisierbarer Automaten.
Abschließend untersuchen wir, inwiefern die Größe eines akzeptierenden Automaten für die synchronisierenden Eingabefolgen
eines gegebenen Automaten mit der Länge einer kürzesten synchronisierenden Eingabefolge in Beziehung steht.

\selectlanguage{english}
\cleardoublepage
\phantomsection
\begin{center}
	\bfseries
	Acknowledgements
\end{center}
\thispagestyle{empty}
%!TEX root = ../thesis-trier.tex



%\begin{center} \large
%	\textbf{ACKNOWLEDGMENTS}
%\end{center}

%\todo{trennlinie oben wegbekommen.}

%\bigskip 

%Todo

%\begin{comment}
%\begin{center}
%\begin{minipage}{0.7\textwidth}



% me and also for encouragement to develop (and write down) my
%ideas!


First, I thank my supervisor Prof.~Dr.~Henning Fernau for giving me
the freedom and opportunity to pursue my research and the topics that interest me and also for encouragement to develop (and write down) my
ideas!
His open-mindedness, broad knowledge, diplomatic attitude, and willingness to engage into different topics
created just the right (research) environment. I learned a lot from him, both technically and humanly.
%I hope that our paths will still cross in the future.
Furthermore, I express my gratitude to him for having infinite patience with me; in particular with respect
to this thesis which took way more time than intended. In this regard, I also have to thank Prof.~Dr.~Müller-Hannemann
from the Martin-Luther University Halle-Wittenberg for his patience and understanding with my delays; 
even if things did not worked out as planned, for which I am deeply sorry.


%\medskip 

I thank Prof.~Dr.~Mikhail Volkov (Ural Federal University)
for introducing our working group to synchronizing automata at a common workshop in 2019
and suggesting the topics that form the basis
of the present thesis. I also thank him for agreeing to be my co-supervisor, despite this
was not possible in the end due to political reasons beyond our control.
%current world affairs. %political reasons.%never happened...

%\medskip 

I thank Prof.~Igor~Rystsov for agreeing to be my co-supervisor on short notice.
Though I never met him in person, his research and ideas, as is certainly noticeable,
influenced the content of the present work.

%\medskip 

Furthermore, I am grateful to all the friends that made my time in Trier so enjoyable,
in particular Richard Spang, Anna Andres, and Ray Kandar for the parties and moments we shared!
Thanks go to the Gladiators Trier Soccer Team of the Institute of Archeology for letting me 
join their training despite not being an archaeologist  (special thanks to Lukas Märker here!). 
Without the regular physical activity, my intellectual activity would certainly not be possible.

I also thank the university choir for letting me participate, and the concerts
and rehearsals I enjoyed. In particular, I thank our bass leader Steffen Kramer, 
for many fun moments and also, together with Malte Busch, many nice movie nights and subsequent discussions. I am also thankful for all the other moments 
I shared with Malte traveling or spending time at a bar, or the wonderful
people I meet due to him, as, for instance, Jana Thielen.


And I thank my flatmates: Andre Janipour, Florian Höller, Chemseddine Amrani, and Asek Ako-Agbor Ashu
for getting along so well and a relaxed and open
atmosphere, despite me not being very social for the last years due to my graduation efforts.
In particular, I thank Florian Höller for %reminding me how great the 90's were by 
sharing his 
''good mood 90's'' music playlist with me that 
played (and put me in a good mood) during the last weeks of writing up the present thesis. 

Furthermore, I thank Henning Fernau, Richard Spang, Malte Busch, Anna Andres,
and (for a short glance before falling asleep on a plane)
Giuseppe Romana for reading through a preliminary version of the present work
and giving valuable feedback.
 

%\medskip 

Lastly, I thank my parents for raising and supporting me, and making it possible for me to study and graduate.

%\pagenumbering{gobble} \todo{das wird unten seitenzahl angezeigt}

\begin{flushright}
	Stefan Hoffmann, 2025
\end{flushright}

%Without them, this thesis would have never been possible.
%\end{comment}
%\end{minipage} 
%\end{center}
\cleardoublepage
%\end{comment}

\pagestyle{fancy}
\setfancyhf
\tableofcontents
\cleardoublepage

% insert list of figures
\phantomsection
\addcontentsline{toc}{chapter}{\listfigurename}
\listoffigures
\cleardoublepage

% insert list of tables
\phantomsection
\addcontentsline{toc}{chapter}{\listtablename}
\listoftables
\cleardoublepage

\phantomsection
\addcontentsline{toc}{chapter}{Notation Index}
\chapter*{Notation Index} \label{chap:notationindex}






%https://www.resurchify.com/latex_tutorial/latex_glossaries.php
%https://en.wikibooks.org/wiki/LaTeX/Glossary
%https://www.google.com/search?q=latex+overview+mathematical+notation+with+page+numbers&rlz=1C1GCEU_deDE823DE823&ei=Ao-HY_PMBMj6sAermoXoAg&ved=0ahUKEwjz1YWvs9b7AhVIPewKHStNAS0Q4dUDCA8&uact=5&oq=latex+overview+mathematical+notation+with+page+numbers&gs_lcp=Cgxnd3Mtd2l6LXNlcnAQAzIFCAAQogQyBwgAEB4QogQyBQgAEKIEMgUIABCiBDoKCAAQRxDWBBCwAzoFCCEQoAE6CAghEBYQHhAdOgQIIRAVOgcIIRCgARAKSgQIQRgASgQIRhgAUKIGWIgYYKsZaAFwAXgAgAGoAYgB3hCSAQQ1LjEzmAEAoAEByAEIwAEB&sclient=gws-wiz-serp



%https://tex.stackexchange.com/questions/274202/create-a-notation-summary-in-latex
%https://tex.stackexchange.com/questions/274463/feynman-trig-notation-creating-custom-characters?rq=1
%https://tex.stackexchange.com/questions/54720/page-with-symbols-and-notation
%https://www.google.com/search?q=nomenclature&rlz=1C1GCEU_deDE823DE823&oq=nomenclature&aqs=chrome..69i57j0i512l9.280j0j7&sourceid=chrome&ie=UTF-8
%https://www.google.com/search?q=pageref&rlz=1C1GCEU_deDE823DE823&oq=pageref+&aqs=chrome..69i57j0i512l3j0i30l6.2234j0j7&sourceid=chrome&ie=UTF-8
%http://www.weinelt.de/latex/pageref.html

\label{chap:notationindex}

\noindent Standard logical and set-theoretical notation
that is briefly mentioned (and fixed) at the beginnings
of \cref{intro:sec:basic,sec:logic_set_theory_number_theory}
and common %well-known
expressions, as, for example, $\binom{n}{2} = \frac{n(n-1)}{2}$, $u^n$ or $t^n$ for the $n$-fold concatenation product (resp.~composition) of a word $u$ or transformation $t$, etc.,
do not appear in the notation index.
We assume a sufficient level of familiarity
with such notions on the part of the reader to not feel it necessary
to include them in this index. Furthermore, notations that are only mentioned in a single paragraph (for example, in a discussion section)
and are not used beyond that are also not included.

\todo[inline]{long table mit ki prüfen}

% https://tex.stackexchange.com/questions/73283/how-to-use-multirow
% curly braces um den zweizeiler mit modulo notatio?
% https://tex.stackexchange.com/questions/33345/making-a-table-fit-on-a-page-by-moving-it-left
%\begin{longtable}{p{4.6cm}p{8.1cm}p{1.3cm}}
\begin{longtable}{p{0.34\textwidth} p{0.56\textwidth} p{0.1\textwidth}}		
	$f \colon X \to Y$ & (Total) function $f$ from $X$ to $Y$ & \pageref{def:total_function} \\
	$\Sigma$           & A finite alphabet, i.e., set of symbols & \pageref{def:Sigma} \\
	$\Sigma^*$         & Set of words over $\Sigma$ with concatenation & \pageref{def:Sigmastar} \\
	$u\cdot v$ or $uv$ & Concatenation of two words $u,v$ & \pageref{def:Sigmastar} \\
	$|u|$              & Length of word $u$ & \pageref{def:length_of_word} \\
	$\varepsilon$      & The empty word & \pageref{def:empty-word} \\
	$\mathbb N_0$      & $\{0,1,2,3,\ldots\}$ & \pageref{def:N0} \\	
	$\mathbb Z$        & Set of integers & \pageref{def:integers} \\
	$\mathbb R$        & Set of real numbers & \pageref{def:realnumbers} \\
	$[n]$             & The set $\{ 0,1,\ldots, n- 1\}$ & \pageref{def:finite_set_n} \\
	$\max(M) / \min(M)$ & Maximum/minimum of finite $\emptyset \ne M \subseteq \mathbb N_0$ & \pageref{def:min} \\
	$\max(a_1,\ldots,a_n)$ & Maximum of $a_1,\ldots,a_n \in N_0$ & \pageref{def:min} \\
	$\min(a_1,\ldots,a_n)$ & Minimum of $a_1,\ldots,a_n \in N_0$ & \pageref{def:min} \\
	$\semiautomat{A}{Q}{\Sigma}{\delta}$ & Automaton specification of $\mathcal A$ & \pageref{def:automaton} \\
	$\applytrans{\delta}{S}{u}$ & Image of $S \subseteq Q$ in $\semiautomat{A}{Q}{\Sigma}{\delta}$ & \pageref{def:image-delta} \\
	$\Syn(\mathcal A)$ & Set of synchronizing words of  $\mathcal A$ & \pageref{def:set_of_sync_words} \\	
	$\rt{\mathcal A}$  & Reset length of $\mathcal A$ & \pageref{def:rt} \\
	$\mathcal C_n = ([n], \{a,b\}, \delta_n)$   & Automata from the \v{C}ern\'y family & \pageref{def:Cn} \\
	$[p,q]$ & Interval between $p,q \in Q$ in $\mathcal C_n$ & \pageref{def:interval} \\
	$\mathcal P(X)$ & Power set of set~$X$ & \pageref{def:powerset} \\
	$(f(x))_{x \in A}$ & Notation to specify function $f\colon X \to Y$ & \pageref{def:tuple-notation-function} \\
	$\pi_i \colon X_1 \times \cdots \times X_n \to X_i$            & Component projection mapping & \pageref{def:pi} \\
	$a \bmod b$       & Unique $0 \le r < b$ with $a = bn + r$ for $r,n \in \mathbb N_0$ & \pageref{def:bmod} \\
%	$a \equiv_b a'$ or  &  \multirow{ 2}{*}{$a$ and $a'$ are congruent modulo $b$} & \pageref{def:residue_congruence} \\
%	$a \equiv a' \pmod{b}$                  & & \\	
$a \equiv a' \pmod{b}$ or $a \equiv_b a'$                & $a$ and $a'$ are congruent modulo $b$ & \pageref{def:residue_congruence} \\
%	-  & Trivial group &  \pageref{def:trivial group} \\
%	-  & Group of units of a semigroup & \pageref{def:group of units} \\ 
	$\gcd(M), \lcm(M)$    & Greatest com. div. / least common multiple & \pageref{def:gcd} \\
	$\lfloor x \rfloor$ & $\lfloor x \rfloor = \min\{ k\in \mathbb Z \mid k \le x\}$ & \pageref{def:floorfn} \\
	$Y^X$ & Set of functions $f \colon X \to Y$ & \pageref{def:BraisedA} \\ 
	$\applyfunc{S}{f}$ & Extension of $f \colon X \to Y$ to subsets $S \subseteq X$ & \pageref{def:f_set_ext} \\
	$\applyfunc{S}{f^{-1}}, \applyfunc{x}{f^{-1}}$ & Inverse image of $S$ (resp. $x$) for $f \colon X \to Y$ & \pageref{def:preimage} \\
	$\dupl{f}$, $\dupl{u}$ & Duplicate states of $f \colon X \to Y$ / word $u \in \Sigma^*$ & \pageref{def:dupl_t}, \pageref{def:exclword} \\	
	$\excl{f}$, $\excl{u}$ & Excluded states of $f \colon X \to Y$ / word $u \in \Sigma^*$ & \pageref{def:excl_t}, \pageref{def:exclword} \\	
	$\concatcirc{f}{g}$ or $\concat{f}{g}$ & Function composition & \pageref{def:fun_comp} \\	
	$\rank{f}$, $\rank{u}$ & Rank of $f \colon X\to Y$ / word $u \in \Sigma^*$ & \pageref{def:rank}, \pageref{def:rankword} \\
	$\defect{f}$, $\defect{u}$ & Defect of $f \colon X \to Y$ / word $u \in \Sigma^*$ & \pageref{def:defect}, \pageref{def:defectword} \\
	$\excess{f}{A}$ & Excess of~$A$ under~$f$ & \pageref{def:excess} \\
	$(k_1, \ldots, k_l)$ & Kernel-type notation & \pageref{def:kernel-type} \\
	$\fix{t}, \move{t}$ & Set of fixpoints / moved points & \pageref{def:fixmove} \\  
	$\operatorname{id}_X \colon X \to X$  & Identity transformation on $X$ & \pageref{def:identity_transformation} \\
	$(a_0 ~ a_1 ~ \cdots ~ a_{m-1})_X,  ()$ & Cycle notation & \pageref{def:cycle-notation} \\
	$[a_1 \mapsto \cdots \mapsto a_m]_X$ & Tail transformation notation & \pageref{def:SI-mapping} \\
	$[x]_{\sim}$ & Equivalence class of $x$ w.r.t.~$\sim$ & \pageref{def:equivalenceclasses} \\
	$X / {\sim}$ & Set of equivalence classes w.r.t.~$\sim$ & \pageref{def:eqfactor} \\
	$\pi_{\sim} \colon X \to X / {\sim}$ & Canonical projection map & \pageref{def:canprojmap} \\
	$\sim_f$     & Kernel rel.~of $f$ (equiv.~rel.~induced by $f$) & \pageref{def:kernel} \\
	$\upperset{T}$           & Smallest upper set containing $T$ & \pageref{def:genupperset} \\		
	$+_n, \cdot_n$ & Cropped operations on $[n]$ & \pageref{def:plusmaln} \\
	$AB$ & Product on subsets $A,B \subseteq S$, $S$ semigroup & \pageref{def:product_subsets} \\
	$A^n$, $A^{\le n}$, $A^{< n}$, $aA$, $Aa$ & Operations derived from products on subsets & \pageref{def:powers-smaller} \\	
	$\le_S$,$\le_M$, $\le_G$ & Subsemigroup, submonoid, subgroup & \pageref{def:subsemigroup} \\ %,\pageref{def:subgroup} \\   
		$Z(G)$ & Center of a group $G$ & \pageref{def:center} \\	
	$\langle T \rangle_S$, $\langle T \rangle_M$, $\langle T \rangle_G$ & Semigroup, monoid, group generated by $T$ & \pageref{def:gensubstructures} \\
		$A = (a_{ij})_{\substack{i\in \{1,\ldots,n\}, j \in \{1, \ldots, m\}}}$ & $n \times m$ matrix & \pageref{def:matrix} \\	
	%	$v = (v_i)_{i \in \{1,\ldots n\}}$ & $n$-dimensional vector & \pageref{def:vector} \\
	$I_n$, $O_n$ & $n \times n$ identity and zero matrices, respectively.  & \pageref{def:identitymtx} \\	
	    $\mathcal T_X$ & Full transformation semigroup on $X$ & \pageref{def:T_X}  \\
	  $\mathcal T_n$ & Full transformation semigroup on $[n]$ & \pageref{def:T_n} \\
	    $\mathcal S_X$ & Full symmetric group on $X$ & \pageref{def:S_X} \\
	$\Orbit{a}{G}$  & Orbit of $a$ in permutation group $G$      & \pageref{def:orbit} \\
	 $\Stab{a}{G}$   & Point stabilizer of $a$ in permutation $G$ & \pageref{def:stabilizer} \\
	$t^g$           & Conjugation of $t \colon X \to X$ by $g \colon X \to X$ & \pageref{def:conjugation} \\
		$\mathcal S_n$ & Full symmetric group on $[n]$ & \pageref{def:S_n} \\
		$\mathcal A_n$ & Alternating group on $[n]$ & \pageref{def:A_n}  \\
		$L^*, L^+$ & Kleene star and plus & \pageref{def:Kleenestar} \\
		$\pref{u}, \pref{L}$ & Set of prefixes of word $u$ (resp.~$L \subseteq \Sigma^*$) & \pageref{def:pref} \\
	$\suff{u}, \suff{L}$ & Set of suffixes of word $u$ (resp.~$L \subseteq \Sigma^*$) & \pageref{def:suff} \\
	$\fact{u}, \fact{L}$ & Set of factors/infixes of word $u$ (resp.~$L \subseteq \Sigma^*$) & \pageref{def:fact} \\
	$|u|_a$ & Number of occurrences of $a \in \Sigma$ in $u \in \Sigma^*$ & \pageref{def:occurrences_of_letters} \\ 
	$u^{-1}L$ & Language quotient of $L$ by $u$ & \pageref{def:quotient} \\
		$\pi_{\Gamma}\colon \Sigma^* \to \Gamma^*$ & Projection morphism & \pageref{def:projection} \\
		$\mathcal{RE}(\Sigma)$ & Set of regular expressions over $\Sigma$ & \pageref{def:RegSigma} \\
	$L(R)$                   & Language described by $R \in \mathcal{RE}(\Sigma)$ & \pageref{def:language_reg_ex} \\
	$\psi\colon \Sigma^* \to \mathbb N_0^{\Sigma}$ & Parikh map/Parikh morphism & \pageref{def:Parikhmap} \\
	$\operatorname{Aut}(G)$ & Automorphism group of a graph $G = (V,E)$ & \pageref{def:AutG} \\
	$G_t$ / $G_R$ & Functional graph / graph of a relation & \pageref{def:Gf} \\ % \pageref{def:GR} \\
	$N^+(U)$      & Vertices $v$ with edge $(u,v)$ s.t.~$u \in U$, $v \notin U$ & \pageref{def:Nplus} \\
	$O(g)$, $o(g)$  & Big $O$ and small $o$ notation for $g \colon \mathbb N_0 \to \mathbb N_0$ & \pageref{def:bigO} \\
	($\LOGSPACE$-uniform)~$\NC^k$, $\AC^k$                  
	& Circuit families or 
	languages decidable by circuit families in ($\LOGSPACE$-uniform) $\NC^k$/$\AC^k$  &  \pageref{def:AC0circuitclass}, \pageref{def:ACkNCk} \\
	$\LOGSPACE$  & Languages decidable in det. log. space & \pageref{def:L}  \\
	$\NL$        & Languages decidable in nondet. log. space & \pageref{def:NL} 		\\ 
	$\NC$ & Language class $\bigcup_{k\ge 0} \text{\LOGSPACE-uniform } \NC^k$  &  \pageref{def:NC} \\
	$\PTIME$     & Languages decidable in det. poly. time & \pageref{def:P} \\
	$\NP$        & Languages decidable in nondet. poly. time & \pageref{def:NP}  \\
	$\PSPACE$    & Languages decidable in det. poly. space & \pageref{def:PSPACE} \\
	% $\NPSPACE$   & Problems decidable in nondet. poly. space & \pageref{def:NPSPACE}  \\
	$\NPSPACE$   & Languages decidable in nondet. poly. space & \pageref{def:NPSPACE}  \\
	$\lelogm$ / $\eqlogm$   & log-space many-one reduction / equivalent& \pageref{def:lelogm} \\
	$\lePm$ / $\eqPm$      & polynomial time many-one reduction / equivalent & \pageref{def:lePm} \\
	$\leACkm{0}$ / $\leNCkm{1}$ & \LOGSPACE-uniform $\AC^0$ / $\NC^1$ many-one reduction & \pageref{def:leACZEROm} \\ 
	% $\eqlogm$    & log-space many-one equivalent & \pageref{def:eqlogm} \\
	% $\eqPm$      & polynomial time many-one equivalent & \pageref{def:eqPm} \\
	$\eqACkm{0}$ / $\eqNCkm{1}$ & \LOGSPACE-uniform $\AC^0$ / $\NC^1$ many-one equivalent & \pageref{def:eqACZEROm} \\
	$\leNCkT{1}$ / $\eqNCkT{1}$ & \LOGSPACE-uniform $\NC^1$ oracle reduction / equivalent & \pageref{def:eqACZEROm} \\
%	$\eqNCkT{1}$ & \LOGSPACE-uniform $\NC^1$ oracle equivalent & \pageref{def:eqACZEROm} \\
	% $\AC^0$                  & Polynomial size, constant depth \& unbounded fan-in circuit complexity class & \pageref{def:AC0}	
	$G^D$ & Diagonal action on~$G$ & \pageref{def:diagonal-action}\\
	$\applytrans{\delta^{-1}}{q}{u}$ & Inverse transition of~$\delta$ & \pageref{def:inverse_transition} \\
	$G(\mathcal A)$ & Automaton graph of automaton $\mathcal A$ & \pageref{def:aut_graph} \\
		$\lereach$, $\lereach^{\mathcal A}$  & Reachability quasiorder  w.r.t.~$\mathcal A$ & \pageref{def:reach_order} \\
		$\eqreach$               & Equivalence relation induced by $\lereach$ & \pageref{def:induced_reach_order} \\
		$\fword{\delta}{u}$ & $\fword{\delta}{u}\colon Q \to Q$ induced by $u \in \Sigma^*$ in $\semiautomat{A}{Q}{\Sigma}{\delta}$ & \pageref{def:fwordu} \\
	$\TS{\mathcal A}$, $\TM{\mathcal A}$ & Transformation semigroup/monoid of~$\mathcal A$ & \pageref{def:transformation_semigroup} \\
%	 & Transformation monoid of~$\mathcal A$ & \pageref{def:transformation_monoid} \\
			$[a]_{\mathcal A}$, $[u]_{\mathcal A}$, $[a]$, $[u]$  & Transition matrix for letter/word $a$/$u$ w.r.t.~$\mathcal A$ & \pageref{def:transitionmtx} \\	
	$\mathcal A \cong \mathcal B$ & $\mathcal A$ and $\mathcal B$ are isomorphic & \pageref{def:isomorphic} \\
				$\mathcal A_{|S}$,
	$\mathcal A_{|\Gamma}$,
	$\restrictAutGamS{\mathcal A}{\Gamma}{S}$ & Restriction of $\mathcal A$ to upper set $S$, $\Gamma$ or both. & \pageref{def:restriction_aut} \\
		$\semiautomat{A / {\sim}}{Q / {\sim}}{\Sigma}{\delta_{\sim}}$& Quotient automaton w.r.t.~$\TM{\mathcal A}$-congruence $\sim$ & \pageref{def:factor_aut} \\
	$\sim_S$                 & Rees congruence for upper set $S$ & \pageref{def:Rees_congruence} \\
				$\mathcal A_1 \times \cdots \times \mathcal A_m$, $\mathcal A^n$ & Product aut.~of $\mathcal A_1, \ldots, \mathcal A_m$ / $n$-fold product & \pageref{def:product_aut} \\
	$\mathcal P(\mathcal A)$ & Power automaton of $\mathcal A$ & \pageref{def:power_aut} \\
	$\mathcal A^{\le 2}$ & Auxiliary construction for pair automaton & \pageref{def:pair_aut} \\
$\mathcal A^{\le 2} / {\sim_F}$ & Pair automaton of $\mathcal A$ &  \pageref{def:pair_aut} \\	
		$\automat{A}{Q}{\Sigma}{\delta}{q_0}{F}$ & Automaton as language acceptor & \pageref{def:automaton-acceptor} \\
	%$\semiautomat{A_{|S, \Gamma}}{Q}{\Gamma}{\delta_{|S,\Gamma}}$ & Restriction of $\mathcal A$ to upper set $S$ and  $\Gamma\subseteq \Sigma$ & \pageref{def:restriction_aut} \\
	%$\semiautomat{A_{|S}}{Q}{\Sigma}{\delta_{|S}}$ & Abbreviation for $\mathcal A_{|S, \Sigma}$, called subautomaton & \pageref{def:restriction_aut} \\
	$L(\mathcal A)$          & Language accepted by automaton $\mathcal A$ & \pageref{def:accepted_lang_aut} \\	
			$\mathcal A_{q,S}$ & Changing/adding start $s$ \& final set $F$ to $\mathcal A$ & \pageref{def:add_start_final} \\	
	%	$\equiv_{\mathcal A, F}$ & Distinguishability relation w.r.t.~$\mathcal A$ and $F$ & \pageref{def:dist_relation} \\
	$\equiv_L$ & Myhill-Nerode relation & \pageref{def:Nerode} \\
	$\minaut{L}$ & Minimal automaton of $L \subseteq \Sigma^*$ & \pageref{def:minaut} \\
	$\pi_{\mathcal A}$ & Morphism onto $\mathcal A_L$ from accessible $\mathcal A$ & \pageref{def:Nerode-morphism} \\ 
	$\le_{\mathcal R}$ & $\mathcal R$-quasiorder on words w.r.t.~automaton & \pageref{def:R-quasiorder} \\
	$\sim_s$           & Stability relation & \pageref{def:stability} \\
	$\Sigma_I, \Sigma_G$ & Subalphabets w.r.t.~OI-automaton & \pageref{def:subalphabets_SI} \\	
	$\mathcal B_{n,m, r}$ & Strongly conn.~comm. automata & \pageref{def:grid-aut} \\	
	$\Gamma_k(\mathcal A) = (Q_k, E_k)$ & Bondar-Casas-Volkov graph construction & \pageref{def:Gammakgraphs} \\
    $I(\mathcal A) = (Q, E)$ & The Rystsov graph & \pageref{def:Rystsovgraph} \\
    $I_i(\mathcal A) = (Q, E_i)$ & Auxiliary graph introduced by Rystsov & \pageref{def:Rystsovgraph} \\    
	$I(S) = ([n], E)$                   & Rystsov graph extended to $S \le_S \mathcal T_n$ & \pageref{def:graph_IS} \\
	$\Gamma_1(S) = ([n], E)$ & $\Gamma_1(\mathcal A)$-graph extended to $S \le_S \mathcal T_n$ & \pageref{def:GammaS} \\
	$\Gamma_{1,i}(\mathcal A)$ & Variant of $\Gamma_1(\mathcal A)$-graph & \pageref{def:Gamma1i-graph} \\	
	$\mathcal A_{f,a}$ & Adding letter $a$ to $\mathcal A$ defined via $f \colon Q \to Q$ & \pageref{def:Afa} \\
	$a \rightsquigarrow_T b$ / $a \rightsquigarrow b$ & Quasiorder used in~\cref{sync_max:grp:Aprim:sync}. & \pageref{def:quasiorderidemp} \\
	$\mathbf{D},\mathbf{E},\mathbf{F},\mathbf{G},\mathbf{H},\mathbf{I},\mathbf{J},\mathbf{K}$ & Exceptional automata for \v{C}ern\'y bound & \pageref{def:exceptional_aut} \\
	$\mathcal V_n = \mathcal W_n$  & Class of slowly synchronizing automata & \pageref{def:Vn} \\
	$\mathcal L_n = \mathcal D_n'$ & Class of slowly synchronizing automata & \pageref{def:Ln} \\
	$\mathcal F_n$, $n > 3$ odd    & Class of slowly synchronizing automata & \pageref{def:Fn} \\
	$\mathcal K_n = \mathcal G_n$, $n > 5$ odd & Class of slowly synchronizing automata & \pageref{def:Kn} \\
		$\overline{\mathcal A_n}$, $n \ge 8$ even & Class of slowly synchronizing automata & \pageref{def:overlineAn} \\
		$\mathcal M_n$, $n \ge 3$              & Class of slowly synchronizing automata & \pageref{def:Mn}  \\
		$\mathcal B_n$, $n \ge 5$ odd &  Class of slowly synchronizing automata & \pageref{def:Bn} \\
%		$\mathcal{A}_{\operatorname{add}(q, S)}$ & Construction from Ex.~\ref{syncmax:slowly:ex} & \pageref{def:Aadd} \\
		$f \colon X \nrightarrow Y$ & Partial function $f$ from $X$ to $Y$ & \pageref{def:partial_function}  \\
	$\sigma_L \colon \Sigma^* \to \Sigma^*$ & Levenshtein-Romanov mapping & \pageref{def:levenshtein-romanov-mapping} \\
 	$\rho \colon \Sigma^* \to \mathbb N_0$ & Equal-letter runs & \pageref{def:rho} \\
	$\eqLsimU{L}{\sim}{U}$ & Relation derived from $\equiv_L$ & \pageref{def:Nerode-generalization} \\
	$\equiv_{L}^{>n}$  & Relation defined with $\eqLsimU{L}{\sim}{U}$ & \pageref{def:equivLn}  \\
	$\preeqLL{C}{V}$  & Relation used to define $\eqLL{C}{V}$ & \pageref{def:preeqLL}  \\		
	$\eqLL{C}{V}$      & Relation defined with $\eqLsimU{L}{\sim}{U}$  & \pageref{def:eqLL} \\ 
	$\Sigma_{i,j}$ & Subalphabet w.r.t.~small constraint automaton & \pageref{def:Sigmaij} \\
	%$\vec{u}, \vec{v}, \vec{x}, \vec{y}, \ldots$ & Notation for elements in $\mathbb N_0^{\Sigma}$ in Sec.~\ref{sec:constr_sync:commutative} & \pageref{def:vecs} \\
	$\psi_n \colon \Sigma^* \to \{0,1,\ldots,n\}^{\Sigma}$ & $n$-cropped Parikh map & \pageref{def:n-cropped-Parikh-map} \\
	$L_i^F$ & Shorthand $L_i^F = L(\mathcal B_{i,F})$ for Table~\ref{constr_sync:table:three_states_23} & \pageref{def:LiFnotation} \\		
	$\mathcal A_t, \mathcal A_c$ & Sync.~binary SI-automata, see Prop.~\ref{prop:structure_simp_idem_binary} & \pageref{def:AtAc}
\end{longtable}
\cleardoublepage

\pagenumbering{arabic}

%!TEX root = ./springer-mono-format/book.tex

\partwithquote{Introductory Part}{``That fondness for science has encouraged me to compose a short work on calculating by al-jabr (completion) and al-muqabala (balancing), confining it to what is easiest and most useful in arithmetic.''}{---Muhammad ibn Musa Al-Khwarizmi (780--850)}
\label{part:intro}

%!TEX root = ./../springer-mono-format/book.tex
%!TEX root = ./../thesis-trier.tex

\chapter{Short Introduction to Synchronizing Automata}
\label{intro:chap:short_intro_sync_aut}
\label{intro:aut_and_formal_lang}

%!TEX root = ./../../springer-mono-format/book.tex
%!TEX root = ./../../thesis-trier.tex

This chapter gives a short introduction to synchronizing automata for readers unfamiliar with the topic. \Cref{intro:sec:basic} fixes the notation needed here. The notation for the
remainder of the thesis is given in~\cref{chap:prelim}.
In \cref{intro:sec:defsyncaut} 
we formally introduce synchronizing automata. Then, in~\cref{intro:sec:examples}, we present several introductory examples and 
problems and discuss their relation to synchronizing automata.
Finally,
\cref{intro:sec:cerny} defines the reset length, states \Cerny's conjecture,
and presents the \Cerny automata, which serve as a running example throughout
the thesis.

For readers who want to learn more about synchronizing automata, the most recent survey articles are by Kari and Volkov~\cite{KariVolkov2021} (see also its precursor~\cite{DBLP:conf/lata/Volkov08}) and by Volkov~\cite{Volkov22}, 
the latter containing a good coverage
of applications from different fields and types of automata considered in the literature.
 The surveys by Jürgensen~\cite{DBLP:journals/iandc/Jurgensen08}, Kari and Volkov~\cite{KariVolkov2021}, and Béal and Perrin~\cite{BealPerrin2016} discuss connections to codes (see also Chapters~3 and~10 of~\cite{DBLP:books/daglib/0025093}). Sandberg~\cite{Sandberg04} and Kohavi~\cite[Chap.~13]{Kohavi70} include relations to automata with outputs as well as applications in circuit theory and testing. The survey by Bogdanović et al.~\cite{Bogdanovic99} presents structural results not covered in the other surveys. Another survey is by Mateescu and Salomaa~\cite{MateescuS01}. Finally, we recommend~\cite{ArCamSt2017} on synchronizing permutation groups, motivated by synchronizing automata, as results presented in that survey were inspirational for our results in \cref{sync_max:sec:perm_grps} from \cref{part:aut}.

\section{Notions Used in the Introduction}
\label{intro:sec:basic}

\input{../part1/sync/notions}	

\section{Formal Definition of Automata Synchronization}
\label{intro:sec:defsyncaut}

\input{../part1/sync/def}	

\section{Introductory Examples}
\label{intro:sec:examples}

\input{../part1/sync/examples}	

\section{Reset Length and the \Cerny Conjecture} 
\label{intro:sec:cerny}

\input{../part1/sync/resetlength}

\chapter{Outline and Contribution}
\label{intro:chap:overview}

First, in \cref{overview:sec:problems} we introduce the two central 
problems with which the thesis is concerned: the constrained
synchronization problem (\cref{overview:sec:csp})
and the description of automata for sets of synchronizing 
words (\cref{intro:overview:structure_set_of_sync_words}).
\Cref{overview:sec:organization} gives an overview on the 
organization of the present work.
Finally, \cref{overview:sec:papers} lists the papers on which this work is based, 
as well as the unpublished results that we have included.

\section{The Problems}
\label{overview:sec:problems}

\subsection{The Constrained Synchronization Problem}
\label{overview:sec:csp}

\IN{constrained synchronization problem} 
\IN{constraint language}
\index{regular constraint language|see{constraint language}}
The \emph{constrained synchronization problem} asks whether a given
deterministic and complete
semiautomaton has  a synchronizing word contained
in a given regular \emph{constraint language}. 
Formally, let~$\Sigma$ be an alphabet 
and $L \subseteq \Sigma^*$ be a regular language 
(see \cref{prelim:sec:decision_problems,prelim:sec:regularlang}
for formal definitions of decision problems and regular languages, respectively).
The problem \constrsync{L} of \emph{constrained synchronization} 
with respect to the \emph{constraint language $L$} is the following decision problem.

\index{\constrsync{L}|see{constrained synchronization problem}}
\begin{center}
\begin{tcolorbox}[width = 0.9\textwidth,colback=white,arc=0pt,outer arc=0pt,boxrule=0.7pt,left =0.5em,right=0em]
\decisionproblem{\constrsync{L}}
{A deterministic and complete semiautomaton $\semiautomat{A}{Q}{\Sigma}{\delta}$.}
{Does there exist a word $w \in L$ synchronizing $\mathcal A$?} 
\end{tcolorbox}
\end{center}

This problem, systematically studied in~\cite{FernauGHHVW19},
firstly appeared in connection with a reduction procedure introduced in~\cite{Gusev2012},
as well as implicitly in reductions used in~\cite{Martyugin14}.
\Cref{constr_sync:sec:motivation} provides further motivation and details.

\subsection{Automata for Sets of Synchronizing Words}
\label{intro:overview:structure_set_of_sync_words}

The set of synchronizing words of a given automaton can be accepted by another automaton of size exponential in the size of the original automaton (see \cref{chap:prelim} for the definition of automata as language acceptors). In this thesis, we also study automata for sets of synchronizing words.
Our motivation comes from the constrained synchronization problem, whose computational complexity is governed by two parameters: the \emph{structure of the constraint language} and the \emph{structure of the set of synchronizing words} of the given input automata.
If an automaton recognizing the set of synchronizing words can be constructed in polynomial time for a class of input automata, then the constrained synchronization problem (with any constraint language) is solvable in polynomial time for this class. An example is the class of commutative automata over a fixed alphabet (see \cref{constr_input:subsec:comm}).
Consequently, \NP- or \PSPACE-hardness results (see \cref{prelim:sec:compl_classes} for definitions of complexity classes) imply that automata for the corresponding sets of synchronizing words cannot, in general, be constructed in polynomial time. 
This motivates our interest in automata 
needing exponentially larger accepting automata for their sets of synchronizing words.

Beyond its connection to the constrained synchronization problem, the study of automata for sets of synchronizing words is also motivated by the following observations.

\begin{enumerate}
\item For known families of \emph{slowly synchronizing automata}, that is, automata whose reset length is close to the \Cerny bound, the minimal automata recognizing their sets of synchronizing words are exponentially larger~\cite{DBLP:journals/corr/Maslennikova14,DBLP:journals/ijfcs/Maslennikova19}. 
This raises the question of whether there is a deeper connection between the reset length of an automaton and the size of a minimal automaton recognizing its set of synchronizing words.

\item The operation that maps a regular language to the set of synchronizing words of its minimal automaton is regularity-preserving and can therefore be studied from the perspective of the state complexity of operations~\cite{GaoMRY17,YuZhuangSalomaa1994}.

\IN{reset complexity} 
\item The class of sets of synchronizing words of automata coincides with the class of regular ideal languages (\cref{prelim:thm:sync_words_is_ideal}). Thus, every regular ideal language can be represented by an automaton having it as its set of synchronizing words. This representation can be exponentially more succinct than a representation by an accepting automaton, 
as witnessed by known families of slowly synchronizing automata~\cite{DBLP:journals/corr/Maslennikova14,DBLP:journals/ijfcs/Maslennikova19}. From the perspective of the descriptional complexity of regular languages, it is therefore natural to investigate those regular ideal languages for which the gap between these two representations is maximal. In automata-theoretic terms, this leads to the study of automata for which minimal automata accepting their sets of synchronizing words are as large as possible.

Furthermore, these observations motivate the definition of the \emph{reset complexity} of a regular ideal language as the size of a smallest automaton having this language as its set of synchronizing words. The \Cerny conjecture is equivalent to the statement that the reset complexity of a regular ideal language is at least the square root of the length of a shortest word in the language, plus one~\cite{DBLP:journals/corr/Maslennikova14,DBLP:journals/ijfcs/Maslennikova19}. 
This yields a language-theoretic approach to the \Cerny conjecture through the investigation of shortest words in regular ideal languages~\cite{DBLP:journals/jalc/FrigeriR25,
	DBLP:journals/jalc/FrigeriR19,
	GusevMP13,
	DBLP:journals/fuin/MaslennikovaR18,
	ReisR16,
	DBLP:journals/aam/Rodaro18}.

	
\item As we will see, our results reveal connections between subset reachability properties,
the size of automata for sets of synchronizing words, and the well-known class of primitive permutation groups.
\end{enumerate}

\section{Organization of the Thesis}
\label{overview:sec:organization}

The thesis consists of four parts.
In \cref{part:intro} we give a short introduction on synchronizing automata
and an overview of the thesis.
\Cref{part:prelim} introduces further notions, automata classes,
and gives preliminary computational complexity results.
A detailed overview is given in \cref{intro:outline:orga:sec:prelim}.
\Cref{part:aut} is concerned with automata accepting the set of synchronizing words of a given automaton and is outlined in \cref{intro:outline:orga:sec:autsync} below.
\Cref{part:constrsync} deals with the constrained synchronization problem
and we give an overview in \cref{intro:outline:orga:sec:csp}.

Parts~\ref{part:aut} and~\ref{part:constrsync} are the main parts.
Unless stated otherwise, all the results stated in these 
parts are due to the author of this thesis.

Historically, and as outlined in \cref{overview:sec:problems}, 
our research on the constrained synchronization problem
led us to investigate automata for sets of synchronizing words.
However, in the present work, we first discuss automata for sets of synchronizing words in \cref{part:aut} before turning to the constrained synchronization problem in \cref{part:constrsync}.
We chose this order of presentation because \cref{part:constrsync} relies on results from \cref{part:aut}, whereas \cref{part:aut} is independent of \cref{part:constrsync}.
Moreover, the study of automata accepting sets of synchronizing words is a natural topic in its own right, 
and we expand it beyond its application to the constrained synchronization problem.

\subsection{Structure of \Cref{part:prelim}: Preliminaries}
\label{intro:outline:orga:sec:prelim}

Here, we introduce further notions in \cref{chap:prelim}. 
In \cref{chap:class-transformations} we investigate classes of transformations.
Permutation groups are investigated in \cref{prelim:sec:permgrp}.
More automata-theoretical notions are given in \cref{prelim:sec:automata}. 
Classes of automata and their properties
are investigated in \cref{chap:classes_aut}.
Computational complexity results on subset mapping problems
in automata with respect to certain subclasses 
are investigated in \cref{chap:prelim_complexity_results}.

Most results in this part are taken from the literature.
Some results on idempotent and one-point singular
transformations from \cref{chap:class-transformations}
are new. The notion of strong $\frac{3}{2}$-connectivity
for digraphs introduced here is new. Further, some results in \cref{chap:classes_aut}
are new, see also \cref{overview:sec:unpublished} for more details.
Some structural results on solvable automata in \cref{prelim:subclasses:solvable}
appear in~\cite{DBLP:journals/jalc/000124}.
\Cref{prelim:sec:automata} on automata theory contains some results
on morphism and quotient automata that we have not seen elsewhere.
The results in \cref{chap:prelim_complexity_results} 
were first presented in~\cite{DBLP:journals/jalc/000124}.

\subsection{Structure of \Cref{part:aut}: Automata for Sets of Synchronizing Words} 
\label{intro:outline:orga:sec:autsync}


\Cref{part:aut} studies automata that accept the set of synchronizing words of a given automaton. 
We derive new notions and conditions related to the size of 
accepting automata for sets of synchronizing words, also giving conditions 
under which minimal automata for such sets have maximal size. 
We further relate our findings to classical automata-theoretic results 
from the field of descriptional complexity.
This part is based on the papers~\cite{DBLP:conf/lata/Hoffmann21,
	DBLP:journals/tcs/Hoffmann21,
	DBLP:conf/lata/Hoffmann21a,
	DBLP:conf/dcfs/Hoffmann21,
	DBLP:conf/dcfs/Hoffmann22,
	HoffmannCircularExtended,
	HoffmannLATAextended2,
	Hoffmann2024}, 
but the presentation has been extended and unified, with additional new results included.

Except in \cref{sync_max:sec:compl_reach,sync_max:sec:BVC-graphs,sync_max:sec:augmentative,sync_max:sec:rystsov_graph} and the discussion sections, which review notions and results from the literature, all results presented are novel and due to the author. Next, we describe the organization of \cref{part:aut} and give a broad overview of each chapter and its relation to the work on which this thesis is based. A more detailed overview of the chapters, including descriptions of their individual sections, is given at the beginning of each chapter.

In \cref{sync_max:chap:prelim}, we first note that the set of synchronizing words is a regular language acceptable by the power automaton.
To obtain lower bounds on the size of automata accepting the set of synchronizing words, we apply minimality conditions to the power automaton. This leads us to the notions of reachability of subsets and distinguishability of pairs of subsets.
The rest of the chapter investigates these two notions, giving sufficient conditions for reachability and distinguishability. We also introduce completely reachable automata, a class previously studied in the literature, and the new class of completely distinguishable automata, for which we prove several equivalent characterizations.
The results in this chapter are based on~\cite{DBLP:conf/dlt/Hoffmann23,Hoffmann2024}, where completely distinguishable automata were introduced, but it also contains new unpublished material, see \cref{overview:sec:unpublished}.

\Cref{sync_max:chap:sync_max} introduces sync-maximal automata as synchronizing automata for which an automaton accepting the set of synchronizing words is as large as possible, i.e., for which the gap between the state complexity and the reset complexity of the set of synchronizing words is maximal. We investigate sufficient conditions and derive structural properties, showing that sync-maximal automata are closely related to completely reachable automata. We also give various examples. The notion of sync-maximal automata, as well as the structural results, were first introduced in~\cite{DBLP:conf/dlt/Hoffmann23,Hoffmann2024}. 

\Cref{sync_max:chap:graphs} discusses graph-theoretical constructions that yield sufficient conditions for complete reachability and complete distinguishability. First, we review constructions and results from the literature. We then extend some of these constructions to arbitrary transformation semigroups. Finally, we give graph-theoretical conditions ensuring complete reachability or distinguishability of certain pairs of subsets. The results from this chapter are new and have not been published before.

\Cref{sync_max:sec:subclasses} investigates the following subclasses of automata: commutative automata, automata in which every letter has defect at most one, SI-automata, and circular automata.
For commutative automata, we show that an automaton recognizing the set of synchronizing words can be computed in polynomial time when the alphabet is fixed. This result was published in~\cite{DBLP:journals/tcs/Hoffmann21} in the context of the constrained synchronization problem for commutative constraint languages, and will be used later in \cref{part:constrsync}.
We then consider automata in which every letter has defect at most one. In this setting, we show that if the group of units of the transformation monoid is transitive, then synchronizability implies complete distinguishability. As a consequence, a binary automaton with at least three states is completely reachable if and only if it is sync-maximal. These results are taken from~\cite{Hoffmann2024}.
Then we investigate SI-automata. For example,
we show that in synchronizing SI-automata
distinct pairs $2$-sets are distinguishable.
We derive this from more general and unpublished results,
but it also follows from~\cite[Lem.~11]{DBLP:conf/dcfs/Hoffmann22}.
However, the paper~\cite{DBLP:conf/dcfs/Hoffmann22} 
contains a mistake stemming from a false statement~\cite[Prop.~2]{DBLP:conf/dcfs/Hoffmann22}.
We discuss this in detail at the end of \cref{sync_max:sec:subclasses:discussion}.
Lastly, we investigate circular automata and state a sufficient condition for complete reachability and complete distinguishability. A special case of this condition was published in~\cite{HoffmannCircularExtended}, with a preliminary version presented in~\cite{DBLP:conf/lata/Hoffmann21a}.

\Cref{sync_max:chap:decision} studies the decision problems for complete reachability, complete distinguishability, and sync-maximality, as well as related questions.
For example, we prove that these problems 
are \NL-hard for SI-automata, with complete reachability even being \NL-complete for SI-automata,
and in \PSPACE\ for general automata.
The results on SI-automata were published in~\cite{DBLP:conf/dlt/Hoffmann23,Hoffmann2024}.

\Cref{sync_max:sec:perm_grps} 
discusses relations to permutation groups.
We introduce primitive groups and
characterize them in terms of complete reachability,
complete distinguishability, and in terms of sync-maximality
of associated automata.
The relation of primitive groups to complete reachability
was first published in~\cite{DBLP:conf/lata/Hoffmann21},
and the relation to sync-maximality in~\cite{DBLP:conf/dcfs/Hoffmann21}.
The extended journal publication~\cite{HoffmannLATAextended2} 
of~\cite{DBLP:conf/lata/Hoffmann21,DBLP:conf/dcfs/Hoffmann21} contains 
both of these characterizations.
We then introduce $A$-primitivity 
and study its connection to congruences, orbitals,
generating sets, synchronization, reachability,
and homogeneity properties.
Our results generalize existing characterizations
of primitivity from the literature,
and a characterization by
a homogeneity property also gives
a new characterization. 
The notion of $A$-primitivity is introduced here for the first time, and the results of this chapter have not been published previously.

\Cref{sync_max:chap:slowly} discusses families of slowly synchronizing
automata. We observe that, for most of these families of automata,
the minimal automata accepting their sets of synchronizing words
have exponential size. In particular,
we use results from Chapter~\ref{sync_max:sec:subclasses}
on circular automata to deduce sync-maximality
for several families of slowly synchronizing circular automata
from the literature, including one family for which 
sync-maximality has previously only been conjectured.
This confirmation of the conjecture was 
presented in~\cite{DBLP:conf/lata/Hoffmann21a,HoffmannCircularExtended}.
We then show, for three families of slowly synchronizing but not sync-maximal
automata, that minimal automata for their sets 
of synchronizing words are exponentially larger.
Based on our observations, we formulate various conjectures.
Apart from the confirmation of the aforementioned conjecture,
the results from this chapter have not been published.

\subsection{Structure of \Cref{part:constrsync}: The Constrained Synchronization Problem}
\label{intro:outline:orga:sec:csp}

\Cref{part:constrsync} deals with the constrained synchronization problem and is based on the works~\cite{FernauGHHVW19,
	DBLP:conf/cocoon/Hoffmann20,
	DBLP:conf/ictcs/Hoffman20,
	DBLP:conf/fct/Hoffmann21,
	DBLP:journals/tcs/Hoffmann21,
	DBLP:conf/cocoon/Hoffmann21,
	DBLP:conf/wia/Hoffmann22,
	HoffmannCocoonExtended}.
The results from~\cite{DBLP:conf/cocoon/Hoffmann20,DBLP:conf/fct/Hoffmann21,DBLP:journals/tcs/Hoffmann21,DBLP:conf/cocoon/Hoffmann21,HoffmannCocoonExtended}
were revisited with an emphasis on a more streamlined, coherent, and simplified presentation.
This was achieved primarily through the use of general methods introduced in~\cite{DBLP:conf/cocoon/Hoffmann21,HoffmannCocoonExtended}.

In \cref{constr_sync:sec:motivation}, we give motivational examples for the constrained synchronization problem.

\Cref{constr_sync:chap:further_notions} introduces additional notions that we will need in \cref{part:constrsync}. We also present a new characterization of solid codes that, to the best of our knowledge, has not been published before.

\Cref{constr_sync:chap:gen_results} presents general results on manipulating constraint languages while preserving
the complexity or to deduce complexity lower bounds. Furthermore, we give sufficient conditions for \NP- and \PSPACE-hardness. 
Some of these results
were presented in~\cite{DBLP:conf/cocoon/Hoffmann21,HoffmannCocoonExtended},
which generalized results from~\cite{FernauGHHVW19},
and some results are new.

The central question investigated in \cref{constr_sync:chap:classes_constraint_languages}
is the computational complexity of \constrsync{L} for different classes of regular
constraint languages $L \subseteq \Sigma^*$.
While a complete classification remains an \emph{open problem},
we present several partial results by giving complete classifications for the following classes of constraint languages:
letter-bounded languages, languages accepted by partial automata (PDFAs, introduced in
\cref{constr_sync:chap:further_notions}) with at most two states over an arbitrary alphabet and at most three states over a binary alphabet, and commutative languages.
We also obtain results for infinite, sparse, and monoid languages, as well as
for languages accepted by strongly connected PDFAs.
The classification for commutative constraint languages appeared
in~\cite{DBLP:conf/cocoon/Hoffmann20,DBLP:journals/tcs/Hoffmann21}.
Here, we present a simplified formulation and proof.
The results for sparse and letter-bounded languages appeared
in~\cite{DBLP:conf/ictcs/Hoffman20,DBLP:conf/fct/Hoffmann21}.

Lastly, motivated by the question of what features are needed in input automata to realize certain computational complexities, in \cref{constr_sync:chap:constr_input} we investigate specific types of input automata. Some of the results from this chapter appeared in~\cite{DBLP:conf/wia/Hoffmann22}, and some mistakes from this paper have been corrected.

\section{Work by the Author Relevant to the Thesis}
\label{overview:sec:papers} 

\subsection{Published Work}

The present thesis contains and extends material from the following papers and conference contributions published since 2019, in which the author was involved.

\begin{enumerate}
	\item[\cite{FernauGHHVW19}] \fullcite{FernauGHHVW19}
	\item[\cite{DBLP:conf/ictcs/Hoffman20}] \fullcite{DBLP:conf/ictcs/Hoffman20}
	\item[\cite{DBLP:conf/cocoon/Hoffmann20}] \fullcite{DBLP:conf/cocoon/Hoffmann20}
	\item[\cite{DBLP:conf/lata/Hoffmann21a}] \fullcite{DBLP:conf/lata/Hoffmann21a}
	\item[\cite{DBLP:conf/lata/Hoffmann21}] \fullcite{DBLP:conf/lata/Hoffmann21}
	\item[\cite{DBLP:conf/cocoon/Hoffmann21}] \fullcite{DBLP:conf/cocoon/Hoffmann21}
	\item[\cite{DBLP:conf/dlt/Hoffmann21a}] \fullcite{DBLP:conf/dlt/Hoffmann21a}
	\item[\cite{DBLP:journals/tcs/Hoffmann21}] \fullcite{DBLP:journals/tcs/Hoffmann21} 
	(extended version of~\cite{DBLP:conf/cocoon/Hoffmann20})
	\item[\cite{DBLP:conf/fct/Hoffmann21}] \fullcite{DBLP:conf/fct/Hoffmann21}
	\item[\cite{DBLP:conf/dcfs/Hoffmann21}] \fullcite{DBLP:conf/dcfs/Hoffmann21}
	\item[\cite{DBLP:conf/dcfs/Hoffmann22}] \fullcite{DBLP:conf/dcfs/Hoffmann22}
	\item[\cite{DBLP:conf/wia/Hoffmann22}] \fullcite{DBLP:conf/wia/Hoffmann22}
	%\item[\cite{DBLP:conf/wia/Hoffmann22extended}] \fullcite{DBLP:conf/wia/Hoffmann22extended}
	%(extended version of~\cite{DBLP:conf/wia/Hoffmann22})
	\item[\cite{HoffmannCocoonExtended}] \fullcite{HoffmannCocoonExtended} 
	(extended version of~\cite{DBLP:conf/cocoon/Hoffmann21})
	
	%{\sloppy 
		
	\item[\cite{HoffmannLATAextended2}] \fullcite{HoffmannLATAextended2}
	(extended version of~\cite{DBLP:conf/lata/Hoffmann21})
		
		%\fussy}
	\item[\cite{HoffmannCircularExtended}] \fullcite{HoffmannCircularExtended}
	(extended version of~\cite{DBLP:conf/lata/Hoffmann21a})
	
	\item[\cite{DBLP:conf/dlt/Hoffmann23}] \fullcite{DBLP:conf/dlt/Hoffmann23}
	
	\item[\cite{Hoffmann2024}] \fullcite{Hoffmann2024}
	(extended version of~\cite{DBLP:conf/dlt/Hoffmann23})
	
	\item[\cite{DBLP:journals/jalc/000124}] \fullcite{DBLP:journals/jalc/000124}
\end{enumerate}

\subsection{Unpublished Work} 
\label{overview:sec:unpublished}

Here, we discuss results contained in the present thesis that were not published before.

In~\cref{part:prelim}, the results on one-point singulars 
from~\cref{prelim:sec:oneattractor,prelim:sec:idempotentoneattractor}
are new. We were also unable to find the statements 
from~\cref{prelim:sec:idempotent} on compositions
of idempotent transformations in the literature (but most of them are rather technical or perhaps folklore). We also note that our treatment of automata theory in \cref{prelim:sec:automata}, tailored to our needs, 
differs from those found in the literature of which we are aware,
in particular in its usage of the reachability order
and morphisms. Furthermore, many results related to morphisms
and isomorphisms between automata that
we give are new or scattered throughout the literature under different terminology and notation.

In~\cref{part:aut}, the results on reachable subsets from \cref{sync_max:sec:reachable},
on pairs of distinguishable subsets from~\cref{sync_max:sec:pairs_k_k-1,sync_max:sec:shortestdist},
the graph-theoretical results from~\cref{sec:extended_Rystsov_graph}, 
on automata where every letter has defect at most one from~\cref{sync_max:sec:defectone},
and on subset mapping properties for pairs
of equal-sized sets in SI-automata from~\cref{sec:SI:dist} are new (but one preliminary 
result appeared in~\cite[Lem.~11]{DBLP:conf/dcfs/Hoffmann22}).
Furthermore, the notion of $A$-primitivity is novel
and all results from~\cref{sec:A-primitive} 
appear in this thesis for the first time.

In~\cref{part:constrsync}, some characterizations
of solid codes from \cref{constrsync:sec:codes:solid} are (to the best of our knowledge) new.
The complexity lower bounds involving concatenations from \cref{constr_sync:subsec:concat} have not been published previously.
In \cref{constr_sync:chap:constr_input} some published results have
been corrected or extended. 
For example, the classification for partially ordered automata in \cref{constrinput:sec:two-state-pto} refines the results of~\cite{DBLP:conf/dlt/Hoffmann21a}; the bound for monotonic automata claimed in \cref{constr_sync:prop:monotonic_bound} corrects that of~\cite{DBLP:conf/wia/Hoffmann22}; and the claims on solvable automata in \cref{constr_input:subsec:solvable} have likewise been corrected. See \cref{constr_input:sec:discussion} for details.
Also, some material has been presented differently. Further details on new, extended, or revised material are given in the discussion section at the end of each chapter.


\partwithquote{Preliminaries}{``Es steht alles schon bei Dedekind.''}{---Emmy Noether (1882--1935)}
\label{part:prelim}

%!TEX root = ./../springer-mono-format/book.tex
%!TEX root = ./../thesis-format/main.tex

\chapter{General Notions}
\label{chap:prelim}

%!TEX root = ./../../thesis-format/main.tex

Recall the definitions from \cref{intro:sec:basic} regarding set theory, functions, automata, and formal languages. 
Here, we introduce additional notions, facts, and results not covered previously.
\Cref{sec:logic_set_theory_number_theory} gives more basic notions
from logic, set and number theory.
Further notions related to functions, in particular for functions on finite sets,
are introduced in \cref{prelim:sec:fun}.
\Cref{pelim:sec:trans} introduces transformations and gives some results
for transformations on finite sets.
In \cref{prelim:sec:relations} we introduce relations.
Semigroups, monoids, groups, semirings, matrices, and transformation semigroups
are introduced in \cref{prelim:sec:algebra}. Order-theoretical notions are introduced in \cref{prelim:sec:order}.
\Cref{sec:formal_language_theory} gives more background on formal language theory
and introduces regular expressions.
Graphs and connectedness notions are introduced in \cref{prelim:sec:graphs}. 
\Cref{prelim:sec:complexity} introduces computational complexity theory
and our theoretical framework to measure the resource requirements of algorithms.
\Cref{prelim:sec:remarks-reduction,prelim:sec:remarks_func_comp} 
discusses conventions and notational choices.
Lastly, \cref{prelim:prelim:discussion} 
gives an overview of the literature related to the introduced notions.

The present chapter serves to fix notation. We assume some familiarity and prior exposure
with general concepts from logic, set theory, number theory, basic algebra, formal language theory, and complexity theory.

\section{More on Logic, Set Theory, and Number Theory}
\label{sec:logic_set_theory_number_theory}

\input{../part2/notions/logic}

\section{More on Functions} 
\label{prelim:sec:fun}

\input{../part2/notions/functions}

\section{Transformations}
\label{pelim:sec:trans}

\input{../part2/notions/transform}

\section{Relations}
\label{prelim:sec:relations}

\input{../part2/notions/relations}

\section{Algebraic Notions}
\label{prelim:sec:algebra}

\input{../part2/notions/algebra}

\section{Order Theory}
\label{prelim:sec:order}

\input{../part2/notions/order}

\section{Formal Language Theory}
\label{sec:formal_language_theory}

\input{../part2/notions/formallang}

\section{Graph Theory}
\label{prelim:sec:graphs}

\input{../part2/notions/graphs}

\section{Computational Complexity Theory}
\label{prelim:sec:complexity}

\input{../part2/notions/complexity}

\section{Remarks on Different Notions of Reduction} 
\label{prelim:sec:remarks-reduction}

\input{../part2/notions/remark_reduction}

\section{Remarks on the Order of Function Composition}
\label{prelim:sec:remarks_func_comp} 

\input{../part2/notions/order_composition}

\section{Discussion} 
\label{prelim:prelim:discussion}

\input{../part2/notions/discussion}

\chapter{Classes of Transformations}
\label{chap:class-transformations}

%!TEX root = ./../../thesis-format/main.tex

Here, we investigate classes of transformations.
Sections~\ref{prelim:sec:perm}, \ref{prelim:sec:cycles}, \ref{subsec:tail}, \ref{prelim:sec:oneattractor}, and~\ref{prelim:sec:idempotent}
introduce and investigate permutations, cycles, one-point singulars, and idempotent transformations, respectively.
In \cref{prelim:sec:idempotentoneattractor},
we investigate idempotent one-singulars.
We will present different characterizations.
Later, the transformations that arise  in the study of automata theory
result from compositions of transformations induced by letters.
For this reason, we also investigate
resulting transformations when composing different types of transformations.

\section{Permutations and Quasipermutations} 
\label{prelim:sec:perm} 

\input{../part2/transform/perm}

\section{Cycles and Cycles Induced by a Transformation} 
\label{prelim:sec:cycles}

\input{../part2/transform/cycles}

\section{Tail Transformations, Simple Idempotents, and Cycles}
\label{subsec:tail}

\input{../part2/transform/tail}

\section{One-Point Singular Transformations}
\label{prelim:sec:oneattractor} 

\input{../part2/transform/oneattractor}

\section{Idempotent Transformations}
\label{prelim:sec:idempotent} 

\input{../part2/transform/idempotent}

\section{Idempotent One-Point Singulars}
\label{prelim:sec:idempotentoneattractor} 

\input{../part2/transform/idempotentoneattractor}

\section{Discussion}

\input{../part2/transform/discussion}

\chapter{Permutation Groups}
\label{prelim:sec:permgrp}

%!TEX root = ./../../springer-mono-format/book.tex
%!TEX root = ./../../thesis-trier.tex


\begin{flushright}
	\itshape %\setstretch{1.0}
	``It was at this point that Wigner, Hund, Heitler, and Weyl entered the picture with their `Gruppenpest': the pest of the group theory… The authors of the `Gruppenpest' wrote papers which were incomprehensible to those like me who had not studied group theory, in which they applied these theoretical results to the study of the many electron problem. The practical consequences appeared to be negligible, but everyone felt that to be in the mainstream one had to learn about it. Yet there were no good texts from which one could learn group theory. It was a frustrating experience, worthy of the name of a pest.''  \\ \normalfont
	\vspace{.2em}
	\rightskip=.8cm---John Slater (1900--1976)
\end{flushright}

The set $\mathcal S_X$ of all permutations on a set $X$
with function composition
and inversion given by $t^{-1} \colon X \to X$ 
for $t \in \mathcal S_X$ is a group, called the \emph{(full) symmetric group on $X$}.
%
%\bigskip 
%
%For a set $X$, the \emph{(full) symmetric group on $X$} is the group structure on the set~$\mathcal S_X$ of all permutations on~$X$
%with function composition as operation
%and inversion given by $t^{-1} \colon X \to X$ 
%for $t \in \mathcal S_X$.
%
%\bigskip
%
%For a set $X$, the set $\mathcal S_X$ of all permutations on $X$,
%which is a group with function composition
%and inversion given by $t^{-1} \colon X \to X$ 
%for $t \in \mathcal S_X$, called the \emph{(full) symmetric group on $X$}.
A \emph{permutation group~$G$} (over the \emph{(permutation) domain}~$X$) is a subgroup of~$\mathcal S_X$.
Here, we investigate permutation groups.

\Cref{prelim:sec:basic} introduces basic notions. In \cref{prelim:permgrp:sec:transitive}
we collect statements related to transitive permutation groups.
The relation between permutation groups and (abstract) groups
is discussed in \cref{prelim:sec:perm_grp_abstract_grp}. \Cref{prelim:sec:grp-of-units} investigates
the group of units in transformation semigroups.
Lastly, \cref{prelim:sec:graphs-groups} points out some relations between permutation groups and graphs.

\section{Basic Notions}
\label{prelim:sec:basic}

\input{../part2/transform/permgrps}

\section{Permutation Groups and Abstract Groups}
\label{prelim:sec:perm_grp_abstract_grp}

\input{../part2/transform/permvsabstract}

\section{Permutation Groups and Non-Permutations}
\label{prelim:sec:grp-of-units}

\input{../part2/transform/grpofunits}

\section{Graphs and Permutation Groups}
\label{prelim:sec:graphs-groups}

\input{../part2/transform/permgrpgraphs}

\section{Discussion}
\label{prelim:permgrp:sec:discussion}

\todo[inline]{diskussion in datei und oben kommentierten teil raus}
Jordan~\cite{jordan1870} wrote the first textbook on permutation group theory (but with an antique notation and terminology).
Prior to this, permutation groups in the context of polynomial
equations (as permutations of their root)
were investigated by Lagrange, Galois, Abel and Ruffini, see~\cite{Kleiner2007_HistoryAbstractAlgebra,Neumann2011,Wussing1984_GenesisAbstractGroup}.
At the end of the 19th century, Burnside published his classical textbook~\cite{burnside1911} (1st edition 1897) 
which collected and contained most results from 19th century group theory with a focus on permutation groups.
Wielandt~\cite{wielandt1964finite} wrote
a standard textbook on permutation group theory
which established some standard terminology.
Other standard introductions are the textbooks by Dixon \& Mortimer~\cite{dixon1996permutation}
and by Cameron~\cite{cameron_1999}. 
Interestingly, permutation group theory, in its historical development, preceded semigroup theory, 
see~\cite{Hollings2009}.

Before the middle of the 19th century, 
groups were synonymous to permutation groups on a finite set
and were simply thought of as non-empty subsets of permutations closed under function composition (called substitution groups back then). The concept of an abstract group slowly emerged.
Cayley~\cite{Cayley54a,Cayley54b,Cayley59} was the first who
viewed groups more abstract solely in terms of laws, and not in some permutation representation.
However, his conceptualization was still far away from our modern view and still relied
on implicit assumptions. Weber gave an axiom system
for finite groups~\cite{Weber1882_MathAnn} and for general groups~\cite{Weber1896_LehrbuchAlgebra_Vol2}. 
During the 19th century many further famous mathematician besides the ones already mentioned 
were involved into forming
the abstract theory of groups as, for example, Cauchy, Kronecker, Dyck, Frobenius, and Hölder
to name a few. The first definition of an abstract group
in its modern form appears in van der Waerden's textbook~\cite{vanderWaerden_ModernAlgebra_SpringerArchive}, which is based
on lectures by Artin and Noether from the beginning of the 20th century.
For more information on the history and genesis of abstract group theory, see~\cite{Kleiner2007_HistoryAbstractAlgebra,Neumann_AbstractGroups_MacTutor,Wussing1984_GenesisAbstractGroup}.

Not all authors use the term $G$-congruence (some just call them equivalence relations preserved by the group), 
but it is a notion that appears in the literature, for example in~\cite{Burness14}.
The generalization of it to transformation semigroups is obvious, but I have not
seen in the literature. In~\cite{KilpKnauerMikhalev2000_MonoidsActs}, in the more general context
of semigroup actions, the authors introduce the term $S$-act congruences,
and $G$-congruences are also used in the more general context of group actions, also called $G$-sets,
for example, in~\cite{NeumannStoyThompson1994GroupsAndGeometry}.

The books~\cite{Seress2003PermutationGroupAlgorithms,Butler1991AlgorithmsPermutationGroups,Holt2005_ComputationalGroupTheory} contains information on algorithms for groups and permutation groups. Beaudry~\cite{Beaudry1989CharacterizationIdempotentTM} 
investigated the computational complexity of deciding whether all transformations in a 
given transformation monoid are idempotent.

\IN{point-symmetric graph}
Vertex-transitive graphs seemed to be introduced
by Sabidussi~\cite{sabidussi1964vertex}, at least this is the first reference
that we were able to find where this term was introduced.
Vertex-transitive graphs are also called \emph{point-symmetric} 
in~\cite{harary1969graph}.
The study of symmetry in graphs was initiated 
by Foster~\cite{foster1932geometrical}.
More information on vertex-transitive graphs
can be found in~\cite{BabaiSurvey96,Biggs1993,Cameron2004_AutomorphismsOfGraphs,PraegerLiNiemeyer1997_FiniteTransitivePermutationGroups}.



\chapter{Notions From Automata Theory}
\label{prelim:sec:automata}

%!TEX root = ./../../springer-mono-format/book.tex
%!TEX root = ./../../thesis-trier.tex

%Further automata-theoretical notions are introduced in \cref{prelim:sec:automata},
%in particular various constructions and automata as language acceptors.
%%most of this section is devoted to introducing various constructions
%%as sub-, (Rees) factor, product, power, and pair automata. We also introduces automata as language acceptors 
%%and regular languages.

%\Cref{prelim:remarks_trans_semigrp_semiautomata,prelim:sec:remarks-reduction,prelim:sec:remarks_func_comp} 
%discusses conventions and notational choices.

Recall the definition of a (semi-)automaton and the set of synchronizing words
from \cref{intro:sec:basic,intro:sec:defsyncaut}.
Here, we give additional background on automata. %introduce additional notions. 
\cref{prelim:subsec:morenotions} 
collects more basic notions. The reachability quasiorder and related
notions are introduced in~\cref{prelim:subsec:reachorder}.
\Cref{prelim:subsec:transsemigrp} introduces
the transformation semigroup of an automaton
and a representation by transition matrices.
\Cref{prelim:subsubsec:endo_aut} investigates morphisms
on automata. We introduce endo-, iso- and automorphisms, 
and collect several properties, facts, and formulas; in particular
how morphisms act w.r.t.~the reachability order, strongly connected
components, and upper sets.
% The following sections are devoted
Then a sequence of sections follows
devoted to various constructions. 
Specifically, \SEC~\ref{prelim:sec:aut:sub}, \ref{prelim:sec:aut:factor}, \ref{prelim:sec:rees}, \ref{prelim:sec:productaut},
\ref{prelim:sec:aut:power}, and \ref{sec:pair_aut}
introduce the sub- and restricted, quotient, Rees quotient, product, power, and pair automaton, respectively.
\Cref{prelim:sec:aut:sequppersets} discusses
the set of maximal states and shows ways to decompose an automaton
into a sequence of upper sets. In~\cref{prelim:sec:autacceptors} 
we connect automata to formal languages.
We introduce regular languages
as languages accepted by automata when adjoined with a start state
and a set of final states (which we then call automata
as language acceptors).
\Cref{prelim:subsubsec:min_aut} investigates the minimal automaton
accepting a given regular language.

\section{More Basic Notions}
\label{prelim:subsec:morenotions}

\input{../part2/automata/basic}

\section{The Transformation Semigroup of an Automaton}
\label{prelim:subsec:transsemigrp}

\input{../part2/automata/transform}

\section{The Reachability Order}
\label{prelim:sec:reach} 
\label{prelim:subsec:reachorder}

\input{../part2/automata/reach}

\section{Morphisms} 
\label{prelim:subsubsec:endo_aut}

\input{../part2/automata//morphism}

%\section{Automata Constructions}
%\section{More Operations on Automata}
\section{Operations on Automata}

\subsection{Subautomata and Restricted Automata}
\label{prelim:sec:aut:sub}

\input{../part2/automata/subautomata}

\subsection{The Quotient Automaton}
\label{prelim:sec:aut:factor}

\input{../part2/automata/factor}

\subsection{The Rees Quotient Automaton}
\label{prelim:sec:rees} 

\input{../part2/automata/rees}

\subsection{The Product Automaton} 
\label{prelim:sec:productaut} 

\input{../part2/automata/product}

\subsection{The Power Automaton} 
\label{prelim:sec:aut:power}

\input{../part2/automata/power}

\subsection{The Pair Automaton}
\label{sec:pair_aut}

\input{../part2/automata/pair}

\section{Sequences of Upper Sets and Maximal States}
\label{prelim:sec:aut:sequppersets}

\input{../part2/automata/seq}

\section{Automata as Language Acceptors} 
\label{prelim:sec:autacceptors}

\subsection{Regular Languages}
\label{prelim:sec:regularlang}

\input{../part2/automata/acceptor}

\subsection{The Minimal Automaton}
\label{prelim:subsubsec:min_aut}

\input{../part2/automata/minimal}

\section{Remarks on Transformation Semigroups and Automata}
\label{prelim:remarks_trans_semigrp_semiautomata}

\input{../part2/automata/remark_transsemigrp}

\section{Discussion} 
\label{prelim:automata:discussion}

\input{../part2/automata/discussion-aut} 

\chapter{Classes of Automata}
\label{chap:classes_aut}

%!TEX root = ./../../springer-mono-format/book.tex
%!TEX root = ./../../thesis-trier.tex

Here, we discuss the following classes of automata
that appear throughout the present work
and collect results of relevance to us:

\begin{enumerate}
	\item synchronizing automata (\cref{prelim:sec:sync_aut}) and synchronizing automata with a trap state (\SEC~\ref{prelim:sync:trap}),
	
	\item unary automata (\cref{prelim:subclasses:ternarybinaryunary}),
				
	\item automata with one-point singular idempotents (OI-automata) and automata with simple idempotents (SI-automata) (\cref{prelim:subclasses:SI}),
	
	\item permutation-reset automata (\cref{prelim:subclasses:permutation-reset}),
		
	\item monotonic automata (\cref{prelim:subclasses:monotonic}),
	
	\item automata whose strongly connected components are blocks (\cref{prelim:subclasses:block}),
		
	\item partially ordered automata (\cref{prelim:sec:pto}), 
	
	\item $\rho$-automata (\cref{prelim:subclasses:rho-automata}),
	
	\item commutative automata (\cref{prelim:subclasses:commutative}),
	    
	\item solvable automata (\cref{prelim:subclasses:solvable}),
	 which generalize commutative and partially ordered automata.
\end{enumerate}

See \cref{prelim:classes:fig:inclusions} for a Hasse diagram of the inclusions that hold among these classes.

\begin{figure}[ht]
	\centering
	\begin{tikzpicture}
		\matrix (A) [matrix of nodes, row sep=0.8cm, nodes={minimum width=2cm}]
		{
			OI-Aut. &                  & Str.~Conn.~Comp.~Blocks & Monotonic & Synchronizing \\
			SI-Aut. & Perm.-Reset Aut. & Solvable & & Sync.~with trap state \\
			&                  & $\rho$-Automata. \\
			&  Perm.~Aut       & Commutative & Part.~Ordered Aut. & \\
			&                  & Unary \\
		};
		\path (A-1-1)--(A-1-3) node[above=1.25cm] (link) {All Automata};
		
		\foreach \i in {1,3,4,5}
		\path (link.south) edge[-] (A-1-\i.north);
		
		\path (A-2-1.north) edge[-] (A-1-1.south)
		(A-2-2.north) edge[-] (A-1-3.south)
		(A-2-3) edge[-] (A-1-3.south);       
		
		\path (A-3-3.north) edge[-] (A-2-3.south);
		
		\path (A-4-2.north) edge[-] (A-3-3.south)
		(A-4-2.north) edge[-] (A-2-2.south)
		(A-4-2.north) edge[-] (A-2-1.south)
		(A-4-3.north) edge[-] (A-3-3.south)
		(A-4-4.north) edge[-] (A-3-3.south);
		
		\path (A-5-3.north) edge[-] (A-4-3.south);
		
		\path (A-2-5.north) edge[-] (A-1-5.south);
	\end{tikzpicture} 
	\caption[Relations among classes of automata]{Inclusion relations among the classes of automata from the present section.}
	\label{prelim:classes:fig:inclusions}
\end{figure}

The following content is new and did not appear before in the literature:

\begin{enumerate}
\item the novel class of $\rho$-automata,

\item the results on automata whose strongly connected components are blocks,

\item the structural results on solvable automata and the result that a solvable automaton admitting
a word of rank $r$ has a word of rank at most 
$r$ and of length at most $n - r$ and 

\item new structural results on solvable automata.
\end{enumerate}

The remaining results presented are known and taken from the literature (maybe presented a little differently).

Other classes of automata, for example circular automata, completely
reachable, completely distinguishable or
sync-maximal automata are only used in Part~\ref{part:aut} and so not introduced
here but separately in Part~\ref{part:aut}.

\section{Synchronizing Automata}
\label{prelim:sec:sync_aut} 

\input{../part2/classes/synchronizing}

\section{Unary Automata}
\label{prelim:subclasses:ternarybinaryunary}

\input{../part2/classes/ternarybinaryunary}

\section{OI- and SI-Automata}
\label{prelim:subclasses:SI}
	
\input{../part2/classes/SI}

\section{Permutation-Reset Automata}
\label{prelim:subclasses:permutation-reset} 

\input{../part2/classes/permutation}

\section{Monotonic Automata} 
\label{prelim:subclasses:monotonic}

\input{../part2/classes/monotonic}

\section{Automata Whose Strongly Connected Components are Blocks}
\label{prelim:subclasses:block} 

\input{../part2/classes/block}

\section{Partially Ordered Automata}
\label{prelim:sec:pto} 

\input{../part2/classes/pto}

\section{\texorpdfstring{$\rho$}{rho}-Automata}
\label{prelim:subclasses:rho-automata} 

\input{../part2/classes/rho}

\section{Commutative Automata}
\label{prelim:subclasses:commutative}

\input{../part2/classes/commutative}

\section{Solvable Automata}
\label{prelim:subclasses:solvable} 

\input{../part2/classes/solvable}

\section{Discussion}
\label{prelim:subclasses:discussion}

\input{../part2/classes/discussion}

\chapter{Subset Mapping Problems}
\label{chap:prelim_complexity_results}


In this chapter, we investigate various subset mapping and synchronization problems on automata,
which we introduce in \cref{compl:sec:problems}.
\Cref{compl:sec:reductions} explores
reductions and relations between them.
In \cref{compl:sec:subclasses} we investigate the computational
complexity of these problems on general automata
and the subclasses of solvable, commutative, and partially
ordered automata. We give a classification of the (classical) 
complexity landscape and close some gaps from the literature.

\section{The Problems}
\label{compl:sec:problems}

\input{../part2/subset/problems}

\section{Reductions}
\label{compl:sec:reductions}

\input{../part2/subset/reductions-new}

\section{Classes of Automata}
\label{compl:sec:subclasses}

\input{../part2/subset/classes}

\partwithquote{Automata for Sets of Synchronizing Words}{``My best tool for efficiency and performance is abstraction.''}{---Bjarne Stroustrup (1950--present)}
\label{part:aut}

%\chapter{Subset Mapping Properties} % in Automata and Transformation Semigroups} 
\chapter{Reachable Sets and Pairs of Distinguishable Sets}
\label{sync_max:chap:prelim}

%!TEX root = ./../../thesis-trier.tex

In this chapter, we investigate the subset mapping properties in automata and transformation semigroups, motivated by minimality conditions for automata accepting sets of synchronizing words. The results presented here can be combined to derive lower bounds for the size of these automata.

In \cref{sync_max:sec:power_automat_for_SynA},
we introduce and motivate the notions of reachability
and distinguishability. 
We do this by first observing
that the power automaton accepts the set of synchronizing,
and then naturally obtain both notions
via minimality conditions for accepting automata from \cref{prelim:subsubsec:min_aut} applied to the power automaton. 
Thus, these notions derive from necessary conditions
characterizing the size of minimal automata for sets of synchronizing words.

\Cref{sync_max:sec:reachable} investigates the reachability of subsets of states in automata.
We derive synchronizability conditions and 
structural limitations in terms of the reachability
of families of subsets of states.

\Cref{sync_max:sec:distinguish}
is concerned with distinguishability,
where we investigate
its relation to synchronizability and 
consider different families of pairwise distinguishable subsets
and their relations. We also introduce the novel notion of completely distinguishable automaton.

Reachability and distinguishability transfer readily from automata
to transformation semigroups. \Cref{sync_max:sec:from_aut_to_semigroups}
shortly discusses this and some differences that arise as a transformation
semigroup might not contain the identity transformation.

\section{Power Automata and Sets of Synchronizing Words}
\label{sync_max:sec:power_automat_for_SynA}

\input{../part3/subset/poweraut}

\section{Reachable Subsets}
\label{sync_max:sec:reachable}

\input{../part3/subset/reach}

\section{Pairs of Distinguishable Subsets} % in Automata and Transformation Semigroups}
\label{sync_max:sec:distinguish}

\input{../part3/subset/pairsdist}

\section{From Automata to Transformation Semigroups}
\label{sync_max:sec:from_aut_to_semigroups} 

\input{../part3/subset/tsemigrps}

\section{Discussion}
\label{sync_max:chap_subset_mapping:discussion}

\input{../part3/subset/discussion}

\chapter{Sync-Maximal Automata}
\label{sync_max:chap:sync_max}


 Here, in \cref{sync_max:sec:sync_max_aut}, we introduce sync-maximal automata as automata
 whose sets of synchronizing words have maximal state complexity.
 Then, in \cref{sync_max:sec:structure}, we give structural results.
 The main result is \cref{sync_maximal:thm:structure}, 
 which states that these automata contain
 a completely reachable and completely distinguishable subautomaton
 that differs from the whole automaton by at most one ``dangling'' state.
 Lastly, \cref{sync_max:sec:examples} provides 
 various examples illustrating 
 relations between sync-maximal, completely distinguishable, and completely reachable automata.
 For instance, we present sync-maximal automata that are not completely reachable, as well as completely reachable automata that are not sync-maximal, along with other examples illustrating our results.
 
 \section{Definition and First Properties}
\label{sync_max:sec:sync_max_aut}
	
\input{../part3/syncmax/sync_max_intro}

\section{The Structure of Sync-Maximal Automata}
\label{sync_max:sec:structure}

\input{../part3/syncmax/sync_max}

\section{Examples} 
\label{sync_max:sec:examples}

\input{../part3/syncmax/examples}

\section{Discussion}

\input{../part3/syncmax/discussion}

\chapter{Graph-Theoretical Constructions}
\label{sync_max:chap:graphs}

%!TEX root = ./../../thesis-trier.tex

Here, we introduce digraphs
associated to automata and transformation semigroups 
that we use to deduce complete reachability and complete distinguishability.
Section~\ref{sync_max:sec:BVC-graphs},~\ref{sync_max:sec:augmentative}, and~\ref{sync_max:sec:rystsov_graph}
review existing constructions from the literature.
We introduce the Bondar-Casas-Volkov graphs, which apply to general automata, and the Rystsov graphs, which only apply to SI-automata,
from the literature.
\cref{sec:extended_Rystsov_graph} extends these graphs to arbitrary transformation semigroups.
Then we investigate the influence of the group of units
on the introduced digraphs (\cref{sync_max:graphs:sec:grpunits}),
relations between them (Section~\ref{sync_max:graphs:sec:relations}), give 
examples that highlight their differences (\cref{sync_max:graphs:sec:examples}), 
and finally formulate sufficient conditions in terms of connectedness notions to deduce
complete reachability, complete distinguishability, or
 lower bounds for the size of accepting automata of sets of synchronizing words (Section~\ref{sync_max:graphs:sec:cr_cd}).
 In \cref{sync_max:graphs:sec:cr_cd}, we also point out some
 relations to congruences and minimal automata.
 
 \section{The Bondar-Casas-Volkov Graphs}
\label{sync_max:sec:BVC-graphs} 

\input{../part3/graphs/bcvgraphs}

\section{Augmentative Sets of Words}
\label{sync_max:sec:augmentative}

\input{../part3/graphs/augmentative}

\section{The Rystsov Graph}
\label{sync_max:sec:rystsov_graph} 

\input{../part3/graphs/rystsov}

\section{From Automata to Transformation Semigroups}
\label{sec:extended_Rystsov_graph} 

\input{../part3/graphs/transsemigrp}

\section{Relations Between the Different Graphs}
\label{sync_max:graphs:sec:relations}

\input{../part3/graphs/relations}

\section{Relations to the Group of Units}
\label{sync_max:graphs:sec:grpunits}

\input{../part3/graphs/grpunits}

\section{Examples}
\label{sync_max:graphs:sec:examples}

\input{../part3/graphs/examples}

%\section{Further Properties}
%\section{Properties of the Graphs}
%\label{sync_max:graphs:sec:properties}

\input{../part3/graphs/properties}

\section{Deducing Reachability and Distinguishability}
\label{sync_max:graphs:sec:cr_cd}

\input{../part3/graphs/deducingrd}

\section{A Bound on the Reset Length for OI-Automata}
\label{sync_max:graphs:sec:extend-bound}

\input{../part3/graphs/extendbound}

\section{Discussion}

\input{../part3/graphs/discussion}

\chapter{Classes of Automata}
\label{sync_max:sec:subclasses}

%!TEX root = ./../../springer-mono-format/book.tex
%!TEX root = ./../../thesis-trier.tex

Here, we investigate the following subclasses of automata with respect to the notions
of complete distinguishability, complete reachability,
and sync-maximality: solvable automata (\Cref{sync_max:sec:blocks}) 
commutative automata (\Cref{sync_max:sec:comm}), automata where every letter has defect at most one (\Cref{sync_max:sec:defectone}), SI-automata (Section~\ref{sync_max:sec:SI}),
circular automata (Section~\ref{sync_max:sec:circular}), where \cref{sync_max:sec:circular} 
also introduces circular automata, and binary automata (\Cref{sync_max:sec:binary}).

\section{Automata Whose Strongly Connected Components are Blocks}
\label{sync_max:sec:blocks}

\input{../part3/classes/blocks}

\section{Commutative Automata}
\label{sync_max:sec:comm}

\input{../part3/classes/comm}

\section{All Letters Have Defect at Most One}
\label{sync_max:sec:defectone}

\input{../part3/classes/defectone}

\section{Automata with Simple Idempotents}
\label{sync_max:sec:SI}

\input{../part3/classes/simpl_idemp}

\section{Circular Automata}
\label{sync_max:sec:circular}

\input{../part3/classes/circular}

\section{Binary Automata}
\label{sync_max:sec:binary}

\input{../part3/classes/binary}

\section{Discussion}
\label{sync_max:sec:subclasses:discussion}

\input{../part3/classes/discussion}

\chapter{Decision Problems}
\label{sync_max:chap:decision}


In Sections~\ref{sync_max:sec:decision_cd} and \ref{sync_max:sec:decision_cr}
we investigate the complexity of deciding complete distinguishability, complete reachability,
and sync-maximality. 
Also, we prove that the problem to decide whether two subsets are distinguishable is \PSPACE-complete.
In \cref{sync_max:decision:sec:binary-SI}, 
we investigate these problems, and the problem to decide synchronizablity,
for binary SI-automata as input,
which appears in a separate section as  
complete reachability, complete distinguishability, and sync-maximality
are equivalent properties for binary SI-automata by \cref{sync_max:prop:binary-SI}.
An overview of our results and those from the literature on complete reachability, complete distinguishability, and sync-maximality,
is given in \cref{tab:complexity}.

% https://tex.stackexchange.com/questions/478671/how-i-can-write-this-multi-column-and-multi-row-table-in-latex
\begin{table}[htb]
	\centering    
	\caption[Overview decision problems]{Overview of results on deciding
		complete reachability, complete distinguishability, and sync-maximality.
		\NL-hardness holds even for SI-automata without permutational letters.}
	\begin{tabular}{*{5}{l}}
		\toprule
		\multirow{2}{*}{Problem} & \multicolumn{4}{c}{Subclass of Automata} \\
		\cmidrule(lr){2-5}
		& SI-Automata & \parbox{1.9cm}{Binary Automata} & \parbox{1.9cm}{Binary SI-Automata} & \parbox{1.9cm}{All \\ Automata} \\  
		\midrule 
		\\ [-1.5ex] 
		Complete Reach. & \parbox{1.8cm}{\NL-c \\ (Th.~\ref{sync_max:thm:SI-aut_compl_reach_NL-c})} & \parbox{1.5cm}{ \PTIME \\ \cite{DBLP:conf/latin/CasasV22}} & \parbox{1.8cm}{\LOGSPACE-c \\ (Th.~\ref{syncmax:SI-aut:prop:sync-L-complete})}  & \parbox{1.5cm}{\PTIME \\ \cite{DBLP:conf/icalp/FerensS23}}  \\ [2.5ex] 
		Complete Dist. & \parbox{1.8cm}{\NL-hard \\ (Th.~\ref{thm:cd_NL-hard})} & Open & \parbox{1.8cm}{\LOGSPACE-c \\ (Th.~\ref{syncmax:SI-aut:prop:sync-L-complete})} & \parbox{1.9cm}{\PSPACE-c \\ \cite{SzykulaZyzik24}}  \\ [2.5ex] 
		Sync-Maximality & \parbox{1.8cm}{\NL-hard \\ (Th.~\ref{thm:sync_max_NL-hard})} & \parbox{2.7cm}{\PTIME \\ (by Cor.~\ref{cor:char-bin-compl-reach})}   & \parbox{1.8cm}{\LOGSPACE-c \\ (Th.~\ref{syncmax:SI-aut:prop:sync-L-complete})} & \parbox{1.9cm}{\PSPACE-c \\ \cite{SzykulaZyzik24}}   \\  [2.5ex] 
		\bottomrule
	\end{tabular} % open rein
	\label{tab:complexity}
\end{table}

\todo[inline]{personal communication im fliestext}

With the correspondence between automata and transformation semigroups (\cref{prelim:remarks_trans_semigrp_semiautomata}), 
the same complexities for the respective problems for transformation semigroups given as a list of generators follow.

\section{Distinguishability}
\label{sync_max:sec:decision_cd} 

\input{../part3/decision/decision_cd}

\section{Complete Reachability and Sync-Maximality}
\label{sync_max:sec:decision_cr}

\input{../part3/decision/decision_cr}

%\section{Sync-Maximality}
%\label{sync_max:subsec:complexity}
%\input{part3/decision/decision_sm}

\section{Binary SI-Automata}
\label{sync_max:decision:sec:binary-SI}

\input{../part3/decision/binarysi}

\section{Discussion} 
\label{sync_max:decision:discussion}

\input{../part3/decision/discussion}







\chapter{Permutation Groups}
\label{sync_max:sec:perm_grps}

%!TEX root = ./../../thesis-trier.tex
 
Here, we investigate the properties of synchronizability, reachability, and distinguishability
in transformation monoids generated by a permutation group
and some non-permutation.
\cref{sec:prim_groups}
recaps some known characterizations for primitive groups, and gives new characterizations
in terms of minimality, synchronizability, complete reachability,
and complete distinguishability of associated automata. 
The subclass of synchronizing groups
is introduced in \cref{sync_max:perm:sec:sync_grps}. This section also discusses $k$-homogeneous groups.
\Cref{sec:A-primitive}
introduces the new notion of $A$-primitivity for a subset $A$
of the permutation domain as a generalization of primitivity.
We show that several properties of primitive groups carry over to $A$-primitive groups.

\section{Primitive Permutation Groups}
\label{sec:prim_groups}

\input{../part3/permgrps/primgrps}

\section{Synchronizing and $k$-Homogeneous Groups}
\label{sync_max:perm:sec:sync_grps} 

\input{../part3/permgrps/syncgrp}

\section{\texorpdfstring{$A$}{A}-Primitive Permutation Groups}
\label{sec:A-primitive} 

\input{../part3/permgrps/aprimgrps}

\section{Congruences on Transformation Semigroups}
\label{syncmax:perm:sec:congruences}

\input{../part3/permgrps/congruences}

\section{Discussion}
\label{sync_max:app:discussion}

\input{../part3/permgrps/discussion}

\chapter{Slowly Synchronizing Automata}
\label{sync_max:chap:slowly}

%!TEX root = ./../../../thesis-format/main.tex
	
\IN{slowly synchronizing automaton}
\IN{slowly synchronizing family of automata}
In~\cite{DBLP:conf/mfcs/AnanichevGV10,KisielewiczS15} an automaton
was called \emph{slowly synchronizing}
if ``[...][its] reset length [is] close to 
the \v{C}ern\'y bound [...]''~\cite{DBLP:conf/mfcs/AnanichevGV10}. 
We adhere to this intuitive ``non-quantitative'' definition.
%To exclude 
%certain modifications (similar as in \cref{syncmax:slowly:ex}), 
In~\cite{AnaVolGus2013}
the authors focused on slowly synchronizing $n$-state automaton with reset
lengths between $(n-2)^2 + 2$ and $(n-1)^2$.
In~\cite{Martugin08,DBLP:journals/jalc/Volkov19a}
slowly synchronizing automata whose reset lengths
do not fall within these bounds have been given.
%Here, we adopt the first definition, which is, strictly
%speaking, more an intuitive statement as ``closeness'' is not precisely defined.
%The examples we present
%also obey the more precise definition.
As we are dealing mainly with families of automata
and to formulate two conjectures, %in \cref{sync_max:sec:slowly_sync_discussion}, 
we introduce \emph{slowly synchronizing families of automata}
as infinite families of distinct automata $\semiautomat{A_n}{Q_n}{\Sigma}{\delta_n}$ 	
such that $\rt{\mathcal A_n} \ge C\cdot |Q_n|^2$
for some $C > 0$ and all $n \ge 1$ (most of the time we have $|Q_n| = n$).\footnote{We denote
the alternating group also with $\mathcal A_n$, see \pageref{def:A_n}.
But there should no confusion arise since the context makes it clear whether
we denote an automaton family or the alternating group by~$\mathcal A_n$.} 

%infinitely many $n \ge 1$. 
For example, the \Cerny family~$\mathcal C_n$ from 
\cref{intro:sec:cerny} (\cref{intro:fig:cerny})
is slowly synchronizing.
%Without having a precise definition, 
For any reasonable interpretation of slowly synchronizing automata,
we would expect every family of slowly synchronizing automata
to give a slowly synchronizing family of automata.

In \cref{sync_max:sec:cerny_bound}, 
we discuss automata whose reset lengths reach 
the \v{C}ern\'y bound. In \cref{sync_max:sec:slowly_sync}
we investigate certain known families of slowly synchronizing automata
and observe that for them
%we investigate slowly synchronizing automata and observe that for
%the ones we investigate, 
minimal automata for their sets of synchronizing words are exponentially larger (we also confirm
sync-maximality for the automata~$\mathcal K_n$ for which this was previously only conjectured).
\Cref{sync_max:sec:slowly_sync_discussion} concludes by discussing and formulating conjectures based
on our observations.
%based
%on notions that exclude artificial examples that arise out of adding letters.

\section{Automata Reaching the \Cerny Bound}
\label{sync_max:sec:cerny_bound}

Besides the \Cerny family from \cref{intro:fig:cerny},
for $n \ge 3$,
only the eight exceptional synchronizing $n$-state automata from \cref{sync_max:fig:exceptional_automata}
are known to have reset length $(n-1)^2$
and the property of becoming non-synchronizing after the removal of any letter~\cite{DonZB21,DBLP:conf/lata/Volkov08}.
These exceptional automata are given in \cref{sync_max:fig:exceptional_automata}.
The automaton $\mathbf{K}$ was
discovered by Kari~\cite{JKari01}, the automaton~$\mathbf{I}$ by Roman~\cite{DBLP:journals/jalc/Roman08},
and the remaining six automata by Trahtman via a computer search~\cite{DBLP:conf/mfcs/Trahtman06}.

By extensive compute analysis, Don et al~\cite{DonZB21} 
found $19$ additional automata with at least three and at most six states reaching the \v{C}ern\'y bounds,
but they do not have the property of becoming
non-synchronizing after removal of any letter. In this sense, with these $8$ exceptional automata
and the~$4$ automata from the \Cerny family we have $12$ automata
with at least two and at most six states that reach the \Cerny bound
and are minimal w.r.t.~the alphabet.

Here, we investigate the eight automata from \cref{sync_max:fig:exceptional_automata} w.r.t.~sync-maximality.

\begin{figure}
	\centering
	\begin{tabular}{ccc}
		\begin{tikzpicture}
			\node[state] (0) {$0$};
			\node[state, above left  = of 0] (1) {$1$};
			\node[state, above right = of 0] (2) {$2$};
			
			\path[->] (1) edge [bend left,above] node {$a$} (2)
			(2) edge [bend left,right] node {$a$} (0)
			(0) edge [bend left,left] node {$a,c$} (1);
			\path[->]  (1) edge [bend left,above,pos=.7] node {$c$} (0)
			(2) edge [above] node {$c$} (1);
		\end{tikzpicture}   & &
		
		\begin{tikzpicture}
			\node[state] (0) {$0$};
			\node[state, above left  = of 0] (1) {$1$};
			\node[state, above right = of 0] (2) {$2$};
			
			\path[->] (1) edge [bend left,above] node {$a,b$} (2)
			(2) edge [above] node {$a,b$} (1);
			\path[->] (0) edge [above,pos=.4] node {$c$} (2)
			(2) edge [bend left,below] node {$c$} (0);
			\path[->] (0) edge [above,pos=.4] node {$a$} (1)
			(0) edge [loop left] node {$b$} (0)
			(1) edge [loop left] node {$c$} (1);			 
		\end{tikzpicture}
		
		\\ [2pt]
		The automaton \textbf{D} & & The automaton \textbf{E} \\ [2pt]
		
		\begin{tikzpicture}
			\node[state] (0) {$0$};
			\node[state, right = of 0] (1) {$1$};
			\node[state, right = of 1] (2) {$2$};
			
			\path[->] (0) edge [loop above] node {$c$} (0)
			          (1) edge [loop above] node {$a$} (1)
			          (1) edge [bend left,above]  node {$c$} (2)
			          (2) edge [loop above] node {$a,b$} (2);
			          
			\path[->] (2) edge [bend left,below] node {$c$} (1)
			 (1) edge [bend left,below] node {$b$} (0)
			 (0) edge [bend left,above] node {$a,b$} (1);
 		\end{tikzpicture} & &
		
		\begin{tikzpicture}
			\node[state] (0) {$0$};
			\node[state, right = of 0] (1) {$1$};
			\node[state, right = of 1] (2) {$2$};
			\node[state, above = of 1] (3) {$3$};
			
			\path[->] (0) edge [loop above] node {$a$} (0)
			 (0) edge [above] node {$b$} (3)
			 (1) edge [above] node {$b$} (0)
			 (3) edge [right] node {$b$} (1)
			 (1) edge [bend left,above] node {$a$} (2)
			 (2) edge [bend left,below] node {$a$} (1)
			 (2) edge [loop above] node {$b$} (2)
			 (3) edge [above] node {$a$} (2);
			
		\end{tikzpicture} \\ [2pt]
		
		The automaton \textbf{F} & & The automaton \textbf{G} \\ [2pt]
		
		\begin{tikzpicture}
			\node[state] (2) {$2$};
			\node[state, above left  = of 2] (0) {$0$};
			\node[state, below left  = of 2] (1) {$1$};
			\node[state, right = of 2]       (3) {$3$};
			
			\path[->] (0) edge [left] node {$a$} (1)
			(1) edge [bend left,below] node {$a$} (2)
			(2) edge [bend left,below] node {$a$} (1)
			(3) edge [loop above] node {$a,b$} (3);
			
			\path[->] (0) edge [loop left] node {$c$} (0)
			  (1) edge [loop left] node {$b,c$} (1);
			  
			\path[->] (2) edge [bend left,above] node {$c$} (3)
			 (3) edge [bend left,below] node {$c$} (2)
			 (2) edge [bend left,above] node {$b$} (0)
			 (0) edge [bend left,above] node {$b$} (2);
		\end{tikzpicture} & &
		
		\begin{tikzpicture}
			\node[state] (0) {$0$};
			\node[state, right = of 0] (1) {$1$};
			\node[state, right = of 1] (2) {$2$};
			\node[state, above = of 0,yshift=1.4cm] (3) {$3$};
			
			%\path[<->] (0) edge [above] node {$c$} (3); % zwei pfeile
			\path[->] (0) edge [bend left,left] node {$c$} (3)
			 (3) edge [left] node {$c$} (0);
			
			\path[->] (3) edge [left,pos=0.7] node {$a,b$} (1);
			\path[->] (1) edge [bend right, left] node {$b$}   (3);
			\path[->] (0) edge [loop left] node {$b$} (0)
			          (0) edge [above] node {$a$} (1);
			\path[->] (1) edge [bend left, above] node {$a,c$} (2)
			 (2) edge [bend left, below] node {$c$} (1)
			 (2) edge [loop above] node {$b$} (2)
			 (2) edge [bend right,above] node {$a$} (3);
			 			
		\end{tikzpicture} \\ [2pt]
		The automaton \textbf{H} & & The automaton \textbf{I} \\ [2pt]
		
		\begin{tikzpicture}
			\node[state] (0) {$0$};
			\node[state, above right = of 0] (1) {$1$};
			\node[state, below right = of 0] (2) {$2$};
			\node[state, right = of 1] (4) {$4$};
			\node[state, right = of 2] (3) {$3$};
			
			\path[->] (0) edge [loop above] node {$c$} (0)
			(1) edge [bend left,above] node {$c$} (4)
			(4) edge [bend left,below] node {$c$} (1)
			(2) edge [bend left,above] node {$c$} (3)
			(3) edge [bend left,below] node {$c$} (2);
			
			\path[->] (0) edge [above] node {$b$} (2)
			(2) edge [bend left,below]  node {$b$} (0)
			(3) edge [loop right] node {$a,b$} (3)
			(4) edge [loop right] node {$a,b$} (4)
			(1) edge [loop left] node {$b$} (1);
			
			\path[->] (0) edge [above,below] node {$a$} (1)
			(1) edge [bend left,right] node {$a$} (2)
			(2) edge [left] node {$a$} (1);			 
		\end{tikzpicture} & &
		
		\begin{tikzpicture}
		%	\tikzstyle{state}=[minimum size=25pt,inner sep=1pt];
			\node[state] (0) {$0$};
			\node[state, right = of 0] (1) {$1$};
			\node[state, right = of 1] (2) {$2$};
			\node[state, below = of 2, yshift=-1.3cm] (3) {$3$};
			\node[state, left = of 3]  (4) {$4$};
			\node[state, left = of 4]  (5) {$5$};			
			
			\path[->] (0) edge [above] node {$a$} (1) 
			 (1) edge [above] node {$a$} (2) 
			 (2) edge [bend right,above] node {$a$} (0);
			 
			\path[->] (5) edge [above] node {$a$} (4)
			 (4) edge [above] node {$a$} (3)
			 (3) edge [bend left, below] node {$a$} (5);
			 
			\path[->] (0) edge [loop above] node {$b$} (0)
			 (1) edge [loop below] node {$b$} (1)
			 (2) edge [bend left, right]  node {$b$} (3)
			 (3) edge [bend left, right]  node {$b$} (2)
			 (4) edge [left] node {$b$} (2)
			 (5) edge [loop above] node {$b$} (5);
			
		\end{tikzpicture} \\ [2pt]
		
		Roman's automaton \textbf{J} & & Kari's automaton \textbf{K}.
		
	\end{tabular}
	\caption[Eight exceptional automata attaining the \v{C}ern\'y bound]{The eight exceptional automata whose reset lengths equal the \v{C}ern\'y bound $(n-1)^2$. 
	The naming is taken from~\cite{HeCLS21} (a mistake for \textbf{J} in~\cite{HeCLS21} has been fixed here, the correct automaton is written in~\cite{DBLP:conf/lata/Volkov08}).} \label{def:exceptional_aut}
	\label{sync_max:fig:exceptional_automata}
\end{figure}

\begin{figure}
\centering 
\scalebox{0.975}{
\begin{tikzpicture}[shorten >=1pt,->,node distance=1.25cm]
	\tikzstyle{state-subset}=[]
	
	\node[state-subset, initial, initial where=right] (0123) {$\{0,1,2,3\}$};
	\node[state-subset, below = of 0123, yshift=0.5cm] (123) {$\{1,2,3\}$};
	\node[state-subset, below = of 123]  (012) {$\{0,1,2\}$};
	\node[state-subset, right = of 012]  (023) {$\{0,2,3\}$};
	\node[state-subset, above = of 023]  (013) {$\{0,1,3\}$};
	\node[state-subset, right = of 023]  (13)  {$\{1,3\}$};
	\node[state-subset, above = of 13]   (12)  {$\{1,2\}$}; 
	\node[state-subset, right = of 12]   (23)  {$\{2,3\}$};
	\node[state-subset, right = of 13]   (02)  {$\{0,2\}$};
	\node[state-subset, right = of 23]   (01)  {$\{0,1\}$};
 	\node[state-subset, below = of 01]   (03)  {$\{0,3\}$};
 	\node[state, fill=white,minimum size=17pt,inner sep=1pt, below = of 03, accepting]   (1)   {$*$}; % genügt, sagen einermenge kollapsed, vielleicht mit * andeuteen
 	
 	\path[->] (0123) edge [right] node {$a$} (123)
 	          (0123) edge [loop left,looseness=2.5] node {$b,c$} (0123);
 	          
 	\path[->] (123) edge [loop left,looseness=2] node {$a,b$} (123)
 	          (013) edge [above] node {$a$} (12)
 	          (023) edge [above] node {$a$} (13);
 	\path[->] (012) edge [bend left,left] node {$a$} (123);
 	
 	\path[->] (12) edge [bend left, above] node {$a$} (23)
 	          (23) edge [bend left, above] node {$a$} (13)
 	          (13) edge [left] node {$a$} (12);
 	          
   \path[->] (01) edge [bend right,above,out=300,in=230] node {$a$} (12)
             (03) edge [right] node {$a$} (1);
   \path[<->] (02) edge [below] node {$c$} (13);
   \path[->]  (02) edge [bend left, below] node {$a$} (13);
             
    \path[<->] (123) edge [right] node {$c$} (012)
               (023) edge [above] node {$b$} (012)
               (013) edge [left] node {$c$} (023);
     
    \path[->] (013) edge [loop left,looseness=2] node {$b$} (013)
              (13)  edge [loop below] node {$b$} (13)
              (12)  edge [loop above] node {$c$} (12);
              
    \path[<->] (23) edge [below] node {$b$} (12)
               (02) edge [loop below] node {$b$} (02)    
               (01) edge [right] node {$b$} (03)
               (01) edge [above] node {$c$} (23);  
    \path[->] (03) edge [loop right,looseness=3] node {$c$} (03);     
    \path[->] (1) edge [loop left] node {$a,b,c$} (1);           
\end{tikzpicture}
}
\caption[Minimal automaton for $\Syn(\mathbf{I})$]{A minimal automaton for the set of synchronizing word of the automaton $\textbf{I}$ from Figure~\ref{sync_max:fig:exceptional_automata}. See the proof of Theorem~\ref{sync_max:thm:extremal_automata} for details.}
\label{sync_max:fig:I} % final und start adden
\end{figure}

\begin{theorem}
\label{sync_max:thm:extremal_automata}
Consider the exceptional automata for the \v{C}ern\'y bound from \FIG~\ref{sync_max:fig:I}.
	The automata \textbf{D}, \textbf{E}, \textbf{I} and \textbf{F} are sync-maximal.
	\textbf{G}, \textbf{H}, \textbf{J} and \textbf{K} are not sync-maximal.
\end{theorem}
\begin{proof} We handle each automaton individually.
		
	\begin{enumerate}
		\item[\textbf{D}, \textbf{E}, \textbf{F}:] In \textbf{D} the letter $a$ generates a primitive permutation group;
		in \textbf{E} and \textbf{F} the letters~$b$ and~$c$ generate the full symmetric group on $\{0,1,2\}$ in the transformation semigroup,
		which is primitive.		
		So, in all cases, sync-maximality follows from \cref{sync_max:thm:prim}.
				
		For \textbf{E}, we also obtain sync-maximality directly by observing that
		$a$ has image $\{1,2\}$, $ac$ has image $\{0,1\}$, $acb$ has image $\{0,2\}$, $acba$ has image $\{1\}$, 
		and since we have the following distinguishing words for all pairs:
		\begin{center}
		\begin{tabular}{l|c|l|c}
			Pair of Subsets & Dist. Word & Pair of Subsets & Dist. Word \\ \hline 
			$\{0,2\}$, $\{1,2\}$    & $a$ & $\{0,1\}$, $\{1,2\}$    & $ba$   \\
			$\{0,2\}$, $\{0,1\}$    & $a$ &	$\{0,1,2\}$, $\{0,1\}$  & $ba$   \\
			$\{0,2\}$, $\{0,1,2\}$  & $a$ &	$\{0,1,2\}$, $\{1,2\}$  & $cba$ \\
		\end{tabular}
	    \end{center}
		For \textbf{E}, we have that $a$, $ab$, $abc, abca$ map $\{0,1,2\}$ to $\{1,2\}, \{0,2\}$, $\{0,1\}$ and $\{1\}$ respectively.
		Furthermore, $a$ distinguishes $\{0,1\}$ form every other $2$-set,
		and $ca$ distinguishes $\{1,2\}$ from $\{0,2\}$.
		
		\item[\textbf{G}:] The following diagram shows the reachable $3$-sets:
		\begin{center}	
			\begin{tikzpicture}[>=stealth, every node/.style={font=\small}]			
				\matrix (m) [matrix of math nodes, column sep=2em] {
					\{0,1,2,3\} & \{0,1,2\} & \{0,2,3\} & \{1,2,3\} \\
				};
				
				\path[->] (m-1-1) edge[above] node {$a$} (m-1-2);
				\path[->] (m-1-2) edge[loop above, in=60, out=120, looseness=3] node {$a$} (m-1-2);
				
				\path[->] (m-1-2) edge[bend left=10,above] node {$b$} (m-1-3);
				\path[->] (m-1-3) edge[bend left=10,below] node {$a$} (m-1-2);
				
				\path[->] (m-1-3) edge[above] node {$b$} (m-1-4);
			\end{tikzpicture}
		\end{center}
		Since $\{0,1,3\}$ is not reachable, the automaton is not sync-maximal by \cref{sync_maximal:thm:structure}.
				
		\item[\textbf{H}:] The following diagram shows all reachable $3$-sets:
		\begin{center}	
		\begin{tikzpicture}[>=stealth, every node/.style={font=\small}]
			
			\matrix (m) [matrix of math nodes, column sep=2em] {
				\{0,1,2,3\} & \{1,2,3\} & \{0,1,3\} & \{0,1,2\} \\
			};
			
			\path[->] (m-1-1) edge[above] node {$a$} (m-1-2);
			\path[->] (m-1-2) edge[loop above, in=60, out=120, looseness=3] node {$a,c$} (m-1-2);
			\path[->] (m-1-2) edge[bend left=10, above] node {$b$}   (m-1-3);
			\path[->] (m-1-3) edge[bend left=10, below] node {$a,b$} (m-1-2);
			\path[<->] (m-1-3) edge[above] node {$c$} (m-1-4);
			\path[->] (m-1-4) edge[loop above, in=60, out=120, looseness=3] node {$b$} (m-1-4);
		\end{tikzpicture}
		\end{center}
		We see that not all $3$-sets are reachable. So, the strongly connected automaton~$\text{H}$ 
		is not sync-maximal  by \cref{sync_maximal:thm:structure}.
		
		\item[\textbf{I}:] 
		The word $b$ induces a transposition and the word $cb$
		induces the cycle $(0 ~ 1 ~ 2 ~ 3)$.
		Thus by \cref{prelim:thm:cycles_gen_Sn} the 
		transformation semigroup contains
		the full symmetric group.
		Since the full symmetric group is primitive,
		from \cref{sync_max:thm:prim} 
		we find that the automaton is sync-maximal.
				
		For an alternative proof, consider the power automaton
		 with the singleton sets merged into a single state and the empty set omitted
		 shown in \cref{sync_max:fig:I}.
		 We see that $\textbf{I}$ is completely reachable.
		 Now, the words $u_1 = aacba$, $u_2 = au_1$, $u_3 = bu_2$, and $u_4 = cu_3$ 
		 map $\{0,1,3\}$, $\{0,2,3\}$, $\{0,1,2\}$, and $\{1,2,3\}$, respectively,
		 to a singleton while mapping the state not in the respective
		 set to another distinct state.
		 So, by \crefitem{sync_max:thm:char_distinguishable}{dist:item:single_point},
		 the automaton \textbf{I} is completely distinguishable.
		 By \cref{sync_maximal:thm:structure}, $\textbf{I}$ is sync-maximal.
		 			
		\item[\textbf{J}:] 
		We have the following transitions from the whole state set 
		(omitting transitions into sets with less than four elements):		
		\begin{center}	
			\begin{tikzpicture}[>=stealth, every node/.style={font=\small}]
				
				\matrix (m) [matrix of math nodes, column sep=2em] {
					\{0,1,2,3,4\} & \{1,2,3,4\} & \{0,1,3,4\} & \{0,1,2,4\} \\
				};
				
				\path[->] (m-1-1) edge[above] node {$a$} (m-1-2);
				\path[->] (m-1-2) edge[loop above, in=60, out=120, looseness=3] node {$a$} (m-1-2);
				\path[<->] (m-1-2) edge[above] node {$b$} (m-1-3);
				\path[<->] (m-1-3) edge[above] node {$c$} (m-1-4);
				\path[->] (m-1-4) edge[loop above, in=60, out=120, looseness=3] node {$b$} (m-1-4);
			\end{tikzpicture}
		\end{center}
		Since $\{0,1,2,3\}$ is not reachable,~$\textbf{J}$ is not sync-maximal by \cref{sync_maximal:thm:structure}.
	
		\item[\textbf{K}:] All the reachable $5$-sets are 
		$\{0,1,2,3,5\}$,$\{0,1,2,4,5\}$, and $\{0,1,2,3,4\}$.
		Thus, not all $5$-sets are reachalbe.
		Since, by \cref{sync_maximal:thm:structure}, for strongly connected automata,
		complete reachability is a necessary condition for sync-maximality, 
		$\textbf{K}$ is not sync-maximal.~\qedhere 	
	\end{enumerate}
\end{proof}

\begin{remark}
	In \textbf{J}, 
	the letters~$b$ and~$c$ restricted to $\{0,2,3\}$
	generate the full symmetric group on $\{0,2,3\}$.
	Since~$a$ maps $\{0,1,2\}$ to $\{1,2\}$,
	we can conclude that the subsets of the form $A \cup \{1,4\}$
	with $A \subseteq \{0,2,3\}$
	form a family of reachable and pairwise distinguishable subsets.
		
	In \textbf{K}, the letters $a$ and $bb$
	induce a subautomaton on $\{3,4,5\}$
	that is isomorphic to the \v{C}ern\'y automaton on three states
	and since $a$ and $b$ permute $\{0,1,2\}$
	the sets of the form $A \cup \{0,1,2\}$
	with $A \subseteq \{3,4,5\}$
	form a  family of reachable and pairwise distinguishable subsets.
\end{remark}

\section{Families of Slowly Synchronizing Automata}
\label{sync_max:sec:slowly_sync}

Here, we present a couple of slowly synchronizing families of automata from the literature and investigate 
the size of their minimal automata for the set of synchronizing words. 

\Cref{tab:slowly_sync_families} shows slowly synchronizing automata
which we show to be sync-maximal in \cref{sync_max:thm:families}.
We only give the state diagrams, the formal definitions should be clear from these diagrams.
For example, the automata family~$\mathcal  K_n = ([n], \Sigma, \delta_n)$
($n > 5$ odd) is defined\footnote{The numbering of the states w.r.t.~the letter~$a$	has been changed compared to ~\cite{DBLP:journals/ijfcs/Maslennikova19}.}
by
\begin{align*} 
	%\delta(i, b) & = i + 1 \mbox{ for } i \in \{0,\ldots, n-2\}, \mbox{ and } \delta(n-1, b) = 0, \\
	\delta_n(i, b) & = (i + 1) \bmod{n} \mbox{ for } i \in \{0,\ldots,n-1\}, \\
	\delta_n(i, a) & = (i + 1) \bmod{n} \mbox{ for } i \in \{1, \ldots,n-3,n-1\}, \\
	\delta_n(n-2,a) & = 1 \text{ and } 	\delta_n(0,a) = 3.
\end{align*}
For the families from \cref{tab:slowly_sync_families},  except for~$\mathcal K_n$ (for which
this was an open problem), sync-maximality has been shown
in~\cite{DBLP:journals/ijfcs/Maslennikova19} by rather involved arguments.\footnote{For $\mathcal V_n$
	sync-maximality follows readily by considering $a^{n-1}$ and $b$ as letters, as the resulting
	automaton is isomorphic to the \v{C}ern\'y automaton.}
By our general results this result follows readily. 

\input{../part3/slowlysync/tab_slowly_sync_aut}

% https://tex.stackexchange.com/questions/87975/one-figure-having-two-columns-of-subfigures-the-left-column-has-one-image-and-t
\newsavebox{\measurebox}
\begin{figure}[th]
	\centering
	\sbox{\measurebox}{%
		\begin{minipage}[b]{.45\textwidth}
			\subfloat
			[$\overline{\mathcal{M}}_n$ ($n$ even and $n \ge 8$) from~\cite{Martugin08}
			%with $\rt{\mathcal{M}_n} = \frac{n^2 + 6n - 16}{4}$ ($n$ even).
			with reset length $\frac{n^2 + 6n - 16}{4}$.
			]
			{\label{sync_max:fig:An}
		\scalebox{0.85}{
			\begin{tikzpicture}
			%	\tikzstyle{state}=[circle,draw,minimum size=30pt,inner sep=1pt]
				%\tikzstyle{state}=[blur shadow,fill={rgb:black,0.4;white,10},minimum size=30pt,inner sep=1pt]
				\tikzstyle{state}=[circle,draw,blur shadow,fill={rgb:black,0.4;white,10},minimum size=30pt,inner sep=1pt]				
				
				\node[state, xshift=-2.3cm] (0) {$0$};
				\node[state, right = of 0, xshift=-0.5cm] (1) {$1$};
				\node[state, right = of 1, xshift=-0.5cm] (2) {$2$};
				%\node[right = of 2] (dots) {$\ldots$};
				\node at (1.8,0) (dots) {$\ldots$};
				\node[state, right = of dots,xshift=-1cm] (m-2) {$m-2$}; 
				\node[state, right = of m-2, xshift=-0.4cm]  (m-1) {$m-1$};
				%\node[state, below right = of dots] (m) {$m$};
				%\node[state] at (8,-2) (m) {$m$};
				\node[state, below = of dots,xshift=-0.5cm,yshift=-0.3cm] (m) {$m$};
				\node[state, below right = of m] (2m-2) {$2m-2$};
				\node[state, below = of 2m-2] (2m-3) {$2m-3$};
				\node[state, below left = of 2m-3] (m+3) {$m+3$};
				\node[state, above left = of m+3] (m+2) {$m+2$};
				\node[state, below left = of m] (m+1) {$m+1$};
				\node[state, left = of m+1,xshift=0.5cm] (2m-1) {$2m-1$}; 
				
				\node at (2.35,-7.25) {$\reflectbox{$\ddots$}$};
				
				% blau b, a rot?
				
				\path[->] (0) edge [loop below] node {$a,b$} (0);
				\path[->] (1) edge [above] node {$a$} (0);
				\path[->] (2) edge [above] node {$a$} (1);
				\path[->] (m-1) edge [above] node {$a$} (m-2);
				\path[->] (1) edge [below] node {$b$} (m);
				\path[->] (2) edge [left] node {$b$} (m);
				\path[->] (m-2) edge [left] node {$b$} (m);
				%	\path[->] (m-1) edge [below,bend left] node {$b$} (m);
				\path[->] (m)   edge [below,bend right=40] node {$b$} (m-1);
				\path[->] (m-1) edge [above,pos=0.7,bend left=15] node {$b$} (m);
				
				\path[->] (2m-1) edge [loop above] node {$a$} (2m-1);
				\path[->] (m+1)  edge [above] node {$b$} (2m-1);
				\path[->] (2m-1) edge [below] node {$b$} (m+2);
				\path[->] (m+2)  edge [left] node {$a$} (m+1);
				\path[->] (m+2)  edge [above,bend left] node {$b$} (m+3);
				\path[->] (m+3)  edge [below,bend left] node {$a$} (m+2);
				\path[->] (2m-3) edge [left, bend left] node {$b$} (2m-2);
				\path[->] (2m-2) edge [right, bend left] node {$a$} (2m-3);
				\path[->] (2m-2) edge [above] node {$b$} (m+1);
				\path[->] (m) edge [right] node {$a$} (2m-2);
				\path[->] (m+1) edge [above] node {$a$} (m);
				
				%\node at (1.5,-5.5) {$\overline{\mathcal A}_n$ ($n$ even)};
			\end{tikzpicture}
		}	
				
		}	
	\end{minipage}}
	\usebox{\measurebox}\qquad
	\begin{minipage}[b]{.45\textwidth}
		\centering
		\subfloat
		[$\mathcal M_n$ ($n \ge 3$)
		with reset length $n^2 - 3n + 3$  introduced in~\cite{KisielewiczS15}.]
		{\label{sync_max:fig:Mn}
	\scalebox{.85}{\begin{tikzpicture}[shorten >=1pt,->,node distance=5cm]
			%\tikzstyle{vertex}=[circle,draw,minimum size=25pt,inner sep=1pt]
			\tikzstyle{vertex}=[circle,draw,blur shadow,fill={rgb:black,0.4;white,10},minimum size=25pt,inner sep=1pt]
			
			
			\foreach \name/\angle/\text in {P-n-2/234/n-2, P-n-1/162/n-1, 
				P-0/90/0, P-1/18/1, P-2/-54/2}
			\node[vertex] (\name) at (\angle:1.75cm) {$\text$};
			
			\path[->] (P-0)   edge [below, pos=0.3,bend left=10] node {$c$} (P-n-1)
			(P-n-1) edge [above, bend left=10] node {$c$} (P-0);
			
			\path[->] (P-1) edge [loop right]   node {$b,c$} (P-1);
			\path[->] (P-2) edge [loop right]   node {$b,c$} (P-2);
			\path[->] (P-n-2) edge [loop left] node {$b,c$} (P-n-2);
			\path[->] (P-n-1) edge [loop left] node {$b$} (P-n-1);
			
			\path[->] (P-n-1) edge [below] node {$a$} (P-1);
			
			\path[->] (P-0) edge [right,pos=.5] node {$a,b$} (P-1)
			(P-1) edge [right] node {$a$}   (P-2)
			(P-n-2) edge [left] node {$a$} (P-n-1);
			\node at (0, -1.5) {\large $\ldots$};			
	\end{tikzpicture}} 
}
		\vfill
		\subfloat
		[$\mathcal B_n$ ($n \ge 5$ odd) from~\cite{AnaVolGus2013,AnanichevVZ07}
		 with reset length $n^2 - 3n + 2$.		 
		%The slowly synchronizing automaton $\mathcal B_n$ ($n \ge 5$ odd) from~\cite{AnanichevVZ07,AnaVolGus2013}.
		%with a shortest synchronizing word of length $n^2 - 3n + 2$. 
		%As it has no letter of defect one it is not sync-maximal.
		]
		{\label{sync_max:fig:Bn}
					\scalebox{.85}{\begin{tikzpicture}[shorten >=1pt,->,node distance=5cm]
						%\tikzstyle{vertex}=[circle,draw,minimum size=10pt,inner sep=1pt]
						\tikzstyle{vertex}=[circle,draw,blur shadow,fill={rgb:black,0.4;white,10},minimum size=25pt,inner sep=1pt]
											
						\foreach \name/\angle/\text in {P-n-2/234/n-2, P-n-1/155/n-1, 
							P-0/90/0, P-1/25/1, P-2/-18/2, P-3/-60/3}
						\node[vertex] (\name) at (\angle:1.85cm) {$\text$};
						
						\path[->] (P-n-1) edge [above] node {$a$} (P-1)
						(P-0)   edge [pos=.8,left] node {$a$} (P-2)
						(P-1)   edge [loop right] node {$a$} (P-1)
						(P-2)   edge [loop right] node {$a$} (P-2)
						(P-3)   edge [loop right] node {$a$} (P-3)
						(P-n-2) edge [loop left]  node {$a$} (P-n-2);						
						
						\path[->] (P-n-2) edge [left] node {$b$} (P-n-1)
						(P-n-1) edge [above] node {$b$} (P-0)
						(P-0)   edge [above] node {$b$} (P-1)
						(P-1)   edge [right] node {$b$} (P-2)
						(P-2)   edge [right] node {$b$} (P-3);
						\node at (0, -1.55) {\large $\ldots$};
				\end{tikzpicture}} 
		}
	\end{minipage}
\caption[Non-sync-maximal and slowly synchronizing automata families]{Slowly synchronizing families
	whose automata are not sync-maximal, but minimal automata for their sets of synchronizing
	words are exponentially larger in the number of states.
	 See the main text for details. %$\mathcal B_n$ is not sync-maximal as it has no letter of defect one. 
} \label{def:overlineAn} \label{def:Mn} \label{def:Bn}
\label{fig:slowly}
\end{figure}

\begin{theorem}
\label{sync_max:thm:families}
	The automata from the families listed in Figure~\ref{tab:slowly_sync_families} are sync-maximal.
\end{theorem}
\begin{proof}
	Let $\semiautomat{A}{Q}{\{a,b\}}{\delta}$.
	For $\mathcal A \in \{\mathcal C_n, \mathcal V_n, \mathcal L_n \}$
	there exists a state~$q$ (marked in \cref{tab:slowly_sync_families})
	such that $\applytrans{\delta}{q}{a} = \applytrans{\delta}{q}{ba}$.
	As the letter $b$ cyclically permutes the states and~$a$ has defect one,
	these automata are sync-maximal by \cref{sync_max:thm:circular_sync_max} 
	(for $\mathcal C_n$, this was also established in \cref{sync_max:thm:cerny_sync_maximal} before).	
	For $\mathcal A \in \{\mathcal F_n, \mathcal K_n\}$ there exists a state~$q$
	such that $\applytrans{\delta}{q}{a} = \applytrans{\delta}{q}{b^2a}$.
	As in this case $n$ is odd, $a$ has defect one, and the letter $b$ cyclically permutes the states,
	sync-maximality follows as well from \cref{sync_max:thm:circular_sync_max}.
\end{proof}

\begin{theorem}
\label{syn_max:thm:An}
  The automata $\overline{\mathcal{M}}_n$ ($n$ even and $n \ge 8$) from~\cref{sync_max:fig:An} introduced in~\cite{Martugin08}\footnote{In~\cite{DBLP:journals/jalc/AnanichevV19} the notation
  	$\mathcal M_n$ was used for this family, that we already used
  	for the family from \cref{sync_max:fig:Mn}.
  	In~\cite{Martugin08} 
  	the notation $\mathcal A_n$ has been used, but we prefer to use
  	$\mathcal A_n$ to denote a general family of automata.}  	  
  are not sync-maximal
  but a minimal automaton for $\Syn(\mathcal{M}_n)$ has at least $2^{n/2}$ states.
 % sogar noch mehr bis auf konstante, wenn man oben die nicht am anfang "rausschiebt"
\end{theorem}
\begin{proof}
	Let $\overline{\mathcal{M}}_n = (Q_n,\{a,b\},\delta_n)$
	with $Q_n = [2m]$ where $m \ge 4$.	
	Since $0$ is a trap state, $\overline{\mathcal{M}}_n$ is not sync-maximal by \cref{syncmax:cor:three-states-trap-not-syncmax}.
	Set $P = \{ m, m +1, \ldots, 2m-1 \}$.
	In the following, we write~$\delta$ 
	for~$\delta_n$ for simplicity.
	Then
	\[
	( \fword{\delta}{babab} )_{| P } = (m ~ m + 1 ~ 2m-1 ~ m+2 ~ \cdots ~ 2m-2) \text{ and }
	( \fword{\delta}{a} )_{| P } = (2m - 2 ~ \cdots ~ m+1 ~ m).
	\]
	Set $u = ba^{m-1}$.
	Then we have $\delta(m, u) = 0$
	and, since $\delta(q, a^{m-1}) = q$
	for each $q \in P$,
	we have $\delta(q, u) = \delta(q, b)$
	for $q \in P \setminus \{m\}$.	
	
	\medskip 
	\noindent\underline{Claim 1:} Every subset of the form $S \cup \{0\}$
	with $S \subseteq P$
	is reachable.
	
	\medskip	
	\noindent\emph{Proof of Claim 1:} We do induction on $|P \setminus S|$.
	For $|P \setminus S| = 0$, we have $S = P$ and 
	$\delta(Q_n, u) = \{0\} \cup P$.	
	Next let $S \subsetneq P$
	be a proper subset and $q \in P \setminus S$.
    Set $T = \{m+1,\ldots,m+|S|\}$.
	Inductively, 
	we find $\delta(Q_n, v) = T \cup \{m\} \cup \{0\}$
	for some $v \in \{a,b\}^*$.
	This gives $\delta(Q_n, vu) = \delta(T, u) \cup \{0\}$.
	Since $T \subseteq \{m+1, \ldots, 2m-1\}$
	and~$b$ permutes the latter set
	and~$a$ permutes~$P$,
	we have $|\delta(T, u)| = |T| = |S|$.	
	By \cref{prelim:thm:cycles_gen_Sn},  
	there exists a word $w \in \{a,babab\}^*$
	inducing a permutation on~$P$
	mapping~$\delta(T, u)$
	onto~$S$. Then $\delta(Q_n, vuw) = S \cup \{0\}$.~\emph{[End, Proof of Claim~1]}
	
	\medskip 
	\noindent\underline{Claim 2:} Distinct subsets $S \cup \{0\}$, $T \cup \{0\}$
	with $S, T \subseteq P$ are distinguishable.
	
	\medskip 
	\noindent\emph{Proof of Claim 2:} 
	We do induction on $\min\{|S|, |T|\}$.
	If one set is empty, the other is non-empty
	and so the empty word distinguishes $S \cup \{0\}$ and $T \cup \{0\}$.
	Otherwise, suppose both sets are non-empty.	
	Without loss of generality, let $|S| \le |T|$.
	Pick some $q \in S$.
	Since $|S| \le |T|$, in both cases $q \in T$
	or $q \notin T$,
	we have $S \setminus \{q\} \ne T \setminus \{q\}$.
	Note that also $T \setminus \{q\} \ne \emptyset$.
	Let $v \in \{babab\}^*$ be such that $\delta(q, v) = m$.
	Set $S' = \delta(S, v)$
	and $T' = \delta(T, v)$.
	We have $m \in S'$,
	$\delta(S', u) = \delta(S' \setminus \{m\}, u) \cup \{0\}$,
	and $S' \setminus \{m\} \ne T' \setminus \{m\}$.
    Now, since $u$ acts on~$P \setminus \{m\}$ in the same way as~$b$, we find
    \[
    \delta(S' \setminus \{m\}, u) \ne \delta(T' \setminus \{m\}, u)
    \]
    and both subsets are contained in $P \setminus \{m\}$.
    Since $|\delta(S' \setminus \{m\}, u)| < |S|$ as $m \in S'$ and $|S'| = |S|$,
    by induction hypothesis,
    there exists $w \in \{a,b\}^*$
    distinguishing $\delta(S' \setminus \{m\}, u) \cup \{0\}$
    and $\delta(T' \setminus \{m\}, u) \cup \{0\}$.    
    We have  $\delta(T' \setminus \{m\}, u) \cup \{0\}
    = \delta(T', u) \cup \{0\}$
    since $\delta(m, u) = 0$.
    Now $vu$ maps $S \cup \{0\}$
    and $T \cup \{0\}$ to $\delta(S' \setminus \{m\}, u) \cup \{0\}$
    and $\delta(T', u) \cup \{0\}$, respectively,
    and so we find that $vuw$
    distinguishes $S$ and~$T$.~\emph{[End, Proof of Claim~2]}
    
	\medskip
	
	By Claims 1 \& 2, the $2^{n / 2}$ sets of the form $S \cup \{0\}$ for $S \subseteq P$
	are reachable and pairwise distinguishable. Then \cref{sync_max:thm:lower_bound_SynA} 
	gives the statement.
\end{proof}

\begin{theorem} \label{syn_max:thm:Mn}
    Consider the automata $\mathcal M_n$ from \cref{sync_max:fig:Mn} introduced in~\cite{KisielewiczS15}.
	The minimal automaton for $\Syn(\mathcal M_n)$ has precisely $2^{n-1} + 2^{n-2} - n + 1$ states.
\end{theorem}
\begin{proof} 
	Let $\semiautomat{M_n}{Q_n}{\{a,b,c\}}{\delta_n}$
	as in \cref{fig:slowly} where $Q_n = [n]$
	and let $S = \{1,\ldots,n-1\}$,
	$T = \{0,\ldots,n-2\}$,
	and $u = cb$ and $v = ac$. 
	When considering the words $\{a, u\}$
	and $\{v, b\}$ as letters on the state sets $S$
	and~$T$, respectively, we obtain
	automata isomorphic to the \v{C}ern\'y automaton $\mathcal C_{n-1}$.
	Since $\delta_n(Q_n, b) = S$
	and $\delta_n(Q_n, v) = T$, and
	the \v{C}ern\'y automaton
	is completely reachable and completely distinguishable 
	by \cref{sync_max:thm:cerny_compl_reach,sync_max:thm:cerny_compl_dist},
	we obtain that all non-empty subsets
	of~$S$ and~$T$ are reachable
	and pairs of distinct non-empty and non-singleton
	subsets of these sets are distinguishable.
		
	\underline{$Q_n$ is distinguishable from every other set:} Let $A \subseteq Q_n$.
	For $A \subseteq T$, 
	we have $\delta_n(A, b) \subseteq S \setminus \{n-1\}$
	and $\delta_n(Q_n, b) = S$,
	and so both subsets are distinguishable by a word $w \in \{a, u\}^*$.
	Then $bw$ distinguishes $A$ and $Q_n$.
	For $A \subseteq S$, we have $\delta_n(A, c) \subseteq T$
	and $\delta_n(Q_n, c) = Q_n$
	and as shown before there exists a word~$w$ distinguishing both
	subsets and so $cw$ distinguishes $A$
	and~$Q_n$.	
	
	\underline{Distinct $A \subseteq S$	and $B \subseteq T$ with $|A|, |B| \ge 2$ are distinguishable:}
	We have $\delta_n(B, b) \subseteq S$, $n - 1 \notin \delta_n(B, b)$,
	and $\delta_n(A, b) = A$.
	Thus, for $n - 1 \in A$,
    we have $\delta_n(A, b) \ne \delta_n(B, b)$
    and so since $|\delta_n(A,b)| \ge 2$ there exists a word in $\{a, u\}^*$
    distinguishing  $\delta_n(A, b)$ and $\delta_n(B, b)$.
    For $n - 1 \notin A$,
    we have $\delta_n(B, c) = ( B \setminus \{0\} ) \cup \{n - 1\}$
    and $\delta_n(A, c) = A$
    so that $\delta_n(B, c)$ and $\delta_n(A, c)$
    are distinct subsets of~$S$
    of size at least two.
    So some word in $\{a, u\}^*$
    distinguishes both subsets.
    
    In conclusion, we obtain that two non-empty and distinct subsets of states
	of size at least two which contain
	at most one state from $\{0,n-1\}$ or equal the whole state set
	are distinguishable and each such set is reachable.
		
	Now, for $A \subseteq Q_n$ with $0 \notin A$
	or $n - 1 \notin A$, inductively,
	we obtain $0 \notin \delta_n(A, w)$
	or $n - 1 \notin \delta_n(A, w)$ for all $w \in \Sigma^*$.
	Since $\delta_n(Q_n, a) = \delta_n(Q_n, b) = S$
	and $\delta_n(Q_n, c) = Q_n$,
	we find that for each $w \in \Sigma^*$
	with $\delta_n(Q_n, w) \ne Q_n$
	we have $0 \notin \delta_n(Q_n, w)$
	or $n - 1 \notin \delta_n(Q_n, w)$.
	Thus, no proper subset that contains $\{0,n-1\}$
	is reachable.
			
	Summarizing, we conclude that
	$\mathcal F = \{ A \subseteq Q : |A| \ge 2, |\{0,n-1\} \cap A| \le 2 \} \cup \{ \{0\} \} \cup \{Q_n\}$
	is a maximal family of pairwise distinguishable
	and reachable subsets of states.
	Since there are $2^{n - 1} - (n - 1) - 1$ subsets of size at least two
	of~$S$ and $2^{n - 2}$ subsets of~$T$
	that contain~$0$, we find that $\mathcal F$
	has size $2^{n-1} + 2^{n-2} - n + 1$.
	Then we get the claim from \cref{sync_max:thm:lower_bound_SynA}.
\end{proof}

\begin{theorem} \label{thm:Bn}
	Consider the family of automata given by $\mathcal B_n$ with $n \ge 5$ being odd from \cref{sync_max:fig:Bn}
	introduced in~\cite{AnaVolGus2013,AnanichevVZ07}.
	The automata $\mathcal B_n$ are not sync-maximal 
	and for~$n \equiv 0 \pmod{5}$, 
	the minimal automaton for $\Syn(\mathcal{B}_n)$ 
	has at least $2^{n / 5}$ states.
\end{theorem}
\begin{proof}
	Let $\mathcal{B}_n = (Q_n,\{a,b\},\delta_n)$ with $Q_n = [n]$.
	By \cref{sync_maximal:thm:structure},
	$\mathcal B$ is not sync-maximal since it has no letter of defect one.
	Suppose $n$ is odd and $n \equiv 0 \pmod{5}$.
	Set $R_r = \{ q \in Q \mid q \equiv r \pmod{5} \}$
	for $r \in \{0,1,2,3,4\}$.
	Then $bbbbb$ cyclically permutes~$R_r$ for $r \in \{0,1,2,3,4\}$,
	$\delta_n(R_1, a) = R_1$, $\delta_n(R_2, a) = R_2$, and $\delta_n(R_3, a) = R_3$.
	
	\medskip
	\noindent\underline{Claim 1:} A subset of the form
	$
	% S \cup \{ ( q + 1 ) \bmod n \mid q \in S \} \cup R_1 \cup R_2 \cup R_3
	S \cup  \delta_n(S, b) \cup R_1 \cup R_2 \cup R_3
	$
	with $S \subseteq R_4$ is reachable.
		
	\medskip 	
	\noindent\emph{Proof of Claim 1:} We do induction on $|R_4 \setminus S|$.
	For $S = R_4$,
	we have $Q_n = S \cup  \delta_n(S, b) \cup R_1 \cup R_2 \cup R_3$,
	which is reachable by the empty word.
	Next, suppose there exists $q \in R_4 \setminus S$.
	Then for $S' = \delta_n(S, b^{n - 1 - q})$
	we have $n - 1 \notin S'$ and $0 \notin \delta_n(S', b)$.
	Set $T = S' \cup \delta_n(S', b) \cup R_1 \cup R_2 \cup R_3$.
	Since $\{1,2\} \subseteq T$ we get
	\[
	\delta_n^{-1}(T, a) = \{n-1, 0\} \cup T = (S' \cup \{n-1\}) \cup \delta_n(S' \cup \{n-1\}, b) \cup R_1 \cup R_2 \cup R_3.
	\]
	Inductively, since $S' \cup \{n-1\} \subseteq R_4$,
	there exists $w \in \Sigma^*$ with $\delta_n(Q_n, w) = \{ n - 1, 0 \} \cup T$.
	Then $\delta_n(Q_n, wa) = T$ and
	$
	\delta_n(T, b^{q + 1})
	= S \cup \delta_n(S, b) \cup R_1 \cup R_2 \cup R_3.
	$~\emph{[End, Proof of Claim~1]}	
			
	\medskip
	\noindent\underline{Claim 2:} Every pair of distinct subsets of the form
	$
	S \cup  \delta_n(S, b) \cup R_1 \cup R_2 \cup R_3
	$
	with $S \subseteq R_4$ is distinguishable.
	
	\medskip 
	\noindent\emph{Proof of Claim 2:} Let $u \in \Sigma^*$ be a shortest synchronizing word. % u oben anders definiert
	Then we can write $u = av$ with $v \in \Sigma^*$
	where $\delta_n(Q_n \setminus \{n-1,0\}, v) = \{ p \}$
	and $|\delta_n(Q_n, v)| \ge 2$.
	Let $S, T \subseteq R_4$ be distinct
	and $S \ne \emptyset$.
	Pick some $i \ge 0$
	such that $n - 1 \in \delta_n(S, u^i)$.
	Set
	\begin{align*}
	R & = \delta_n(S, u^i) \cup \delta_n(S, u^ib) \cup R_1 \cup R_2 \cup R_3, \\
	K & = \delta_n(T, u^i) \cup \delta_n(T, u^ib) \cup R_1 \cup R_2 \cup R_3.
	\end{align*}
 	As $0 \in \delta_n(S, u^ib)$,
 	we have $\{n - 1, 0\} \subseteq R$, $R \setminus \{n-1,0\} \ne \emptyset$,
 	and $\{n - 1, 0\} \cap K = \emptyset$.
	Hence $\delta_n(K, v) = \{p\}$
	and $\delta_n(R, v) = \{\delta_n(n-1, v), \delta_n(0, v), p\} = \delta_n(Q_n, v)$.
	Then, $u^iv$  distinguishes 
	$
	S \cup \delta_n(S, b) \cup R_1 \cup R_2 \cup R_3
	$
	from 
	$
	T \cup \delta_n(T, b) \cup R_1 \cup R_2 \cup R_3.
	$
	\emph{[End, Proof of Claim 2]}
	
	\medskip
	
	Finally, by Claim 1 \& 2 and as $|R_4| = n / 5$
	we find that $2^{n/5}$ subsets are reachable and pairwise
	distinguishable. So,
	\cref{sync_max:thm:lower_bound_SynA} 
	gives the claim. 
\end{proof} 

\begin{remark}
	We remark that two distinct subsets $S, T$ 
	which are unions of sets from 
	\[
	\{ \{0,1\}, \{2,3\}, \ldots, \{n-3,n-2\}\}
	\]
	with $|S|, |T| \ge 3$ are distinguishable in $\mathcal B_n$.
	To show this, consider a shortest synchronizing
	word, which must have the form $av$ for some $v \in \Sigma^*$.
	By the minimal choice of~$av$, 
	with the notation from the proof of \cref{thm:Bn},
	we have $\delta_n(Q_n \setminus \{n - 1, 0\}, v) = \{p\}$ for some state~$p$
	and $|\delta_n(Q_n, v)| \ge 2$,
	and so at least one of both states $0$ or $n-1$
	is not mapped to~$p$.	
    Now suppose, without loss of generality,
    $\{q,q+1\} \subseteq S$
    and $\{q, q+1\} \cap T = \emptyset$
    for some state~$q$ with $q \equiv 0 \pmod{n}$.
    Then for $i = (n - 1) - q$ we have
    \[
    \{n-1, 0\} \subseteq \delta_n(S, b^i) \text{ and }
    \{n-1, 0\} \cap      \delta_n(T, b^i) = \emptyset.
    \]
    Thus $\delta_n(T, b^iv) = \{p\}$
    and, since $|\delta_n(S, b^i)| \ge 3$ 
    we have $\delta_n(S, b^i) \cap \{1,\ldots,n-2\} \ne \emptyset$
    and so $|\delta_n(S, b^iv)| \ge 2$.
\end{remark}

\begin{remark}
	Every non-empty set that is a union 
	of sets from the family
	\[
	\{ \{i,j\} \subseteq [n] \mid \text{$i \equiv 2x\pmod{n}$ and $j \equiv 2x - 1 \pmod{n}$ for some $x \ge 0$} \}
	\]
	is reachable in $\mathcal B_n$ (the given family equals
	the orbit of $\langle \fword{\delta}{bb} \rangle$
	on $\{0, n - 1\}$).
	Next, we show this inductively using the notation from the proof of \cref{thm:Bn}.
	It is clearly true for the whole state set.
	Consider a proper and non-empty subset $A$ of states
	which is a union of these sets.
	Then there exist states $i_1, i_2, j_1, j_2$
	with $i_2 \equiv i_1 - 1 \pmod{n}$
	and $j_2 \equiv j_1 - 1 \pmod{n}$
	such that $\{i_1, i_2\} \cap A = \emptyset$ and
	$\{j_1, j_2\} \subseteq A$
	with $j_1 \equiv i_1 + 2 \pmod{n}$
	and $j_2 \equiv i_2 + 2 \pmod{n}$.
	Since~$n$ is odd, the word $bb$ cyclically permutes the states.
	Consider the word $u = (bb)^{n - i_1} a (bb)^{i_1}$.
	We have 
	$\delta_n( \{i_1, i_2\}, u ) = \{ j_1, j_2 \}$
	and $\delta_n(i, u) = i$
	for all $i \in Q_n \setminus \{i_1, i_2\}$.
	Thus, $\delta_n^{-1}(A, u) = A \cup \{ i_1, i_2 \}$
	and $\delta_n^(A \cup \{i_1, i_2\}, u) = A$.
	Inductively, we have $\delta_n(Q_n, v) = A \cup \{i_1, i_2\}$
	for some $v \in \Sigma^*$.
	Hence, $vu$ maps the whole states set onto~$A$.	
\end{remark}

%\subsection{Two Conjectures About Slowly Synchronizing Automata}
\section{Discussion}
\label{sync_max:sec:slowly_sync_discussion}

\Cerny~\cite{Cer64} introduced the \Cerny family from \cref{intro:fig:cerny},
which forms a family of slowly synchronizing automata. Only the eight automata from \cref{sync_max:sec:cerny_bound} 
are known to have the same reset length, see also~\cite{DBLP:conf/lata/Volkov08}.
A little more than forty years later, in~\cite{AnanichevVZ07} the authors gave two slowly synchronizing families
of automata have with a letter of defect two, one of them was the family $\mathcal B_n$ given in \cref{sync_max:fig:Bn}.
In~\cite{Martugin08} the slowly synchronizing family of binary automata with a sink state
$\overline{\mathcal M}_n$ from \cref{sync_max:fig:An} was introduced. 
For this class of automata, an improved lower bound was given in~\cite{DBLP:journals/jalc/AnanichevV19} (based on a construction
from~\cite{Vorel2015_SynchronizationRoadColoringJumps}).
A construction for an arbitrary alphabet was given in~\cite{Pribavkina2011_SlowlySynchronizingZeroNoncomplete}.
A general method to obtain slowly synchronizing families of automata
based on labeling edges of primitive digraphs with alphabet symbols was investigated 
in~\cite{DBLP:conf/mfcs/AnanichevGV10,AnaVolGus2013}.
Slowly synchronizing automata families having three, four, and five letters
were given in~\cite{DonZB21}.
Volkov~\cite{DBLP:journals/jalc/Volkov19a} gave a slowly synchronizing family of $n$-state automata
all of whose letters are idempotents of rank $n / 2$.

The relation of slowly synchronizing automata and the size of automata for their sets
of synchronizing words was first systematically studied by Maslennikova~\cite{DBLP:journals/corr/Maslennikova14,DBLP:journals/ijfcs/Maslennikova19}.
She proved \cref{sync_max:thm:families}, except for $\mathcal K_n$, which was left
open as a conjecture in~\cite{DBLP:journals/ijfcs/Maslennikova19} (which we confirmed here,
a result published in~\cite{HoffmannCircularExtended}).
We note that the proofs given in~\cite{DBLP:journals/corr/Maslennikova14,DBLP:journals/ijfcs/Maslennikova19}
were all quite involved, whereas here this result followed rather easily
from more general statements we gave.

\IN{almost sync-maximal}
We proved that many families of slowly synchronizing automata
have exponentially larger automata for their sets of synchronizing words.
We call an infinite family of distinct automata $\semiautomat{A_n}{Q_n}{\Sigma}{\delta_n}$
\emph{almost sync-maximal}, if there exists $C > 0$
such that the minimal automata for $\Syn(\mathcal A_n)$
have at least~$2^{C \cdot |Q_n|}$ states. %, i.e., the growth of these automata is not in $2^{o(|Q_n|)}$.
Similarly, we call an individual $n$-state automaton \emph{almost sync-maximal} if the minimal automaton
has size close to $2^n - n$ (this being an informal definition similarly to the definition for individual
slowly synchronizing automata).

Note that since the individual automata are over a fixed alphabet,
an almost sync-maximal infinite family of distinct automata
never consists solely of commutative automata by \cref{prop:aut_for_sync_words}.

The evidence suggests % a relation between
that being slowly synchronizing is related to 
being almost sync-maximal. 
%However, this connection is yet to be uncovered and seems to be subtle.
As a shortest synchronizing
word is a path in an accepting automaton for the set of synchronizing words,
when we need long synchronizing words, %this might be linked
%to accepting automata for the sets of synchronizing words
we might expect accepting automata for the sets of synchronizing words
to be ``complex'' or ``large''. %, leading us to our first conjecture.
These consideration lead us to our first conjecture.
%A shortest synchronizing
%word is a path in an accepting automaton for the set of synchronizing words.
%So, when we need long synchronizing words, %this might be linked
%%to accepting automata for the sets of synchronizing words
%we might expect accepting automata for the sets of synchronizing words
%to be ``complex'' or ``large'', leading us to our first conjecture.

\begin{center}
	\fbox{ \begin{minipage}{30em}
			Conjecture: Every slowly synchronizing family of automata 
			is almost sync-maximal.
		\end{minipage}
	}
\end{center}

However, the converse does not hold. \Cref{syncmax:discussion:item:alphabet,syncmax:discussion:item:inside}
show how to obtain sync-maximal automata with short synchronizing words.

\begin{example}
	\label{syncmax:discussion:item:alphabet} Adding letters to a sync-maximal or almost sync-maximal automaton retains these properties,
	but may alter the reset length. For example, adding a letter that maps everything to a single 
	state yields a reset length of one. 
\end{example} 

\begin{example}
	\label{syncmax:discussion:item:inside} 
	Consider the family of 
	binary $n$-state automata from \cref{fig:rev_cerny}.
	The automata from this family are sync-maximal since by
	considering the word $u = b^{n-1}$ as a new letter and $\{u,a\}$ as an alphabet, 
	the resulting automata over $\{u,a\}$
	are isomorphic to the \v{C}ern\'y automata~$\mathscr C_n$
	and induce the same transformation monoids as $b$ is a permutational letter.
	However, this family is not slowly synchronizing
	since $(ab)^n$ is a synchronizing word.
	Note that removal of any letter results in a non-synchronizing automaton for $n \ge 3$.
	
	\begin{figure}[ht]
		\centering
		\scalebox{1.0}{\begin{tikzpicture}[shorten >=1pt,->,node distance=4cm]
				%\tikzstyle{vertex}=[circle,color=red!60,fill=black!10,minimum size=25pt,inner sep=1pt,very thick]
				%\tikzstyle{vertex}=[fill=white, circle, draw=black, scale=1,minimum size=25pt,inner sep=1pt]
				\tikzstyle{vertex}=[circle,draw, blur shadow,fill={rgb:black,0.4;white,10},minimum size=30pt,inner sep=1pt]
				
				\foreach \name/\angle/\text in {P-n-2/234/n-2, P-n-1/162/n-1, 
					P-0/90/0, P-1/18/1, P-2/-54/2}
				{
					\node[vertex] (\name) at (\angle:1.75cm) {$\text$};
				}
				% \node at (0, -1.5) {$\ldots$};
				
				\path[->] (P-n-2) edge [-latex, left] node {$b$} (P-n-1);
				\path[->] (P-n-1) edge [-latex, above,pos=0.35] node {$b$} (P-0);
				\path[->] (P-0) edge [-latex, above] node {$b$} (P-1);
				\path[->] (P-1) edge [-latex, right] node {$b$} (P-2);
				
				\path[->] (P-n-2) edge [-latex,out = 215, in=145, distance=1cm, left] node {$a$} (P-n-2);
				\path[->] (P-n-1) edge [-latex,out = 215, in=145, distance=1cm, left] node {$a$} (P-n-1);
				\path[->] (P-1) edge [-latex,loop right, out = 325, in=35, distance=1cm] node {$a$} (P-1);
				\path[->] (P-2) edge [-latex,loop right, out = 325, in=35, distance=1cm] node {$a$} (P-2);
				\path[->] (P-0) edge [-latex, out = 260, in = 0, below, pos=.4] node {$a$} (P-n-1);
				
				\node at (0, -1.5) {\large $\ldots$};
		\end{tikzpicture}}%
		
		\caption[Binary sync-maximal automaton with short reset length]{Sync-maximal $n$-state automata with synchronizing words $(ab)^n$, see \cref{syncmax:discussion:item:inside}.}
		\label{fig:rev_cerny}
	\end{figure}
\end{example} 

Now, for two automata $\mathcal A$ and $\mathcal B$
with $\TM{\mathcal A} = \TM{\mathcal B}$
the state complexities of $\Syn(\mathcal A)$ and $\Syn(\mathcal B)$
are equal as the properties of reachability of subsets of states 
and distinguishability of pairs of subsets of states only
depend on the transformation monoid.
For example, adding a letter synchronizing the automaton in \cref{syncmax:discussion:item:alphabet} 
does not change the transformation monoid,
or considering $\{u,a\}$ as an alphabet in \cref{syncmax:discussion:item:alphabet}.
In view of these observations, 
we conjecture
that a slowly synchronizing family ``hides inside'' an almost sync-maximal family.

\begin{center}
	\fbox{ \begin{minipage}{30em}
			Conjecture: For every almost sync-maximal $n$-state family of automata~$\mathcal A_n$
			there exists a slowly synchronizing family of automata~$\mathcal B_n$ with state sets of size at most~$n$
			such that $\TM{\mathcal B_n}$ is isomorphic to a subsemigroup of $\TM{\mathcal A_n}$ for all $n \ge 1$?
			%where $\TM{\mathcal B_n} = \TM{\mathcal A_n}$ for all $n \ge 1$?
		\end{minipage}
	}
\end{center}

By \cref{sync_max:cor:structure_letters}, sync-maximal automata have at least two letters and binary
sync-maximal automata are circular and have a letter of defect one. %and are letter-reduced sync-maximal.
So, as the simplest class of automata containing sync-maximal automata, 
it is natural to first try to confirm both conjectures from the present section
for binary circular automata having a letter of defect one. 

The following statement confirms the above conjecture for families of binary sync-maximal automata,
where we even have an isomorphism to the whole transformation monoid of $\mathcal A_n$
if one letter is a simple idempotent.

\begin{proposition} \label{prop:binary-sync}
	Let $\mathcal A_n$ be a binary sync-maximal automaton $\mathcal A_n$ with state set~$Q = [n]$ where $n \ge 3$.
	
	\begin{enumerate}
		\item If one letter is a simple idempotent,
		then $\TM{\mathcal A_n}$ is isomorphic to $\TM{\mathcal C_n}$
		where $\mathcal C_n$ denotes the $n$-state automaton from the \v{C}ern\'y family.
		
		\item $\TM{\mathcal A_n}$ contains a subsemigroup isomorphic
		to $\TM{\mathcal C_m}$ for some divisor $m$ of~$n$ with $m > 1$.
		
		\item If~$n$ is prime, then $\TM{\mathcal C_n}$ is isomorphic to a subsemigroup of $\TM{\mathcal A_n}$.
	\end{enumerate}
\end{proposition}
\begin{proof}
	Since $\mathcal A_n$ is sync-maximal, by \cref{sync_max:cor:structure_letters}
	we have $\Sigma = \{a,b\}$ where~$a$ has defect one
	and~$b$ cyclically permutes the states.
	First, assume~$a$ is a simple idempotent
	mapping $q$ to a distinct state $\delta(q, b^i)$ for some $i \in \{1,\ldots,n-1\}$.
	Let $u = b^i$. The word $u$ cyclically permutes the states
	and the automaton we obtain when considering $\{a,u\}$
	as letters equals, up to renaming of the states and letters,~$\mathcal C_n$ (and so has a transformation monoid isomorphic to~$\TM{\mathcal C_n}$)
	and, since~$b$ permutes the states, has the same transformation monoid as~$\mathcal A_n$.
		
	Next, consider the general case. By \cref{sync_words:lem:idempotent} 
	there exists a word $v$ inducing a simple idempotent on the states
	that maps a state $q$ to a distinct state $\delta(q, b^j)$ for some $j \in \{1,\ldots,n-1\}$.
	Set $w = b^j$. The word~$w$ cyclically permutes the states~$S \subseteq Q$
	that are multiples of~$j$. So, $|S|$ divides~$|Q|$ and we have $|S| > 1$.
	The automaton resulting when considering the words $\{w, v\}$
	as letters on the state set~$S$
	equals, up to renaming of the states and letters, $\mathcal C_m$
	for $m = |S|$.
	Thus the transformation semigroup of~$\mathcal C_m$
	can be identified with the subsemigroup of~$\TM{\mathcal A_n}$
	generated by~$w$ and~$v$.
		
	Now, if~$n$ is prime, we have $S = Q$ in the previous argument, 
	so that $m = n$ and $\mathcal C_n$
	also acts on the same state set~$Q$.
	Note that $\TM{\mathcal C_n}$
	does not has to be a subset of $\TM{\mathcal A_n}$
	since the states in the automaton over $\{w,v\}$
	are numbered according to~$w$ and not~$b$,
	but is only isomorphic to a subsemigroup.
\end{proof}

The second statement from \cref{prop:binary-sync} for non-prime~$n$
can be seen in the automaton from \cref{ex:I_vs_Gamma1} where we find
a subsemigroup in its transformation monoid isomorphic to $\TM{\mathcal C_2}$.

We obtain a stronger variant of the last conjecture if we demand that the 
automata~$\mathcal B_n$ also have $n$ states.

Next, we give a family of binary automata where the state complexities
of the sets of synchronizing words have superpolynomial growth
and the reset lengths of all automata having the same transformation monoids
are linear. So, the above conjecture is false 
when we only demand superpolynomial growth.

\begin{example}\label{syncmax:slowly:ex} 
	Starting from the  \v{C}ern\'y automata  $\semiautomat{C_n}{[n]}{\{a,b\}}{\delta_n}$ for $n \ge 3$,
	here, we construct a family of $m$-state automata~$\semiautomat{A_m}{[m]}{\{a,b\}}{\mu_m}$ with $m = (n-1)^2$
	where the state complexity of $\Syn(\mathcal A_m)$
	is $2^{\sqrt{m}+1} - \sqrt{m} - 1$ and such that for every family of automata~$\mathcal B_m$
	with $\TM{\mathcal B_m} = \TM{\mathcal A_m}$
	we have $\rt{\mathcal B_m} \le m$.
	For $q \in [m]$ and $x \in \{a,b\}$ we set
	\[
	\mu_m(q, x) =
	\begin{cases}
		\delta_n(q, x) & \text{if } q \in [n], \\
		1              & \text{otherwise.}		
	\end{cases}
	\]
	We have $\Syn(\mathcal A_m) = \Syn(\mathcal C_n)$.
	So, a minimal automaton for $\Syn(\mathcal A_m)$
	has size $2^{\sqrt{m} + 1} - (\sqrt{m} + 1)$
	and the reset length of $\mathcal A_m$
	is $m$.
	Now, let $\semiautomat{B}{[m]}{\Gamma}{\rho}$
	be with $\TM{\mathcal B} = \TM{\mathcal A_m}$.
    Set $S = [n]$. Then
    \begin{enumerate}
    \item $S$ is an upper set in  $\mathcal A_m$
    and in $\mathcal B$,
    
    \item $t(q) \in S$ for every $t \in \TS{\mathcal A_m}$,
    
    \item $\mathcal A_{|S} = \mathcal C_n$,
    
    \item $\TM{\mathcal B_{|S}} = \TM{( \mathcal A_m )_{|S}}$.    
    \end{enumerate}
    The group of units of $\TM{\mathcal C_n}$
    is generated by a permutation cyclically permuting~$S$.
    As $\TM{\mathcal B_{|S}} = \TM{( \mathcal A_m )_{|S}}$,
    we find that $\mathcal B_{|S}$ must be a circular
    automaton.
    So, the reset length of $\mathcal B_{|S}$
    is at most $(n-1)^2$
    as the \v{C}ern\'y conjecture
    has been confirmed for circular automata~\cite{Dubuc98}.		
    Hence, the reset length of~$\mathcal B$
    is at most~$m + 1$.	
\end{example}

\partwithquote{The Constrained Synchronization Problem}{``Du glaubst nicht, wie müde ich bin, wie satt ich die Chemie habe, wie namentlich gewisse Teile mich ordentlich anekeln, mir wenigstens so langweilig sind, dass ich gähnen muss, wenn ich daran denke. Sind wir denn schon so alt, oder wie ist das? Diese nervenschwächende Wirkung muss wirklich der Chemie eigentümlich sein. Ich glaube, die materiellen Influenzen, die Dämpfe, Gerüche und all die Teufelsstinkereien haben großen Anteil daran. Besonders ist es das Praktikum, was einen so herunterbringt.''}{--Friedrich Wöhler (1800--1882) \\ Aus einem Brief an Justus von Liebig (22.07.1847)}
\label{part:constrsync}


\chapter{Motivational Examples}
\label{constr_sync:sec:motivation}


Vladimir Gusev~\cite{Gusev2012} introduced the constrained synchronization problem
in connection with a reduction procedure on automata
that preserves the reset length up to plus/minus one w.r.t.~a given constraint language. 
We discuss this in \cref{constrsync:sec:motivation:reduction}.
The problem of constrained synchronization also appears implicitly 
in a reduction used by Martyugin~\cite{Martyugin14}. %, for which we give details below
Here, we give details on the procedure from~\cite{Gusev2012} and on
the reduction used in~\cite{Martyugin14} in \cref{constrsync:sec:motivation:implicit}.
Then we outline a relation to the \emph{Skolem problem}~\cite{BLONDEL2002,DBLP:conf/rp/OuaknineW12} and the \emph{matrix orbit problem}~\cite{KannanLipton1986} in \cref{constrsync:sec:motivation:skolem}.
Finally, we sketch possible applications that naturally extend
areas where synchronizing automata have been used previously
in \cref{constrsync:sec:motivation:operating,constrsync:sec:motivation:part-orienters,constrsync:sec:motivation:supervisory}.

\section{Use in a Reduction Procedure}
\label{constrsync:sec:motivation:reduction}

\IN{$r$-decomposition}
\IN{rank factorization}
Let $\semiautomat{A}{Q}{\Sigma}{\delta}$ be a semiautomaton with $n$ states
over a binary alphabet $\Sigma = \{a,b\}$. A set of matrices $X_1, X_2, Y_1, Y_2$
over $\{0,1\}$ %Boolean semiring
is an \emph{$r$-decomposition}~\cite{Gusev2012} of $\mathcal A$
if
\begin{enumerate}
\item $X_1$ and $Y_1$ are $n \times r$ matrices,
\item $X_2$ and $Y_2$ are $r \times n$ matrices,
\item $[a] = X_1 X_2$, $[b] = Y_1 Y_2$,
\item ev ery row of $X_1, X_2, Y_1, Y_2$ has only one occurrence of~$1$.
\end{enumerate}
The automaton $\mathcal A$ admits an $r$-decomposition if and only if $r \ge \max\{\rank{a},\rank{b}\}$~\cite{Gusev2012}.
By the first three conditions $X_1, Y_1$ and $X_2,Y_2$ constitute
\emph{rank factorizations}
of $[a]$ and $[b]$, respectively, as discussed in~\cite{piziak99}.
Since the matrices $X_1, X_2, Y_1, Y_1$ have one~$1$ in each row,
the same holds for their products.
Hence, we obtain a complete and deterministic $r$-state automaton
$\semiautomat{A_{X_1, X_2, Y_1, Y_2}}{Q_{X_1, X_2, Y_1, Y_2}}{\Gamma}{\delta_{X_1,X_2, Y_1, Y_2}}$
by setting $Q_{X_1, X_2, Y_1, Y_2} = \{1,\ldots,r\}$ and $\Gamma = \{ a_{X_2Y_1}, a_{X_2X_1}, a_{Y_2X_1}, a_{Y_2Y_1} \}$
where the action of the letters is given
by $[a_{X_2Y_1}] = X_2 Y_1$, $[a_{X_2X_1}] = X_2X_1$, $[a_{Y_2X_1}] = Y_2 X_1$
and $[a_{Y_2 Y_1}] = Y_2 Y_1$. % (w.r.t.~some fixed ordering of the states $Q_{X_1, X_2, Y_1, Y_2}$, see \cref{prelim:subsec:transmtx}).
The following was shown by Gusev~\cite[Prop.~5]{Gusev2012}.

\begin{proposition}
		 Let $\semiautomat{A}{Q}{\{a,b\}}{\delta}$ be an automaton with $r$-decomposition $X_1, X_2, Y_1, Y_2$
		 and $\semiautomat{A_{X_1, X_2, Y_1, Y_2}}{Q_{X_1, X_2, Y_1, Y_2}}{\Gamma}{\delta_{X_1,X_2, Y_1, Y_2}}$.
		 Then $\mathcal A$ is synchronizing if and only if 
		 \[ Syn(\mathcal A_{X_1, X_2, Y_1, Y_2}) \cap R \ne \emptyset
		 \] 
		 with
		 \begin{align*}
		 	R = \Gamma^* \setminus \big(
		 	& \Gamma^* ( a_{X_2 X_1} ( a_{Y_2 X_1} + a_{Y_2 Y_1} ) ) \Gamma^* + 
		 	\Gamma^* ( a_{X_2 Y_1} ( a_{X_2 Y_1} + a_{X_2 X_1} ) ) \Gamma^* + \\ 
		 	& \Gamma^* ( a_{Y_2 Y_1} ( a_{X_2 Y_1} + a_{X_2 X_1} ) ) \Gamma^* + 
		 	\Gamma^* ( a_{Y_2 X_1} ( a_{Y_2 X_1} + a_{Y_2 Y_1} ) ) \Gamma^* \big),
		 \end{align*}
	     where, if this is the case, $m - 1 \le \rt{\mathcal A} \le m + 1$
	     for  $m = \min\{ |u| : u \in \Syn(\mathcal A_{X_1, X_2, Y_1, Y_2}) \cap R  \}$.
\end{proposition}

So, the reset length of an automaton $\mathcal A$ without permutational letters can be bounded by the length of a shortest synchronizing word 
of a smaller automaton $\mathcal A_{X_1, X_2, Y_1, Y_2}$ contained in a constraint language independent of $\mathcal A$.

Gusev~\cite{Gusev2012} used $r$-decompositions to construct $n$-state automata 
with reset length at least $r^2 - r - 1$. This was the sole purpose of this work
and while the above statement used constrained synchronization, it was not introduced
as a decision problem and questions of computational complexity were not 
investigated.

\section{Implicit in a Reduction}
\label{constrsync:sec:motivation:implicit}

Martyugin~\cite{Martyugin14} 
proved that it is \PSPACE-complete to decide whether a given partial automaton
has a carefully synchronizing word (for the definition of these notions, see \cref{constr_sync:subsec:pdfa}).
His reduction constructs a partial automaton over $\Sigma$
where a shortest carefully synchronizing 
word has the form $x\Gamma^*y$ for $x,y \in \Sigma$ and $\Gamma = \Sigma \setminus \{x,y\}$~\cite[Lem.~1]{Martyugin14}.
By adding self-loops to obtain 
a complete automaton in the construction,
one can easily see that it has a synchronizing word in $x\Gamma^*y$
if and only if the original partial automaton
has a carefully synchronizing word.
In this sense, \constrsync{x\Gamma^*y} is implicit
in the construction. 

\section{The Matrix Orbit Problem and the Skolem Problem}
\label{constrsync:sec:motivation:skolem}

\IN{Skolem problem} 
\IN{Pisot's problem} 
\IN{matrix orbit problem}
The \emph{(matrix) orbit problem} is the problem
which asks for an $n \times n$ matrix $A$ over $\mathbb Q$
and
$x,y \in \mathbb Q^n$ %a matrix $A \in \mathbb Q^{n \times n}$
whether $xA^i = y$
for some $i \ge 0$. This problem was shown to be solvable
in polynomial time by Kannan \& Lipton~\cite{KannanLipton1986}.
In~\cite[Cor.~1.3]{BLONDEL2002}
the problem which asks for an $n \times n$ matrix $A$ over~$\mathbb Z$
and $1 \times n$ and $n \times 1$ vectors $x$ and $y$ over $\mathbb Z$, respectively, 
whether $xA^iy = 0$ for some $i \ge 0$ was shown to be \NP-hard.
As noted in~\cite{BLONDEL2002},
this implies \NP-hardness of the \emph{Skolem problem} (also called \emph{Pisot's problem})
which ask whether an integer linear recurrent sequence has zero as a value,
see \cite{BLONDEL2002}, and also~\cite{DBLP:conf/rp/OuaknineW12} for more information on the Skolem problem.
Decidability of the Skolem problem is unknown~\cite{DBLP:conf/rp/OuaknineW12}.

Next, we point out a connection
to the two problems above.\footnote{These connections were originally pointed out by Vladimir V. Gusev (personal communication).}
Let $\semiautomat{A}{Q}{\{a,b\}}{\delta}$ with $Q = \{q_1,\ldots,q_n\}$
and $q \in Q$.
Set
\[
V_q = \{ (y_1, \ldots, y_n) \mid y_j > 0 \text{ implies } q_j \in \delta^{-1}(q, a) \text{ for } j \in \{1,\ldots,n\} \},
\]
and define numbers $x_1, \ldots, x_n$ and $z_1, \ldots, z_n$
by 
\[
x_i = \begin{cases}
	1 & \text{if } q_i \in \delta(Q, a), \\
	0 & \text{otherwise,}
\end{cases}
\quad\text{ and }\quad
z_i = \begin{cases}
	1 & \text{if } \delta(q_i, a) \ne q, \\
	0 & \text{otherwise,}
\end{cases}
\]
for $i \in \{1,\ldots,n\}$. Then for $i \ge 0$ we have
\[
	\delta(Q, ab^ia) = \{q\}
	 \Leftrightarrow x \cdot [b]_{\mathcal A}^i \in V_q 
	 \Leftrightarrow  x \cdot [b]_{\mathcal A}^i \cdot \begin{pmatrix} z_1 \\ \vdots \\ z_n \end{pmatrix} = 0
\]
We will show later in \crefitem{constr_sync:thm:introductory_examples}{constrsync:ex:item:abstara}
that $\constrsync{ab^*a}$ is \NP-complete. The above construction
then also yields \NP-hardness of the problem introduced above
that asks whether $xA^iy = 0$ for a matrix~$A$ and vectors~$x,y$ over~$\mathbb Q$.
The construction also yields another proof
that $\constrsync{ab^*a}$ is in \NP\  since checking $x \cdot [b]_{\mathcal A}^i \in V_q$ for some $q \in Q$
is in \NP\  by guessing $q \in Q$ and then checking $y \in V_q$ and $x[b]^i_{\mathcal A} = y$,
which can be done in polynomial time since the matrix orbit problem 
can be solved in polynomial time.
 
\begin{remark}
	In \cref{prelim:compl:thm:reach-comm-fixed} 
	we proved that \REACHF
	is solvable in polynomial time for commutative automata over a fixed alphabet.
	The result of Kannan \& Lipton~\cite{KannanLipton1986}
	stating polynomial-time solvability 
	of existence of some $i \ge 0$
	such that $xA^i = y$ for a matrix~$A$
	and vectors $x,y$ over~$\mathbb Q$
	also implies polynomial-time solvability
	of \REACHF for unary automata
	similarly as above by considering the transition matrix and encoding state
	sets $S$ into vectors $(x_1, \ldots, x_n)$ 
	by setting $x_i = 1$ if $q_i \in S$ and $x_i = 0$ otherwise.
	Note that the related problem \ST 
	is \NP-complete for unary automata by \cref{compl:thm:comm_ST}.	
\end{remark}

\section{Operating Restrictions}%ursprungsbeispiel, aber auch homogene experimente sternberg paper
\label{constrsync:sec:motivation:operating}

Fernau et al.~\cite{FernauGHHVW19} gave the following motivating example. Suppose that a system, modeled by an automaton, 
has a ``normal'' and a ``directing`` operating mode, and only in the latter mode
synchronizing instructions are allowed. The switching between both modes 
can be modeled, in its simplest form, 
by the constraint $ab^*a$, where $a$ switches into the ``directing'' mode.
As written in the previous paragraph, \constrsync{ab^*a} is \NP-complete, see \crefitem{constr_sync:thm:introductory_examples}{constrsync:ex:item:abstara}.
Even more generally, it might be possible that a system, for example, a remotely controlled
rover on a distant planet or a satellite in orbit, is not allowed to execute all commands in every possible
order. For example, some commands are only allowed in a certain order or after
other commands have been executed. All of this imposes constraints
on the possible reset sequences.

% bestimmte operationen nur wenn bestimmte andere, oder erst vertikal
\section{Part Orienters}
\label{constrsync:sec:motivation:part-orienters}

Generalizing the part orienters example
from \cref{intro:sec:examples:partorienters}, suppose some parts arrive at a manufacturing site and they
need to be sorted and oriented before assembly. Practical considerations
favor methods which require little or no sensing, employ simple devices,
and are as robust as possible. 
As in \cref{intro:sec:examples:partorienters},
the parts are put on a conveyor belt
so that the stream of the parts encounter a series of passive obstacles
placed along the belt before reaching the assembly point.
Now, furthermore, assume, due to the availability of components, requirements how to lay things out
or physical restrictions,
the passive components can %could 
not be placed at random along
the belt, but have to obey some restrictions, or restrictions
in what order they are allowed to happen. 
These constraints on the
permissible sequences of applicable operations 
can be modeled by the constrained synchronization problem.


\IN{discrete event system}
\section{Supervisory Control} 
\label{constrsync:sec:motivation:supervisory}

Ramadge and Wonham~\cite{RamadgeWonham87} modelled
 a \emph{discrete event system} (DES) by a partial automaton 
 generating an event sequence (see also \cite{DBLP:books/daglib/0034521,wonham2019}). 
 %The authors distinguished between controllable and uncontrollable
 %events.
 Then, a supervisor,  depending on the event sequence \todo[inline]{supervisor or controller}
 that occurred so far,
 is able to restrict the set of events
 that are possible in the next step, where, however,
 he can only limit the use of so-called controllable events.
 So, if we want to (globally) reset a finite state DES~\cite{Alves2020} under supervisory control, 
 this can be seen as an instance of a constrained synchronization problem
 in case the permissible event sequences form
 a regular language.
 
 

% % noch etwas umschreiben, sonst identsich mit Mischa-Text
% \item[Biocomputing.] In~\cite{Benenson2003,Benenson2001} DNA molecules have been used as both
% hardware and software for finite automata of nanoscaling size, see also~\cite{Vol2008}. 
% For instance, Benenson et al~\cite{Benenson2003} produced ``a `soup of automata', that is, a solution
% containing $3 \times 10^{12}$ identical automata per $\mu1$. All these molecular
% automata can work in parallel on different inputs, thus ending up in different
% and unpredictable states. In contrast to an electronic computer, one cannot
% reset such a system by just pressing a button; instead, in order to synchronously bring
% each automaton to its start state, one should spice the soup with (sufficiently
% many copies of) a DNA molecule whose nucleotide sequences encodes a reset word''~\cite{Vol2008}.
% Now, it might be possible that certain sequences, or subsequences,
% are not possible as they might have unwanted biological side-effects, 
% or might destroy the molecules at all.


\chapter{Further Notions and Preliminaries}
\label{constr_sync:chap:further_notions}

%!TEX root = ./../../springer-mono-format/book.tex
%!TEX root = ./../../thesis-trier.tex

% https://english.stackexchange.com/questions/186434/is-notations-a-proper-english-word
Here, we introduce additional concepts and notations %that is only 
needed in \cref{part:constrsync}, and give two new results in coding theory.
Specifically, in \cref{constr_sync:subsec:pfun}, we introduce partial functions. 
\Cref{constr_sync:subsec:pdfa} introduces
partial automata (PDFAs), which generalize our automata model by allowing
the transition relation to be partial function.
These automata give a convenient specification mechanism
for constraint languages (see the discussion in \cref{constrsync:notions:sec:discussion}).
\Cref{constr_sync:subsec:levenshtein-romanov-map} introduces
the Levenshtein-Romanov mapping.
In \cref{constr_sync:chap:classes_constraint_languages} we use prefix codes in some hardness conditions
and solid codes for a statement on the preservation of the computational complexity.
\Cref{constr_sync:subsec:codes} introduces the relevant notions
related to codes, various subclasses of codes as, for example, solid codes,
and overlap-free languages. We also present two new characterizations of solid codes.
 % that have not been noted before.
Later, in \cref{thm:constr_sync_hom_strongly_self_sync}, 
we will prove, using the Levenshtein-Romanov mapping, 
that the computational complexity
of the constrained synchronization is preserved for reencodings
of the constraint language by solid codes.
\Cref{constr_sync:sec:bounded,constr_sync:subsec:monoid} introduce bounded and monoid languages, respectively.
%both will be later investigated 
%as classes of constraint languages in \cref{sec:constr_sync:letter_bounded,constr_sync:sec:monoid}, respectively.
Lastly, in \cref{constr_sync:sec:growth}, we investigate
the growth of languages, and state connections to bounded and monoid languages.
In particular, we classify the growth of monoid languages.

\section{Partial Functions}
\label{constr_sync:subsec:pfun}

\input{../part4/notions/partialfunc}

\section{Partial Automata (PDFAs)}
\label{constr_sync:subsec:pdfa}

\input{../part4/notions/pdfa}

\section{The Levenshtein-Romanov Mapping}
\label{constr_sync:subsec:levenshtein-romanov-map}

\input{../part4/notions/levenshtein}

\section{Codes and Overlap-Free Languages}
\label{constr_sync:subsec:codes}

\input{../part4/notions/codes}

%\section{Characterizing Solid Codes}
%\label{constr_sync:subsec:charsolidcode}

\section{Bounded Languages}
\label{constr_sync:sec:bounded} 

\input{../part4/notions/bounded}

\section{Monoid Languages}
\label{constr_sync:subsec:monoid}

\input{../part4/notions/monoid}

\section{The Growth of Languages}
\label{constr_sync:sec:growth} 

\input{../part4/notions/growth}

%\section{General Assumptions on Constraint Languages and Constraint PDFAs}
\section{Assumptions on Constraint Languages and Constraint PDFAs}
\label{constr_sync:sec:assumptions}
 
\input{../part4/notions/remarkconstraut}

\section{Discussion}
\label{constrsync:notions:sec:discussion}

\input{../part4/notions/discussion}

\chapter{General Results}
\label{constr_sync:chap:gen_results}

%!TEX root = ./../../springer-mono-format/book.tex
%!TEX root = ./../../thesis-trier.tex

In~\cref{constr_sync:sec:prelim_results} we show that \constrsync{L} is decidable and in \PSPACE.
Furthermore, we give example constraint languages
yielding constrained synchronization 
problems with computational complexities 
in $\AC^0$, or \LOGSPACE-, \NL-, \NP-, or \PSPACE-complete.

In~\cref{constr_sync:sec:manip} we present results on altering and simplifying 
the constraint language and how certain operations as forming factor sets, unions, intersections, morphisms, and concatenation
affect the computational complexity.
This also yields
general methods to deduce complexity lower bounds between different instances of constrained synchronization problems.

General hardness conditions for \NL, \NP, and \PSPACE\  are given in \cref{constr_sync:sec:hardness}.
This section also develops a general framework to deduce these hardness results introducing
the notion of auxiliary state avoiding words and the introduction of relations derived from the Myhill-Nerode relation.

In~\cref{constr_sync:sec:aut} we discuss the representation of the constraint by PDFAs.
We show that we have some freedom in the choice of a start state and a set of final states.

\section{Basic Complexity Results and Examples}
\label{constr_sync:sec:prelim_results}

\input{../part4/genresults/basicexamples}

\section{Manipulating Constraint Languages}
\label{constr_sync:sec:manip}

\input{../part4/genresults/manipulating}

\section{Hardness Conditions}
\label{constr_sync:sec:hardness}

\input{../part4/genresults/new_pspace}


\section{Remarks on Constraint PFDAs}
\label{constr_sync:sec:aut}

\input{../part4/genresults/startfinal}

\section{Discussion}
\label{constr_sync:sec:gen_results:discussion}
	
\input{../part4/genresults/discussion} 

\chapter{Classes of Constraint Languages}
\label{constr_sync:chap:classes_constraint_languages}

%!TEX root = ./../../thesis-trier.tex

\input{../part4/constrlang/intro}

\section{Infinite Languages}
\label{sec:constr_sync:infinite}

\input{../part4/constrlang/infinite.tex}

\section{Sparse Languages}
\label{sec:constr_sync:letter_bounded}

\input{../part4/constrlang/sparse}

\section{Monoid Languages}
\label{constr_sync:sec:monoid}

\input{../part4/constrlang/monoid}

\section{Small Constraint Automata}
\label{sec:constr_sync:small_aut}

\input{../part4/constrlang/new_small_aut}

\section{Commutative Languages}
\label{sec:constr_sync:commutative}

\input{../part4/constrlang/comm_aut}

\section{Discussion} 
\label{constrsync:sec:classes:discussion}

\input{../part4/constrlang/discussion}

\chapter{Classes of Input Automata}
\label{constr_sync:chap:constr_input}


\input{../part4/constr_input/introduction}

\section{Monotonic Automata}%\todo{oder Monotone?}
\label{constr_sync:subsec:monotonic_aut}

\input{../part4/constr_input/monotonic}

\section{Partially Ordered Automata}
\label{constr_input:subsec:pto}

\input{../part4/constr_input/waa}

\section{\texorpdfstring{$\rho$}{rho}-Automata With a Trap State}
\label{constr_input:subsec:rho} 

\input{../part4/constr_input/rho} 

\section{Commutative and Unary Automata}
\label{constr_input:subsec:comm}

\input{../part4/constr_input/comm}

\section{Solvable \texorpdfstring{$\rho$}{rho}-Automata}
\label{constr_input:subsec:solvable}

\input{../part4/constr_input/solvable}

\section{Automata with Simple Idempotents}
\label{constr_input:subsec:SI-automata}

\input{../part4/constr_input/simple_idemp}

\section{Discussion}
\label{constr_input:sec:discussion}

\input{../part4/constr_input/discussion}


% https://www.overleaf.com/learn/latex/Questions/BibTeX_isn%27t_working%3B_my_%5Ccite_are_showing_up_as_question_marks_(%3F)
%\nocite{*}
%%\citeall
\printbibliography[heading=bibintoc, title={References}]
\selectlanguage{english}
\printindex

%\todo[inline]{nochmal mit prüfungsordnung abklären, und fochs.pdf prüfen.}
\todo[inline]{kommas richtig?}

\selectlanguage{ngerman}
\clearpage 
\vspace*{8cm}
\begin{description} % https://www.uni-trier.de/studium/beratung-service/service/ordnungen/promotionsordnungen/fachbereich-iv-promotionsordnung
	

	\item[$\Box$] Hiermit erkläre ich an Eides statt, dass ich diese Arbeit selbstständig verfasst und keine anderen als die angegebenen Quellen und Hilfsmittel verwendet habe.
	Es wurden keine Werkzeuge oder Programme, die auf %, wie man sie gegenwärtig bezeichnet,
	künstlicher Intelligenz beruhen, bei der Erstellung dieser Arbeit benutzt
	oder sonstwie herangezogen. 
	Mit der Überprüfung der Arbeit durch einschlägige Plagiatserkennungssoftware bin ich einverstanden.
	Ich habe vor diesem Antrag noch kein Promotionsverfahren beantragt. \\
	
	% okay "\\" hier zu schrieben liefert Fehler....
\end{description}

%Diese Arbeit habe ich selbstständig verfasst und keine anderen als die angegebenen Quellen und Hilfsmittel verwendet.
%Es wurden keine Werkzeuge oder Programme, die auf %, wie man sie gegenwärtig bezeichnet,
%künstlicher Intelligenz beruhen, bei der Erstellung dieser Arbeit benutzt
%oder sonstwie herangezogen. Sie ist alleiniges Produkt einer, nach Ansicht vieler, bald obsoleten und gebrechlichen, durch Emotionen,
%sowie metaphysischen und epistemologsichen Fragestellungen getriebenen, menschlichen  	\glqq Intelligenz \grqq{}.	 \todo{kommas, sapiens statt intelligenz}
%Mit der Überprüfung der Arbeit durch einschlägige Plagiatserkennungssoftware bin ich einverstanden.
%Ich habe vor diesem Antrag noch kein Promotionsverfahren beantragt. \\

%\noindent {\Large Selbstständigkeitserklärung}

%Hiermit versichere ich, dass die vorliegende 



\vspace{3cm}

\begin{minipage}[t]{3cm}
	\rule{3cm}{0.5pt}
	Datum
\end{minipage}
\hfill
\begin{minipage}[t]{9cm}
	\rule{9cm}{0.5pt}
	Unterschrift der Kandidatin/des Kandidaten
\end{minipage}
%\clearpage
%\begin{center} 
	\textbf{Wissenschaftlicher Werdegang} 
\end{center} 

%\hspace*{-1.2cm}
\begin{tabular}{lp{8cm}p{6cm}}
	% 2011 & Intermediate Diploma in Mathematics                  & MLU Halle-Wittenberg, Halle (Saale), Germany \\
%	2004 & Abitur, High School Diploma                            & Gymnasium im Bildungszentrum, Halle (Saale), Germany \\
	2014 & Diploma in Computer Science (with Minor in Psychology) & MLU Halle-Wittenberg, Halle (Saale), Germany \\
	2025 & Ph.D. in Theoretical Computer Science                  & Trier University, Trier, Germany \\
	& \emph{Advisor:} Prof. Dr. Henning Fernau                    &
\end{tabular}
%
%\begin{center} \textbf{Professional Experience}
%\end{center} 
%
%\hspace*{-1.2cm}
%\begin{tabular}{lp{12.8cm}}
%	2010-2014 & \emph{It Services / System Administrator}, Bioscience Center Weinberg Campus, Halle (Saale) \\ 
%	2014-2016 & \emph{Full Stack Database Programmer}, Psychiatric Hospital, Halle (Saale) \\
%	2016-2022 & \emph{Research and Teaching Assistant}, Trier University, Faculty of Computer Science, Trier \\
%	& Supervised Courses and Seminar: \\
%	& $\quad$ Foundations of Discrete Mathematics (each winter term) \\
%	& $\quad$ Automata and Formal Language Theory (each summer term) \\
%	& $\quad$ Student Seminars on Various Topics (Kernelization, Parameterized \\ 
%	& $\quad$ Algorithms, Graph Drawing, Data Compression, Classical \\ 
%	& $\quad$ Works from Theoretical Computer Science) \\
%	& $\quad$ Specialized Courses on Formal Language Theory \\ 
%	& $\quad$ (Descriptional Complexity, Formal Systems, Learning Algorithms)
%\end{tabular}
%\input{note2}
\end{document}
